\chapter{问答与对话}
\label{chap:nlp-question-answering-dialogue}

一份通知能够被准确复述，并不意味着系统已经能够回答围绕它提出的问题。用户问“我明天下午到可以吗”，系统要先在资料中找到时间和报名条件，弄清咨询发生在哪一天、用户属于哪类参加者，再判断哪些条件尚未满足；若下一轮用户更正了日期，系统必须撤掉旧的要求；若查询报名记录失败，回复也不能把预期结果说成已经成立的事实。第\ref{chap:nlp-text-generation}章讨论怎样从给定材料写出文本，本章在此之上增加两种约束：回答必须有可检查的依据，多轮交流必须维护当前有效的目标。

本章的例子均为教学构造。共同通知发布于北京时间2025年4月15日，内容为：“机器学习读书会原定于4月17日举行，现改至4月18日（周五）14:00，地点为海棠楼201。校内师生可直接参加，校外来宾须在4月16日18:00前提交报名信息。”活动编号为\code{ML-0418}。通知没有提供实际报名人数、主讲人或报名成功记录；后文需要这些数据时，另行声明报名快照、统计表或假设的业务状态。三者来源不同：回答的证据来自资料，业务状态来自服务的真实返回，两者不能互换。

围绕这份材料，本章要回答九个问题，见表\ref{tab:nlp-dialogue-challenges}。前五个属于问答：在大规模资料中找到相关证据；从证据中定位答案，或把多个片段、表格单元格与计算组合成答案；证据不足时弃答；回答附上可核对的引用；对已有回答逐项核验事实。后四个属于对话：跨轮记录用户需求，并处理更正与撤回；选择下一步动作并调用真实服务；在没有明确任务的开放交流中，避免“我不知道”一类通用乏味的回复，并保持人物设定与知识一致；在用户目标含混时决定是否澄清、怎样澄清。

\begin{table*}[htbp]
  \centering
  \small
  \begin{tabularx}{\linewidth}{@{}>{\raggedright\arraybackslash}p{38mm}
    >{\raggedright\arraybackslash}X>{\raggedright\arraybackslash}p{24mm}@{}}
    \toprule
    挑战 & 通知中的例子 & 主要处理位置\\
    \midrule
    在大规模资料中找到证据 & 时间在片段A，报名条件在片段B，干扰片段C词面相近 & 第\ref{sec:nlp-dialogue-evidence}节\\
    定位或组合答案 & 时间是原文跨度；“截止前校外报名几人”要在表上筛选求和 & 第\ref{sec:nlp-dialogue-evidence}节\\
    证据不足时弃答 & 通知没有主讲人 & 第\ref{sec:nlp-dialogue-evidence}、\ref{sec:nlp-dialogue-verification}节\\
    回答附可核对的引用 & 时间引用A，报名截止引用B & 第\ref{sec:nlp-dialogue-evidence}节\\
    核验已有回答的逐项事实 & 回答把日期写成4月17日 & 第\ref{sec:nlp-dialogue-verification}节\\
    跨轮记录需求与更正 & 用户先说明天、再改口后天 & 第\ref{sec:nlp-dialogue-task}节\\
    选择动作并调用真实服务 & 查询报名记录，接口超时 & 第\ref{sec:nlp-dialogue-task}节\\
    开放交流中避免通用回复、保持一致 & 闲聊读书会时只回答“挺好的” & 第\ref{sec:nlp-dialogue-open}节\\
    目标含混时澄清 & “把通知改简单些”没有说保留什么 & 第\ref{sec:nlp-dialogue-collaboration}节\\
    \bottomrule
  \end{tabularx}
  \caption{本章的挑战清单。同一份通知在不同任务中暴露不同困难；最后一列给出主要讨论的位置，多项挑战会在后续各节中反复出现。}
  \label{tab:nlp-dialogue-challenges}
\end{table*}

方法谱系在问答侧与对话侧各有一条主线。问答侧依次是：给定段落的抽取式阅读，由手工设计的问文交互层走到预训练编码器；检索加阅读，再让检索器与阅读器联合训练；检索增强生成，把检索结果作为生成条件，并在查询、排序、检索时机、多跳与引用各环节逐一改进；证据为数据库、表格或文本与表格混合时，答案由执行查询或数值程序得出。对话侧依次是：模块化流水线，把一轮交流拆成理解、状态追踪、策略与回复生成；各模块的神经化与端到端序列生成；大语言模型加函数调用，把这些模块合进一次生成，再用显式状态与权限约束把它拉回可检查的范围。开放对话与协作对话沿用这些部件，但完成标准不同：前者看相关、信息量与一致性，后者看最终产物是否满足用户确认过的要求。

各节的安排也按这两条线展开。第\ref{sec:nlp-dialogue-evidence}节从资料中取得证据并形成答案；第\ref{sec:nlp-dialogue-verification}节反过来，把一个已有回答拆成主张，逐条检索证据并修订，它处理的正是证据问答留下的“回答对不对”的问题。第\ref{sec:nlp-dialogue-task}节转入多轮，用户需求在轮次之间变化，回答的依据除了资料还有服务返回的业务状态。第\ref{sec:nlp-dialogue-open}节去掉明确任务，第\ref{sec:nlp-dialogue-collaboration}节让目标在交流中逐步确定。第\ref{sec:nlp-dialogue-comparison}节在固定资料、服务与预算下比较各类方法的评价对象与分母。本章用到的预训练编码器与跨度头见第\ref{chap:nlp-classical-models}章，检索器训练见第\ref{chap:nlp-classification-matching-organization}章，编码器—解码器生成见第\ref{chap:nlp-text-generation}章，偏好对齐见第\ref{chap:nlp-preference-alignment}章；各节在首次用到时说明所用结论。

\FloatBarrier
\section{证据问答}
\label{sec:nlp-dialogue-evidence}

证据问答的输入不只有问题$q$，还包括允许访问的资料、资料版本、查询时点和计算预算；输出可以是原文跨度、生成的回答、表格单元格或计算值，也可以是“给定材料不足以回答”。这里的“不足”是对可用证据的判断，不是断言世界上没有答案。本节固定一份教学资料：片段A写改期后的时间与地点，片段B写校内、校外的参加条件，片段C是无关的教室借用说明；必要时另行声明教学表格。

这一任务的挑战可以列为五项。其一，资料规模大而答案只在少数片段中，找到证据本身就是问题。其二，答案的形式不一：可能是原文跨度、需要综合多个片段的生成文本、表格单元格，或要经过计算的数值。其三，一部分问题需要多跳组合，第二步要查什么取决于第一步查到了什么。其四，证据不足时系统应当弃答，而不是给出一个类型相似的短语。其五，回答应能追溯到来源，读者可以逐句核对。

方法谱系大致沿这五项挑战推进：给定段落的抽取式阅读只解决定位与弃答（第\ref{sec:nlp-dialogue-text-qa}小节）；开放域检索问答在阅读前加入检索，并让检索器与阅读器联合训练（第\ref{sec:nlp-dialogue-open-domain}小节）；检索增强生成把检索结果作为生成条件，再逐环节改进（第\ref{sec:nlp-dialogue-rag-family}小节）；证据为数据库或表格时转为执行查询或选格聚合（第\ref{sec:nlp-dialogue-structured-qa}小节），文本与表格混合时再加入数值程序（第\ref{sec:nlp-dialogue-hybrid-qa}小节）。评价时，检索遗漏、阅读错误、组合错误与引用错误分别统计：只召回B而缺活动时间是覆盖失败，把报名截止当成活动时间是组合错误，答案正确却引用了C是引用错误。表\ref{tab:nlp-dialogue-evidence-methods}汇总各方法解决的问题与新代价。

\subsection{抽取式阅读理解}
\label{sec:nlp-dialogue-text-qa}

\term{抽取式阅读理解}（extractive reading comprehension）给定问题$q=(q_1,\ldots,q_J)$和一段材料$c=(c_1,\ldots,c_T)$，要求输出材料中的一个连续跨度$(s,e)$，$1\leq s\leq e\leq T$；材料不含答案时输出空答案。SQuAD把这一设定做成大规模数据集，答案由标注者在维基百科段落中圈出\scite{nlp-qa-rajpurkar2016squad}；SQuAD 2.0又加入了看似可答、实则材料中没有答案的问题，要求模型学会弃答\scite{rajpurkar2018squad2}

用完整通知构造三类问题。“读书会何时举行”可以定位到跨度“4月18日（周五）14:00”；“谁可直接参加”定位到“校内师生”，它是类别而不是出席名单；“主讲人是谁”在材料中没有答案，正确输出是弃答。这三个问题依次检验定位、答案类型与证据边界。下面两个模型都只回答“答案已在这段材料中时怎样定位”；在全文资料中能否找到这段材料，留给第\ref{sec:nlp-dialogue-open-domain}小节。

\subsubsection{BiDAF问文交互阅读}
\label{sec:nlp-dialogue-bidaf}

定位跨度时，材料中每个词是否属于答案取决于问题问的是什么。早期的注意力阅读模型通常把问题或材料先汇总成一个固定向量，再用它去匹配另一方，问题中的不同词对材料不同位置的作用被过早平均掉了。\term{双向注意力流}（bidirectional attention flow, BiDAF）的思想是：不做早期汇总，在每个材料位置都计算一份与问题相关的表示，并让这些逐位置的表示连同原有的上下文表示一起“流”入后续层\scite{nlp-dialogue-seo2017}。模型分为六层，结构见图\ref{fig:nlp-dialogue-bidaf}。

\begin{figure}[htbp]
  \centering
  % 第九卷 13.1.1 节：BiDAF的六层结构（正文栏宽，约108 mm）。
% 蓝色：输入与嵌入；黄色描边：编码与交互层；红色：输出分布。
% 左列为材料（长度T），右列为问题（长度J），在注意力流层汇合。
\begin{tikzpicture}[
  x=1mm,y=1mm,
  figure picture,
  emb/.style={figure node,draw=FigureInput,fill=FigureInputFill,
    minimum height=8mm,minimum width=40mm,inner ysep=2pt,
    font=\figurefont\scriptsize},
  enc/.style={figure node,draw=FigureModel,fill=FigureModelFill,
    minimum height=8mm,minimum width=40mm,inner ysep=2pt,
    font=\figurefont\scriptsize},
  wide/.style={enc,minimum width=88mm},
  res/.style={figure node,draw=FigureOutput,fill=FigureOutputFill,
    minimum height=8mm,minimum width=40mm,inner ysep=2pt,
    font=\figurefont\scriptsize},
  lab/.style={font=\figurefont\scriptsize,text=FigureMuted,anchor=east,align=right},
  fl/.style={figure flow}
]
  % 输入
  \node[font=\figurefont\scriptsize] (xin) at (22,0) {材料词$x_1,\ldots,x_T$};
  \node[font=\figurefont\scriptsize] (qin) at (70,0) {问题词$q_1,\ldots,q_J$};
  % 1-2 字符与词嵌入
  \node[emb] (xe) at (22,11) {字符CNN与词向量$\to$Highway\\$\bm X\in\mathbb R^{d\times T}$};
  \node[emb] (qe) at (70,11) {字符CNN与词向量$\to$Highway\\$\bm Q\in\mathbb R^{d\times J}$};
  % 3 上下文层
  \node[enc] (xh) at (22,25) {双向LSTM\\$\bm H\in\mathbb R^{2d\times T}$};
  \node[enc] (qu) at (70,25) {双向LSTM（共享参数）\\$\bm U\in\mathbb R^{2d\times J}$};
  % 4 注意力流层
  \node[wide,minimum height=13mm] (att) at (46,42)
    {相似矩阵$\bm S\in\mathbb R^{T\times J}$，$S_{t,j}=\bm w_{\mathrm S}^{\top}[\bm h_t;\bm u_j;\bm h_t\odot\bm u_j]$\\
     材料到问题：逐行softmax得$\widetilde{\bm U}$；\quad 问题到材料：行最大后softmax得$\widetilde{\bm H}$\\
     $\bm G_{:,t}=[\bm h_t;\tilde{\bm u}_t;\bm h_t\odot\tilde{\bm u}_t;\bm h_t\odot\tilde{\bm h}_t]$，$\bm G\in\mathbb R^{8d\times T}$};
  % 5 建模层
  \node[enc] (m) at (22,60) {两层双向LSTM\\$\bm M\in\mathbb R^{2d\times T}$};
  \node[enc] (m2) at (70,60) {再一层双向LSTM\\$\bm M_2\in\mathbb R^{2d\times T}$};
  % 6 输出层
  \node[res] (p1) at (22,75) {起点$\bm p^{(1)}=\operatorname{softmax}(\bm w_1^{\top}[\bm G;\bm M])$};
  \node[res] (p2) at (70,75) {终点$\bm p^{(2)}=\operatorname{softmax}(\bm w_2^{\top}[\bm G;\bm M_2])$};

  \draw[fl] (xin) -- (xe); \draw[fl] (qin) -- (qe);
  \draw[fl] (xe) -- (xh); \draw[fl] (qe) -- (qu);
  \draw[fl] (xh) -- (xh |- att.south); \draw[fl] (qu) -- (qu |- att.south);
  \draw[fl] (att.north -| m) -- (m);
  \draw[fl] (m) -- (m2);
  \draw[fl] (m) -- (p1);
  \draw[fl] (m2) -- (p2);
  % G 直接进入输出层
  \draw[fl,draw=FigureMuted,dash pattern=on 2pt off 1.5pt]
    (att.east) -- ++(4,0) |- (p2.east);
  \draw[fl,draw=FigureMuted,dash pattern=on 2pt off 1.5pt]
    (att.west) -- ++(-4,0) |- (p1.west);

  \node[lab] at (-1,11) {嵌入};
  \node[lab] at (-1,25) {上下文};
  \node[lab] at (-1,42) {注意力流};
  \node[lab] at (-1,60) {建模};
  \node[lab] at (-1,75) {输出};
\end{tikzpicture}

  \caption{BiDAF的六层结构。左列处理长度为$T$的材料，右列处理长度为$J$的问题，两列在注意力流层由同一个相似矩阵连接；交互后的表示$\bm G$既送入建模层，也经虚线直接送入输出层。起点与终点分别由两组表示打分，终点一侧多经过一层双向LSTM。张量形状按原文约定写为“特征维度$\times$序列长度”。}
  \label{fig:nlp-dialogue-bidaf}
\end{figure}

前三层分别编码两段文本。字符嵌入层把每个词的字符序列送入一维卷积并在宽度方向做最大池化，得到对未登录词也有效的字符级向量；词嵌入层取固定的预训练词向量。两者拼接后经两层Highway网络\cite{srivastava2015highway}，得到材料表示$\bm X\in\mathbb R^{d\times T}$和问题表示$\bm Q\in\mathbb R^{d\times J}$。上下文嵌入层用一个双向LSTM分别编码两者，拼接两个方向的输出，得到$\bm H\in\mathbb R^{2d\times T}$与$\bm U\in\mathbb R^{2d\times J}$；记它们的第$t$列与第$j$列为$\bm h_t$与$\bm u_j$。双向LSTM的结构与训练可参考\BookChapterRef{chap:recurrent-networks}{卷五的循环神经网络章}，这里只用到它的输出：每个位置的向量同时概括了左右两侧的上下文。原文取$d=100$，字符卷积用100个宽度为5的滤波器。

第四层是注意力流层。两个方向的注意力共用一个相似矩阵$\bm S\in\mathbb R^{T\times J}$，它的元素衡量材料第$t$个词与问题第$j$个词的匹配程度：
\begin{equation}
  S_{t,j}=\bm w_{\mathrm S}^{\top}
  [\bm h_t;\bm u_j;\bm h_t\odot\bm u_j],
  \qquad \bm w_{\mathrm S}\in\mathbb R^{6d}.
  \label{eq:nlp-dialogue-bidaf}
\end{equation}
其中$\odot$表示逐元素乘积，$[\cdot;\cdot]$表示拼接。三段拼接分别让打分看到材料词本身、问题词本身和两者逐维的乘积；只用内积$\bm h_t^{\top}\bm u_j$相当于把$\bm w_{\mathrm S}$固定为只作用于乘积项的全1向量，学习的$\bm w_{\mathrm S}$更灵活。

由$\bm S$可以从两个方向读出注意力。材料到问题（context-to-query, C2Q）的注意力回答“对材料第$t$个词而言，问题中哪些词最相关”：对第$t$行做softmax得到$\bm a_t=\operatorname{softmax}(\bm S_{t,:})\in\mathbb R^{J}$，再加权问题向量得$\tilde{\bm u}_t=\sum_j a_{t,j}\bm u_j$，各列组成$\widetilde{\bm U}\in\mathbb R^{2d\times T}$。问题到材料（query-to-context, Q2C）的注意力回答“材料中哪些词与问题的某个词最接近”：先对每行取最大值，再在材料位置间做softmax，$\bm b=\operatorname{softmax}(\max_{j}S_{t,j})_{t=1}^{T}\in\mathbb R^{T}$，得材料摘要$\tilde{\bm h}=\sum_t b_t\bm h_t$，复制$T$次成为$\widetilde{\bm H}\in\mathbb R^{2d\times T}$。两方向的结果与原表示逐位置拼接：
\begin{equation}
  \bm G_{:,t}=[\bm h_t;\tilde{\bm u}_t;\bm h_t\odot\tilde{\bm u}_t;\bm h_t\odot\tilde{\bm h}_t]\in\mathbb R^{8d}.
  \label{eq:nlp-dialogue-bidaf-g}
\end{equation}
这里每个材料位置都得到自己的问题相关表示，注意力权重在不同时刻之间互不依赖，原文称之为无记忆注意力；$\bm h_t$本身也保留在$\bm G$中，交互层不会抹去材料的原始上下文。

用一个两行两列的教学子矩阵$\bm S=\left[\begin{smallmatrix}2&0\\0&1\end{smallmatrix}\right]$看两个方向的计算。逐行softmax得到C2Q权重$(0.881,0.119)$和$(0.269,0.731)$：第一个材料词主要读取第一个问题词，第二个材料词主要读取第二个问题词。行最大值为$(2,1)$，softmax后得Q2C权重$(0.731,0.269)$，材料摘要偏向第一个词。这些权重决定送往后续层的表示，本身并不是答案概率。

第五层是建模层，用两层双向LSTM读取$\bm G$，得$\bm M\in\mathbb R^{2d\times T}$。它与上下文层的区别在于输入：上下文层只看材料自身，建模层看到的是已经带着问题信息的表示，因而能建模“在这个问题下，材料各词之间的关系”。第六层是输出层，起点与终点分布分别为
\begin{equation}
  \bm p^{(1)}=\operatorname{softmax}\bigl(\bm w_1^{\top}[\bm G;\bm M]\bigr),\qquad
  \bm p^{(2)}=\operatorname{softmax}\bigl(\bm w_2^{\top}[\bm G;\bm M_2]\bigr),
  \label{eq:nlp-dialogue-bidaf-output}
\end{equation}
其中$\bm w_1,\bm w_2\in\mathbb R^{10d}$，$\bm M_2\in\mathbb R^{2d\times T}$是把$\bm M$再送入一层双向LSTM的输出，softmax在$T$个材料位置上进行。终点一侧多一层循环编码，使终点预测能够参考起点附近已形成的表示。训练最小化真实起点与终点的负对数概率之和在样本上的平均；推断时在$s\leq e$的约束下选使$p^{(1)}_s p^{(2)}_e$最大的$(s,e)$。固定$e$时只需知道$\max_{s\leq e}p^{(1)}_s$，从左到右扫描一遍维护这个前缀最大值，就能在$O(T)$时间内找到最优跨度。

BiDAF确立了“逐位置交互、交互表示再建模、起止两头输出”的阅读器框架，后续许多阅读器在此基础上修改注意力形式或叠加自注意力。它的局限也来自这套手工设计：交互只在一层中发生；所有参数在问答数据上从头训练，词向量之外没有预训练知识；LSTM沿序列串行计算；输出层只能给出单个跨度，原始设置也不包含弃答。

\subsubsection{BERT跨度阅读与弃答}
\label{sec:nlp-dialogue-bert-span}

BiDAF的交互层是为阅读任务专门设计的。若把问题与材料拼成一个序列交给预训练的双向编码器，每一层自注意力都允许材料位置读取问题位置、问题位置读取材料位置，C2Q与Q2C两类交互就成为自注意力的一部分，而且在每一层反复发生。BERT正是这样做的：输入写成\code{[CLS] 问题 [SEP] 材料 [SEP]}，问题和材料分别用两种段嵌入区分，编码器之上只保留起点向量与终点向量组成的跨度头，起止分布与训练目标见第\ref{sec:nlp-classical-bert-heads}小节的式\eqref{eq:nlp-classical-span-head}\scite{devlin2019}。预训练取代了手工交互层，模型结构反而更简单；在SQuAD 1.1开发集上，单模型F1由BiDAF的约77提高到BERT-large的约91。

阅读任务在通用跨度头之外还需要处理三件事。第一是长材料。编码器有最大长度，材料超长时按固定步长切成重叠窗口，每个窗口都与问题拼接；BERT公开实现的SQuAD脚本取最大长度384、步长128。训练时，答案不在当前窗口内的样本把起止都标到\code{[CLS]}。每个窗口还要保存子词到原文字符的偏移，以便把预测的子词区间还原为原文字符串，而不是把子词拼接起来当作答案。

第二是合法跨度与弃答。记起点向量与终点向量为$\bm w_{\mathrm s},\bm w_{\mathrm e}$，位置表示为$\bm h_i$。合法跨度要求$i\leq j$、长度不超过$\ell_{\max}$且两端都在材料段内，最佳跨度分数和空答案分数分别为
\begin{equation}
  s_{\mathrm{span}}=\max_{\substack{i\leq j\leq i+\ell_{\max}-1\\ i,j\in\text{材料}}}
  \bigl(\bm w_{\mathrm s}^{\top}\bm h_i+\bm w_{\mathrm e}^{\top}\bm h_j\bigr),
  \qquad
  s_{\mathrm{null}}=\bm w_{\mathrm s}^{\top}\bm h_{\mathrm C}+\bm w_{\mathrm e}^{\top}\bm h_{\mathrm C},
  \label{eq:nlp-dialogue-span}
\end{equation}
其中$\bm h_{\mathrm C}$是\code{[CLS]}位置的表示。当$s_{\mathrm{span}}-s_{\mathrm{null}}>\tau$时返回跨度，否则弃答；阈值$\tau$在开发集上按F1选定。分数差是未归一化的对数尺度，同一个$\tau$在不同窗口间可比，这一点比比较各窗口内单独归一化的概率更可靠。

用压缩后的材料词元“读书会／改至／4月18日／14:00／地点／海棠楼201／校内师生／可直接参加”演算，以下概率均为教学假设。对时间问题，位置3、4的起点概率为$0.60,0.10$，终点概率为$0.15,0.70$。跨度$(3,4)$的乘积为$0.42$，两个单词元候选分别为$0.60\times0.15=0.09$和$0.10\times0.70=0.07$。设空答案的乘积为$0.02$，同一归一化下分数差为$\log(0.42/0.02)=3.045$，在教学阈值$\tau=1$时返回“4月18日14:00”。对“主讲人是谁”，若最佳跨度的乘积为$0.03$、空答案为$0.30$，分数差为$-2.303$，模型弃答。

第三是跨窗口汇总。每个窗口给出自己的最佳非空候选和空答案分数；汇总时取所有窗口中最大的$s_{\mathrm{span}}-s_{\mathrm{null}}$，再与$\tau$比较。若答案分散在两个窗口，例如“校内学生和教师”需要两个不相邻的短语，单跨度头无法恢复，需要多跨度标注或改为生成。与BiDAF相比，BERT阅读器的优势来自预训练与深层交互，代价是长度上限带来的分窗、每个窗口一次完整的编码计算，以及对输出类型的同样限制：它仍只能抽取，不能综合。

\subsection{开放域检索问答}
\label{sec:nlp-dialogue-open-domain}

\term{开放域问答}（open-domain question answering）不给定段落，系统要先从大规模文集中找出可能含答案的片段，再阅读。DrQA确立了这条\term{检索器}（retriever）加\term{阅读器}（reader）的流水线：检索器用词袋统计从维基百科中取回少量文章，阅读器在其中抽取答案\scite{nlp-qa-chen2017drqa}。Natural Questions用真实搜索查询配合维基百科页面标注长短答案，成为这一设定的常用数据\scite{nlp-qa-kwiatkowski2019nq}。流水线把问题分成两半：检索器决定答案是否有机会被读到，阅读器决定读到之后能否定位。

这一设定新增两个困难。第一，词汇检索依赖词项重合，“校外怎么参加”与“须提交报名信息”几乎没有共同词，答案片段可能根本进不了候选；第\ref{sec:nlp-dialogue-dpr}小节用稠密检索处理它。第二，检索器与阅读器分开训练时，检索器并不知道阅读器实际需要什么；第\ref{sec:nlp-dialogue-joint-retriever}小节把检索当作潜变量、用答案似然训练检索器。

\subsubsection{DPR检索与阅读}
\label{sec:nlp-dialogue-dpr}

先看词汇检索在教学资料上的表现。BM25的打分公式与参数含义见第\ref{sec:nlp-classification-bm25}小节。表\ref{tab:nlp-dialogue-retrieval}直接规定一个教学索引的统计：查询只保留“读书会、报名”，$k_1=1.2$，$b=0.75$，$n=3$个片段，平均长度$40/3$。两个词的逆文档频率为$\log1.6=0.470$和$\log(8/3)=0.981$，例如B的分数为$0.470\times2.2/2.65+0.981\times4.4/3.65=1.573$。取前两项得到B、A；只取第一项会缺活动时间。问题若写成“校外怎么参加”，“报名”一词不出现，B的分数也随之消失。

\begin{table}[htbp]
  \centering
  \begin{tabular}{lrrrr}
    \toprule
    片段 & 长度 & “读书会”词频 & “报名”词频 & BM25分数\\
    \midrule
    A：时间地点 & 10 & 1 & 0 & 0.524\\
    B：参加条件 & 20 & 1 & 2 & 1.573\\
    C：借用说明 & 10 & 0 & 0 & 0\\
    \bottomrule
  \end{tabular}
  \caption{固定教学索引上的词汇召回。长度为教学分词统计；片段A、B互补，排序靠前与证据覆盖完整不是同一个条件。}
  \label{tab:nlp-dialogue-retrieval}
\end{table}

DPR用两个BERT编码器分别把问题和片段映射为向量$\bm q=E_q(q)$、$\bm z=E_z(z)\in\mathbb R^{d}$，以内积$s(q,z)=\bm q^{\top}\bm z$打分；训练目标是第\ref{sec:nlp-classification-dpr}小节的对比损失，即式\eqref{eq:nlp-matching-infonce}取内积打分、以批内其他问题的正片段和BM25难负例为负例\scite{nlp-karpukhin2020dpr}。难负例挖掘与蒸馏也在那里讨论，这里只看DPR在问答流水线中的位置。

离线阶段，文集按固定长度切成片段并逐一编码。DPR把英文维基百科切成每段100词的片段，共约2100万个，向量存入近似最大内积索引；倒排索引与近似近邻检索的工程实现可参考\BookChapterRef{chap:rs-search}{卷十一的搜索相关性章}。切分时每个片段保留标题和对象编号，否则“18:00前”离开了“校外来宾”和“ML-0418”就失去了所指。在线阶段只编码问题，取内积最高的前$k$个片段。教学向量取$\bm q=(1,1)$，$\bm z_A=(2,0)$，$\bm z_B=(0,3)$，$\bm z_C=(-1,0)$，内积为$2,3,-1$，前两项为B、A。这组二维向量只用来核对计算，不代表真实训练所得的表示。

阅读器面对的是$k$个片段，而不是一个。DPR的阅读器对每个片段$z_i$单独编码\code{[CLS] 问题 [SEP] 片段}，得表示$\bm P_i\in\mathbb R^{L\times h}$，并同时做两件事：在片段内给跨度打分，在片段间做选择。三个分布为
\begin{equation}
\begin{aligned}
  P_{\mathrm{start},i}(s)&=\operatorname{softmax}(\bm P_i\bm w_{\mathrm{start}})_s,\qquad
  P_{\mathrm{end},i}(t)=\operatorname{softmax}(\bm P_i\bm w_{\mathrm{end}})_t,\\
  P_{\mathrm{sel}}(i)&=\operatorname{softmax}(\widehat{\bm P}^{\top}\bm w_{\mathrm{sel}})_i,
\end{aligned}
  \label{eq:nlp-dialogue-dpr-reader}
\end{equation}
其中矩阵$\widehat{\bm P}\in\mathbb R^{h\times k}$的第$i$列是片段$z_i$的\code{[CLS]}表示，$\bm w_{\mathrm{start}},\bm w_{\mathrm{end}},\bm w_{\mathrm{sel}}\in\mathbb R^{h}$是可学习向量。片段选择分布在$k$个片段之间归一化，它的作用相当于一个交叉编码的重排器：问题与片段在编码时充分交互，比双编码器的内积更精细，但只能在少量候选上计算。训练时，从检索结果前100项中取一个含答案的正片段和若干负片段，最大化正片段中全部答案出现位置的边缘对数似然，并加上正片段被选中的对数似然；推断时先选$P_{\mathrm{sel}}$最高的片段，再在其中取$P_{\mathrm{start}}\times P_{\mathrm{end}}$最大的跨度。

在教学资料上，问“校外来宾报名何时截止”，阅读器应选B并返回“4月16日18:00前”；若同时问活动时间与报名截止，单个片段只能回答一半，两个子答案要分别提问或由外部模板拼合。一次输出综合多个片段，要等到第\ref{sec:nlp-dialogue-rag-family}小节的生成式方法。

第\ref{sec:nlp-classification-dpr}小节讨论过的假负例，在问答中还有一种特有形式：需要两项事实的问题，A和B都是有效证据，不能因为某轮训练把B选作正例就把A当作负例；反过来，“含答案字符串即正例”的弱监督也会把同名但无关的片段当作正例。另一个代价来自索引：片段向量出自某个版本的片段编码器，编码器一旦更新就必须重新编码全文，否则新的问题向量与旧的片段向量不在同一空间。下一小节的联合训练正要面对这一代价。

\subsubsection{检索器与阅读器的联合训练}
\label{sec:nlp-dialogue-joint-retriever}

DPR的检索器只从“哪个片段含答案”学习，而逐题标注这一点本身就昂贵。一个片段即使含答案字符串，也可能上下文不足、阅读器读不出来；另一个片段即使不含原词，也可能恰好提供了推出答案所需的事实。联合训练的思路是：不给检索器逐片段的标签，而把检索到哪个片段当作潜变量，以最终答案的似然为目标，让答案的好坏反过来决定检索器应该提高谁、压低谁。

REALM把这一思路用于预训练\scite{nlp-qa-guu2020realm}。输入$x$是一段被遮住部分词元的文本，目标$y$是被遮住的内容，片段$z$取自文集$\mathcal Z$。模型分两步：检索器给出$p(z\mid x)$，读取器在$x$与$z$的拼接上预测$p(y\mid z,x)$；答案似然对片段求和，
\begin{equation}
\begin{aligned}
  p(y\mid x)&=\sum_{z\in\mathcal Z}p(y\mid z,x)\,p(z\mid x),\\
  p(z\mid x)&=\frac{\exp f(x,z)}{\sum_{z'}\exp f(x,z')},\qquad
  f(x,z)=\bm e_{\mathrm{in}}(x)^{\top}\bm e_{\mathrm{doc}}(z).
\end{aligned}
  \label{eq:nlp-dialogue-realm}
\end{equation}
其中$\bm e_{\mathrm{in}}$与$\bm e_{\mathrm{doc}}$是两个BERT编码器的\code{[CLS]}表示经线性投影所得的向量。对全文求和不可行，REALM用最大内积搜索取出前$k$个片段近似求和。

为什么这样训练能教会检索器？对检索器参数求$\log p(y\mid x)$的梯度。由$\nabla\log p(z\mid x)=\nabla f(x,z)-\sum_{z'}p(z'\mid x)\nabla f(x,z')$，代入后整理得
\begin{equation}
  \nabla\log p(y\mid x)=\sum_{z}r(z)\,\nabla f(x,z),\qquad
  r(z)=\Bigl[\frac{p(y\mid z,x)}{p(y\mid x)}-1\Bigr]p(z\mid x).
  \label{eq:nlp-dialogue-realm-grad}
\end{equation}
推导的关键一步是：$\sum_z p(y\mid z,x)p(z\mid x)\nabla\log p(z\mid x)/p(y\mid x)$中，$p(y\mid z,x)p(z\mid x)/p(y\mid x)$正是后验$p(z\mid x,y)$，梯度因此等于$\sum_z[p(z\mid x,y)-p(z\mid x)]\nabla f(x,z)$。$r(z)>0$当且仅当$p(y\mid z,x)>p(y\mid x)$：一个片段让答案的概率高于按检索分布随机取片段时的平均水平，它的分数就被提高，否则被压低。教学例子中，设A、B的检索概率都是0.5，读取器以A为条件时答案概率0.8、以B为条件时0.2，则$p(y\mid x)=0.5$，$r(A)=(0.8/0.5-1)\times0.5=0.3$，$r(B)=-0.3$。检索器没有见过任何片段标签，却得到了“提高A、压低B”的信号。

这个信号也带来索引问题。$\bm e_{\mathrm{doc}}$每更新一次，索引中的向量就过时了。REALM的做法是异步刷新：训练任务持续更新参数，另一个索引任务每隔一段时间取参数快照，重新编码全部片段并重建索引，完成后交回训练任务，在其原文的预训练设置中约每500步刷新一次。两次刷新之间索引略有滞后，但它只用于选出前$k$个候选，候选的$p(z\mid x)$及其梯度仍用当前参数重新计算。预训练还用到两个设计：\term{显著跨度遮蔽}（salient span masking）优先遮住命名实体和日期这类需要世界知识才能补全的跨度；候选中加入一个空片段，供不需要检索的样本使用。微调到开放域问答时，REALM只建一次索引，固定片段编码器，只更新问题编码器和阅读器。

Atlas把联合训练用于少样本学习\scite{nlp-qa-izacard2023atlas}。检索器是Contriever式双编码器，阅读器是T5上的解码器融合模型（见第\ref{sec:nlp-dialogue-fid}小节）。它不对片段求和，而是用阅读器提供一个目标分布，让检索器分布向它靠拢。检索器在前$K$个片段上的分布为$p_{\mathrm{retr}}(z_k\mid q)\propto\exp(s(z_k,q)/\theta)$，$\theta$为温度。\term{困惑度蒸馏}（perplexity distillation）以“只给这个片段时阅读器生成答案$a$的概率”作为目标：
\begin{equation}
  p_k=\frac{\exp\log p_{\mathrm{LM}}(a\mid z_k,q)}{\sum_{i=1}^{K}\exp\log p_{\mathrm{LM}}(a\mid z_i,q)},\qquad
  L_{\mathrm{PDist}}=\mathrm{KL}\bigl(p\,\Vert\,p_{\mathrm{retr}}\bigr),
  \label{eq:nlp-dialogue-pdist}
\end{equation}
目标分布不回传梯度，这一损失只更新检索器。\term{注意力蒸馏}（attention distillation）把阅读器对各片段的交叉注意力汇总成片段分数作为目标，源自Izacard与Grave的阅读器到检索器蒸馏\cite{nlp-qa-izacard2021distilling}；Atlas另外比较了EMDR\textsuperscript{2}式边缘化目标和留一困惑度目标，最终采用较稳定且计算较省的困惑度蒸馏。

Atlas还量化了刷新索引的代价。设文集有$N$个片段，检索器与阅读器参数量为$P_{\mathrm{retr}}$与$P_{\mathrm{lm}}$，每步批量为$B$、每题读$K$个片段，前向加反向与激活重算约为一次前向的4倍，则每$R$步全量重建索引的相对开销约为$N P_{\mathrm{retr}}/(4BKP_{\mathrm{lm}}R)$。按其设置，3700万片段、批量64、每题20个片段、每1000步刷新一次，开销约为30\%。两种替代做法以精度换计算：每步多取$L$个候选并用最新检索器重新编码排序；或者只微调问题编码器、固定片段编码器，索引从此不必更新。在大多数少样本设置中，后一种做法对效果影响不大。Atlas在只用64条训练样本时，在Natural Questions上的精确匹配率超过42\%。

联合训练省去了片段标注，并让检索器的目标与阅读器一致；代价是训练需要重复检索、读取器要对每个候选单独计算以提供信号，以及索引刷新与陈旧性的工程负担。阅读器本身的偏差也会经由梯度或蒸馏目标传给检索器。

\subsection{检索增强生成}
\label{sec:nlp-dialogue-rag-family}

抽取式阅读器只能返回原文跨度，无法写出“活动于4月18日14:00举行，校外来宾须在4月16日18:00前报名”这样综合两个片段的回答。\term{检索增强生成}（retrieval-augmented generation, RAG）把检索结果作为生成器的条件，答案不再受跨度限制。第\ref{sec:nlp-dialogue-rag}与第\ref{sec:nlp-dialogue-fid}小节的RAG与FiD回答“多个片段怎样接入生成器”，二者的结构差别见图\ref{fig:nlp-dialogue-rag-fid}。

\begin{figure*}[htbp]
  \centering
  % 第九卷 13.1.3 节：RAG与FiD接入多个检索片段的位置（跨栏，约150 mm宽）。
% 蓝色：输入与片段；黄色描边：模型部件；红色：输出概率或答案。
\begin{tikzpicture}[
  x=1mm,y=1mm,
  figure picture,
  inp/.style={figure node,draw=FigureInput,fill=FigureInputFill,
    minimum height=7mm,minimum width=18mm,inner xsep=2pt,inner ysep=2pt,
    font=\figurefont\scriptsize},
  mod/.style={figure node,draw=FigureModel,fill=FigureModelFill,
    minimum height=7mm,minimum width=18mm,inner xsep=2pt,inner ysep=2pt,
    font=\figurefont\scriptsize},
  res/.style={figure node,draw=FigureOutput,fill=FigureOutputFill,
    minimum height=8mm,minimum width=60mm,inner xsep=2pt,inner ysep=2pt,
    font=\figurefont\scriptsize},
  note/.style={font=\figurefont\scriptsize,text=FigureMuted,align=center},
  fl/.style={figure flow}
]
  \figurePanel{0}{0}{72}{86}{RAG：在概率层对片段边缘化}{FigureInputFill}
  \figurePanel{78}{0}{72}{86}{FiD：在解码器内融合片段}{FigureInputFill}

  % ---------- RAG ----------
  \node[inp,minimum width=24mm] (rq) at (36,8) {问题$x$};
  \node[mod,minimum width=60mm] (rr) at (36,20)
    {DPR检索器：$p_\eta(z\mid x)\propto\exp(\bm d(z)^{\top}\bm q(x))$，取前$k$个};
  \foreach \i/\xx/\zz in {1/14/{z_1},2/36/{z_2},3/58/{z_k}}{
    \node[inp] (rz\i) at (\xx,32) {片段$\zz$};
    \node[mod] (rg\i) at (\xx,46) {生成器$p_\theta$\\条件$x,\zz$};
    \draw[fl] (rz\i) -- (rg\i);
    \draw[fl] (rr.north -| rz\i) -- (rz\i);
  }
  \node[note] at (47,32) {$\cdots$};
  \node[res] (ro) at (36,62)
    {$\sum_{z}p_\eta(z\mid x)\times$（整句概率或每个词元概率）};
  \foreach \i in {1,2,3}{\draw[fl] (rg\i) -- (rg\i |- ro.south);}
  \draw[fl] (rq) -- (rr);
  \node[note] at (36,73) {每个分支只读一个片段；\\检索概率作为权重参与求和};

  % ---------- FiD ----------
  \foreach \i/\xx/\zz in {1/92/{z_1},2/114/{z_2},3/136/{z_k}}{
    \node[inp] (fz\i) at (\xx,8) {$x$与片段$\zz$};
    \node[mod] (fe\i) at (\xx,21) {编码器\\$\bm H_{\i}\in\mathbb R^{\ell\times d}$};
    \draw[fl] (fz\i) -- (fe\i);
  }
  \node[note] at (125,8) {$\cdots$};
  \node[mod,minimum width=60mm] (fc) at (114,35)
    {拼接$[\bm H_1;\ldots;\bm H_k]\in\mathbb R^{k\ell\times d}$};
  \foreach \i in {1,2,3}{\draw[fl] (fe\i) -- (fe\i |- fc.south);}
  \node[mod,minimum width=60mm] (fd) at (114,49)
    {单个解码器：每步交叉注意全部$k\ell$个位置};
  \node[res] (fo) at (114,62) {$\prod_i p_\theta(y_i\mid y_{<i},[\bm H_1;\ldots;\bm H_k])$};
  \draw[fl] (fc) -- (fd);
  \draw[fl] (fd) -- (fo);
  \node[note] at (114,73) {编码代价随$k$线性增长；\\一个词元可同时依据多个片段};
\end{tikzpicture}

  \caption{RAG与FiD接入多个检索片段的位置。左：RAG让生成器在每个片段条件下各算一次概率，再以检索概率为权重求和，求和可以在整句层面或每个词元层面进行。右：FiD把问题与每个片段分别编码，拼接全部编码结果，由一个解码器在每步对所有片段做交叉注意力；它不对片段做概率加权。$k$为片段数，$\ell$为每个问题—片段输入的长度。}
  \label{fig:nlp-dialogue-rag-fid}
\end{figure*}

接入方式确定后，问题转移到流水线的其他环节，图\ref{fig:nlp-dialogue-rag-pipeline}标出了各改进方法作用的位置。问题的表述与答案文档的表述不一致，检索就会遗漏，查询改写与假设文档检索在检索之前修改查询（第\ref{sec:nlp-dialogue-query-rewrite}小节）；双编码器的初检排序不够精细，交叉编码器在少量候选上重排（第\ref{sec:nlp-dialogue-rerank}小节）；多跳问题一次检索不够，迭代检索让推理与检索交替（第\ref{sec:nlp-dialogue-iterative}小节）；并非每个问题、每句话都需要检索，Self-RAG让模型自己决定何时检索并评判结果（第\ref{sec:nlp-dialogue-self-rag}小节）；回答无法追溯，带引用的生成让每个陈述附上来源并可自动核验（第\ref{sec:nlp-dialogue-citation}小节）。最后讨论另一条路线：把全部资料直接放进长上下文，能否取代检索（第\ref{sec:nlp-dialogue-long-context}小节）。

\begin{figure*}[htbp]
  \centering
  % 第九卷 13.1.3 节：检索增强生成流水线及各改进方法作用的环节（跨栏，约152 mm宽）。
% 实线：一次检索增强生成的数据流；虚线：由生成过程触发的再检索。
% 蓝色：输入；黄色描边：检索与生成部件；红色：交付的回答与核验。
\begin{tikzpicture}[
  x=1mm,y=1mm,
  figure picture,
  inp/.style={figure node,draw=FigureInput,fill=FigureInputFill,
    minimum height=11mm,minimum width=20mm,inner xsep=2pt,inner ysep=2pt,
    font=\figurefont\scriptsize},
  mod/.style={figure node,draw=FigureModel,fill=FigureModelFill,
    minimum height=11mm,minimum width=21mm,inner xsep=2pt,inner ysep=2pt,
    font=\figurefont\scriptsize},
  res/.style={figure node,draw=FigureOutput,fill=FigureOutputFill,
    minimum height=11mm,minimum width=21mm,inner xsep=2pt,inner ysep=2pt,
    font=\figurefont\scriptsize},
  note/.style={font=\figurefont\scriptsize,text=FigureMuted,align=center},
  fl/.style={figure flow},
  back/.style={figure flow,draw=FigureMuted,dash pattern=on 2.5pt off 1.5pt}
]
  \node[inp] (q) at (10,20) {问题$q$};
  \node[mod] (rw) at (36,20) {查询构造\\$q\to\tilde q$};
  \node[mod] (rt) at (62,20) {初检\\BM25或DPR\\前$K$项};
  \node[mod] (rr) at (88,20) {重排\\交叉编码器\\前$k$项};
  \node[mod] (gn) at (114,20) {生成器\\RAG或FiD};
  \node[res] (ct) at (140,20) {带引用回答\\逐陈述核验};

  \draw[fl] (q) -- (rw);
  \draw[fl] (rw) -- (rt);
  \draw[fl] (rt) -- (rr);
  \draw[fl] (rr) -- (gn);
  \draw[fl] (gn) -- (ct);

  % 再检索回路
  \draw[back] (gn.north) -- ++(0,7) -| node[pos=.25,above,note] {Self-RAG：反思标记决定是否检索，并评判片段与生成} (rt.north);
  \draw[back] ([xshift=6mm]gn.north) -- ++(0,15) -| node[pos=.25,above,note] {IRCoT：上一句推理作为新查询，片段并入上下文} ([xshift=-6mm]rt.north);

  % 方法标注
  \node[note] at (36,7) {查询改写\\HyDE假设文档};
  \node[note] at (62,7) {检索器训练\\见第11章};
  \node[note] at (88,7) {以计算换\\排序精度};
  \node[note] at (114,7) {片段边缘化\\或解码器融合};
  \node[note] at (140,7) {ALCE引用\\召回与精确率};
\end{tikzpicture}

  \caption{检索增强生成的流水线及各改进方法作用的环节。实线是一次检索增强生成的数据流：构造查询、初检取前$K$项、重排保留前$k$项、生成、逐陈述核验引用。两条虚线是由生成过程触发的再检索：Self-RAG在每个片段前以反思标记决定是否检索，IRCoT把刚生成的推理句作为下一次检索的查询。下方注明各环节对应的方法。}
  \label{fig:nlp-dialogue-rag-pipeline}
\end{figure*}

\subsubsection{RAG检索条件生成}
\label{sec:nlp-dialogue-rag}

RAG的原模型把文档看作潜变量：不预先指定唯一正确的文档，而把各文档条件下的答案概率按检索概率加权求和\scite{lewis2020rag}。检索器是DPR，$p_\eta(z\mid x)\propto\exp(\bm d(z)^{\top}\bm q(x))$；生成器是BART-large，输入为问题与片段的直接拼接。设检索到的前$k$个片段为$Z_k$，$\pi_z=p_\eta(z\mid x)$在候选内归一化，生成器给出$p_\theta(y_i\mid x,z,y_{<i})$。求和可以放在两个层面：
\begin{equation}
\begin{aligned}
 p_{\mathrm{seq}}(y\mid x)
 &=\sum_{z\in Z_k}\pi_z
      \prod_{i=1}^{m}p_\theta(y_i\mid x,z,y_{<i}),\\
 p_{\mathrm{tok}}(y\mid x)
 &=\prod_{i=1}^{m}
      \sum_{z\in Z_k}\pi_z\, p_\theta(y_i\mid x,z,y_{<i}).
\end{aligned}
\label{eq:nlp-dialogue-rag}
\end{equation}
RAG-Sequence在一次完整答案中使用同一片段，再对片段求和；RAG-Token在每个位置都对片段求和，允许不同位置依赖不同片段。后者并不在每个词元处重新检索，候选集合与检索概率仍只由输入问题决定。

表\ref{tab:nlp-dialogue-rag}给出完整的教学计算。为了突出求和与求积的位置，只比较两个双词元候选，省略两者相同的结束符概率；表中数值是给定各自前缀时的条件概率，不是文献实验。取$\pi_A=0.6$，$\pi_B=0.4$，候选甲为“周五／14:00”，候选乙为“周四／14:00”。

\begin{table}[htbp]
  \centering
  \begin{tabular}{llrr}
    \toprule
    候选 & 条件片段 & 第1词元概率 & 第2词元概率\\
    \midrule
    甲 & A & 0.8 & 0.2\\
    甲 & B & 0.2 & 0.9\\
    乙 & A & 0.1 & 0.5\\
    乙 & B & 0.4 & 0.4\\
    \bottomrule
  \end{tabular}
  \caption{RAG的教学条件概率。第2列表示生成时读取哪个片段；第4列还以该候选的第1词元为条件。}
  \label{tab:nlp-dialogue-rag}
\end{table}

序列级甲得分为$0.6\times0.8\times0.2+0.4\times0.2\times0.9=0.168$，乙为$0.6\times0.1\times0.5+0.4\times0.4\times0.4=0.094$。词元级甲得分为$(0.6\times0.8+0.4\times0.2)(0.6\times0.2+0.4\times0.9)=0.56\times0.48=0.2688$，乙为$(0.06+0.16)(0.30+0.16)=0.1012$。两种模型在这两个候选中都选甲，但分数含义不同。在词元级模型下，甲第一位置的片段后验为$(0.48,0.08)/0.56=(0.857,0.143)$，第二位置为$(0.12,0.36)/0.48=(0.25,0.75)$：第一词元主要依据A，第二词元主要依据B。后验的变化来自生成概率，不是新的检索调用。

训练用问题与参考答案对，最小化式\eqref{eq:nlp-dialogue-rag}中相应边缘概率的负对数。由于$\pi_z$出现在似然中，梯度同样经过$p_\eta(z\mid x)$流向问题编码器，机制与式\eqref{eq:nlp-dialogue-realm-grad}相同：让答案更可能的片段得到更高的检索分数。与REALM不同，RAG在微调中固定片段编码器与索引，只更新问题编码器和生成器，因而不必刷新索引；代价是片段表示不能随任务调整。解码时，RAG-Token的每步分布本身就是对片段求和后的分布，可直接接束搜索；RAG-Sequence的似然不能按词元分解，需要先在每个片段条件下分别束搜索得到候选，再把每个候选在所有片段下重新打分求和，或者只对已出现在某片段束中的候选求和以节省计算。

RAG的输出不附带来源。生成器可能凭参数知识给“14:00”较高概率，即使所读片段只有B；片段后验也只说明生成时依据了哪个片段，不说明该片段支持了哪句话。引用需要另外的输出约定和核验（第\ref{sec:nlp-dialogue-citation}小节）。

\subsubsection{FiD多片段融合生成}
\label{sec:nlp-dialogue-fid}

RAG的每个生成分支只直接读到一个片段，片段之间的信息只在概率加权时汇合。若回答必须同时用到A的活动时间和B的报名截止，更直接的做法是让生成器一次看到全部片段。\term{解码器融合}（fusion-in-decoder, FiD）把这件事放在编码器—解码器结构里完成：每个问题—片段对独立编码，全部编码结果拼接后交给同一个解码器\scite{nlp-dialogue-izacard2021}。编码器—解码器与交叉注意力的计算见第\ref{sec:nlp-generation-encoder-decoder}小节。

具体地，对第$j$个片段构造输入“question: 问题 title: 标题 context: 片段”，经编码器得$\bm H_j\in\mathbb R^{\ell\times d}$，再沿序列维度拼接：
\begin{equation}
 \bm H_j=E_\theta(q,z_j),\qquad
 p(y\mid q,Z_k)=
 \prod_i p_\theta\bigl(y_i\mid y_{<i},[\bm H_1;\ldots;\bm H_k]\bigr),
 \label{eq:nlp-dialogue-fid}
\end{equation}
其中$[\bm H_1;\ldots;\bm H_k]\in\mathbb R^{k\ell\times d}$是解码器交叉注意力的记忆。解码器每步可以同时注意所有片段的所有位置，一个词元能依据多个片段生成。这里没有检索概率$\pi_z$，也没有对片段求和；片段的相对重要性由交叉注意力权重隐式决定。原方法以T5初始化，用BM25或DPR检索，用参考答案的词元负对数似然训练；检索器固定，不与阅读器联合训练。

编码阶段的代价是FiD可以处理大量片段的原因。设每个问题—片段输入长为$\ell$，片段数为$k$，只看自注意力的位置对数，独立编码约为$k\ell^2$，把所有片段直接拼成一条长输入则约为$(k\ell)^2$。在$k=100$、$\ell=250$的教学设置下，前者为$6.25\times10^{6}$，后者为$6.25\times10^{8}$，相差$k$倍。解码器的交叉注意力仍要访问全部$k\ell$个位置，这部分代价随$k$线性增长。原文报告，片段数从10增加到100时，答案准确率持续上升，说明生成器确实能从更多片段中汇集证据。

与RAG相比，FiD的融合发生在解码器内部，适合需要组合多个片段的回答，也不必为每个片段各解码一次；但它失去了显式的片段后验，无法直接说出哪个片段贡献了哪个词元，而且片段在编码阶段彼此不可见。Atlas正是利用FiD解码器对各片段的注意力或困惑度信号，反过来训练检索器（第\ref{sec:nlp-dialogue-joint-retriever}小节）。

\subsubsection{查询改写与假设文档检索}
\label{sec:nlp-dialogue-query-rewrite}

RAG与FiD都直接用原问题检索。用户的问题往往口语化、省略对象，或者与答案文档的写法相差很远：“我是外校的，明天下午能直接过去吗”与片段B“校外来宾须在4月16日18:00前提交报名信息”几乎没有共同词，语义也隔了一层推断。两种方法在检索之前修改查询。

\term{查询改写}（query rewriting）让一个模型先把问题改写成更适合检索的查询，再检索、阅读。Ma等提出的“改写—检索—阅读”框架把改写器作为可训练部件，阅读器是冻结的大语言模型，检索器是网络搜索引擎\scite{nlp-qa-ma2023rewrite}。改写器是T5-large，分两步训练。预热阶段，提示大语言模型改写训练集问题，只保留改写后阅读器答对的样本$\{(x,\tilde x)\}$，以标准对数似然微调改写器。强化学习阶段，把改写器看作策略，每生成一个查询词元是一步动作；改写完成后依次检索、阅读，以阅读器答案的精确匹配、F1以及检索结果是否含参考答案（含则$+1$，否则$-1$）的加权和作为奖励$R_{\mathrm{lm}}$，并减去与初始改写器的KL惩罚，
\begin{equation}
  R(x,\tilde x)=R_{\mathrm{lm}}(\tilde x,y)-\beta\,\mathrm{KL}\bigl(\pi_\theta\,\Vert\,\pi_0\bigr),
  \label{eq:nlp-dialogue-rewrite-reward}
\end{equation}
用PPO优化。PPO的目标与优势估计见第\ref{sec:nlp-align-ppo}小节，这里的区别只在于策略输出的是查询而不是回答，奖励来自下游检索与阅读的结果。这样，改写器不必知道“好查询”长什么样，只需从阅读器能否答对中学习；代价是每个训练样本都要真正调用一次检索与阅读。

另一种做法不改写问题，而是先猜答案。HyDE（hypothetical document embeddings）用指令跟随的大语言模型针对问题写一份假设文档，再用无监督对比学习得到的稠密编码器$f$（Contriever）编码它，以其向量检索真实文档\scite{rs-s30}。若生成$N$份假设文档$\hat d_1,\ldots,\hat d_N$，查询向量取它们与问题向量的平均：
\begin{equation}
  \hat{\bm v}_q=\frac{1}{N+1}\Bigl[\sum_{k=1}^{N}f(\hat d_k)+f(q)\Bigr],
  \qquad s(q,d)=\hat{\bm v}_q^{\top}f(d).
  \label{eq:nlp-dialogue-hyde}
\end{equation}
思路是：答案文档与答案文档之间的相似，比问题与答案文档之间的相似更容易被无监督编码器捕捉。对上面的问题，模型可能写出“外校人员参加讲座一般需提前一天在线报名，凭确认信息入场”。这段文字的具体日期与流程可能都不对，但它在写法上与B相近，向量也更接近B。假设文档只用于检索，不进入证据，回答仍以检索到的真实片段为依据；其中的错误事实因此只影响检索结果，不直接成为答案。

两种方法的位置相同，代价不同。查询改写需要训练信号与在线调用，适合有问答对的任务；HyDE无需任何相关性标注，每次查询多一次生成调用。若语言模型对问题领域完全陌生，假设文档会偏离方向，检索也随之偏离。

\subsubsection{检索结果重排}
\label{sec:nlp-dialogue-rerank}

双编码器让问题与片段在编码时互不可见，只在最后做一次内积，适合在全文中召回，却难以分辨“原定4月17日”与“现改至4月18日”这类只差关系词的片段。第\ref{sec:nlp-classification-rerank}小节的交叉编码器把两者拼接后联合编码，每一层都让问题与片段交互，分辨力更高；代价是片段表示不能预先计算，每个候选都要做一次完整编码。问答系统因此采用“初检取前$K$项、重排保留前$k$项”的两阶段结构。monoBERT是这一做法的代表\scite{rs-s20}：输入\code{[CLS] 问题 [SEP] 片段 [SEP]}，相关分为$s(q,z)=\sigma(\bm w^{\top}\bm h_{\mathrm C}+b)$，以初检候选中的正例与非正例训练二元交叉熵；成对、列表式目标与排序指标见第\ref{sec:nlp-classification-rerank}小节。

问答中的特殊之处在于训练样本的来源。负例应取自初检排名靠前的非答案片段，使重排器学习初检分不清的情形；正例常由“片段含答案字符串”自动构造，同名或旧版本片段可能被误标为正例，需要用答案位置或版本信息过滤。

重排的计算量可以直接估算。初检取$K=100$、重排后保留$k=5$，每个问题要做100次长为$|q|+|z|$的编码，而双编码器在线只编码问题一次；换来的是送入生成器的前5项更准确。重排还能把最相关的片段放到开头或紧挨问题的位置，间接缓解长上下文中的位置效应（第\ref{sec:nlp-dialogue-long-context}小节）。它只能调整已召回候选的次序，初检遗漏的片段无法找回；候选中没有合适片段时，仍需第\ref{sec:nlp-dialogue-bert-span}小节那样的分数阈值决定弃答。

\subsubsection{迭代检索与多跳问答}
\label{sec:nlp-dialogue-iterative}

\term{多跳问答}（multi-hop question answering）的答案需要串联多条事实，后一条查什么取决于前一条查到了什么。HotpotQA为这类问题标注了支持事实，使答案与证据组合可以分别评价\scite{yang2018hotpotqa}。看一个教学问题：“4月18日读书会所在的楼由哪个单位管理？”另设片段D：“海棠楼由信息学院负责管理和借用审批。”用原问题检索，与“读书会”“4月18日”重合的是A；D中没有这些词，第一次检索很难找到它。只有读过A、知道地点是海棠楼，才知道下一步该查“海棠楼由谁管理”。

\term{迭代检索}（iterative retrieval）让检索随推理推进。IRCoT把思维链的生成与检索交替进行\scite{rs-s40}：先用问题检索$K$个片段；然后反复执行两步，推理步在问题、已收集片段和已有推理句的条件下只生成下一句思维链，检索步以这一句为查询再取$K$个片段并入已收集集合；当生成的句子中出现“答案是”或步数达到上限时停止，见算法\ref{alg:nlp-dialogue-ircot}。思维链提示本身见第\ref{sec:nlp-execution-cot}小节。原文用BM25检索，步数上限为8，已收集片段总数上限为15，以便在上下文中留出示例的位置。

\begin{algorithm}[IRCoT：推理与检索交替]
  \label{alg:nlp-dialogue-ircot}
  \begin{algorithmic}[1]
    \Require 问题$q$，检索器$\mathcal R$，语言模型$M$，每步检索数$K$，步数上限$T_{\max}$，片段总数上限$C_{\max}$
    \Ensure 已收集片段集合$\mathcal D$与思维链$r_1,\ldots,r_t$
    \State $\mathcal D\gets\mathcal R(q,K)$；$t\gets0$
    \Repeat
      \State $t\gets t+1$
      \State $r_t\gets M$在$(q,\mathcal D,r_1,\ldots,r_{t-1})$条件下生成的下一句，只保留第一句
      \If{$r_t$含“答案是”}
        \State \textbf{break}
      \EndIf
      \State $\mathcal D\gets\mathcal D\cup\mathcal R(r_t,K)$，超出$C_{\max}$时截断
    \Until{$t=T_{\max}$}
    \State \Return $\mathcal D$与$r_1,\ldots,r_t$；由阅读器在$\mathcal D$上作答，或直接取思维链末句的答案
  \end{algorithmic}
\end{algorithm}

在教学问题上，第一次检索得到A。推理步生成$r_1$“读书会在海棠楼201举行。”；以$r_1$为查询，“海棠楼”与D重合，D被检索进来。第二个推理步读到D，生成$r_2$“海棠楼由信息学院管理，所以答案是信息学院。”，循环停止。检索的查询从原问题变成了推理中刚出现的中间实体，这正是一次性检索无法做到的。

迭代检索的收益依赖于中间推理句是否正确。若$r_1$把地点写错，下一次检索就会沿错误方向展开，而且这种错误不会被后续步骤自动纠正。每多一步，就多一次语言模型调用和一次检索，延迟随步数线性增长。步数上限与片段上限是必须预先设定的预算；达到上限仍未得到答案时，应把剩余子问题作为未解决项报告，而不是继续搜索直到碰巧得到一个答案。

\subsubsection{自适应检索与反思标记：Self-RAG}
\label{sec:nlp-dialogue-self-rag}

前面的方法对每个问题都检索固定次数。可是“写一段读书会的欢迎词”几乎不需要检索，而“读书会改到哪天”必须检索；同一回答中，有的句子陈述事实，有的只是过渡。总是检索浪费计算，也可能把无关片段引入上下文；从不检索又会让事实性陈述失去依据。\term{自适应检索}（adaptive retrieval）让模型自己决定何时检索。Self-RAG进一步让模型在生成过程中评判检索结果与自己的输出\scite{nlp-qa-asai2024selfrag}。

Self-RAG在词表中加入四类\term{反思标记}（reflection token）。\code{Retrieve}在给定输入或已有输出时判断是否需要检索，取值为是、否或继续使用已有片段；\code{IsRel}判断片段$d$对输入是否相关；\code{IsSup}判断当前输出段是否被$d$完全支持、部分支持或不被支持；\code{IsUse}给整个回答的有用程度打1到5分。生成以句子为段，每段之前可以输出\code{Retrieve}，检索后在片段之后输出\code{IsRel}，生成该段后输出\code{IsSup}，回答结束时输出\code{IsUse}。

训练分两步，都是普通的语言模型目标。第一步训练评判模型$C$：从训练集中抽样，用类型专属的指令提示GPT-4为四类反思标记分别打标签，每类收集数千到两万条，再以条件语言模型目标$\max_C\mathbb E\log p_C(r\mid x,y)$训练一个Llama 2-7B评判模型。第二步构造生成器的训练数据：对每条输入—输出对$(x,y)$的每个输出段，由$C$判断是否需要检索；需要时检索前$K$个片段，由$C$为每个片段预测\code{IsRel}与\code{IsSup}，选出相关且有支持的片段插入，并在段后写上相应标记；最后由$C$预测\code{IsUse}。生成器$M$在这样增强过的语料上以下一词元目标$\max_M\mathbb E\log p_M(y,r\mid x)$训练，学习同时生成正文与反思标记；插入的片段文本不计入损失。评判在离线标注时完成，训练时不需要像RLHF那样在线调用奖励模型。

推断时，反思标记的概率成为可调的控制量。当\code{Retrieve=是}在该类标记内归一化后的概率超过阈值$\delta$时触发检索。检索到$K$个片段后，生成器对每个片段并行生成一个候选段$y_t$，按下式打分，并以段为单位做束宽为$B$的束搜索：
\begin{equation}
  f(y_t,d)=p(y_t\mid x,d,y_{<t})+\sum_{G\in\{\mathrm{Rel},\mathrm{Sup},\mathrm{Use}\}}w^{G}s_t^{G}.
  \label{eq:nlp-dialogue-selfrag}
\end{equation}
其中$s_t^{G}$由第$G$类标记的概率在该类内归一化得到：$s^{\mathrm{Rel}}$为“相关”的概率；$s^{\mathrm{Sup}}$为“完全支持”的概率加上0.5倍“部分支持”的概率；$s^{\mathrm{Use}}$把1至5分依次赋权$-1,-0.5,0,0.5,1$后求期望。权重$w^{G}$在推断时设定：重视证据支持就加大$w^{\mathrm{Sup}}$，开放写作则降低检索阈值的作用、更看重有用性。原文默认取$w^{\mathrm{Rel}}=w^{\mathrm{Sup}}=1$、$w^{\mathrm{Use}}=0.5$，段级束宽为2。

用一个教学例子看打分。问题是“读书会改到哪天”，检索到A与C。以A为条件的候选段“读书会改至4月18日（周五）。”生成概率0.6，\code{IsRel}中“相关”占0.9，\code{IsSup}中完全、部分支持分别占0.8与0.1，\code{IsUse}的期望权重为0.75，得分为$0.6+0.9+(0.8+0.05)+0.5\times0.75=2.725$。以C为条件的候选段“读书会仍在4月17日举行。”生成概率0.5，相关0.2，完全、部分支持为0.1、0.1，有用性期望0.25，得分为$0.5+0.2+0.15+0.125=0.975$。束搜索保留以A为依据的段。若生成器只看生成概率，两个候选的差距小得多；反思标记把“片段是否相关、陈述是否被支持”纳入了选择。

Self-RAG把检索决策、相关性判断与支持判断都内化为生成器的输出，推断时不需要另外的评判模型，并能通过阈值与权重在事实性和创造性之间调节。代价有三项：训练数据依赖强模型的离线标注，标注偏差会被生成器学去；每个需要检索的段要为$K$个片段各生成一个候选；反思标记是模型对自己的评价，\code{IsSup}给出高分并不等于外部核验通过，第\ref{sec:nlp-dialogue-verification}节的核验仍然需要。

\subsubsection{带引用的回答生成}
\label{sec:nlp-dialogue-citation}

前述方法得到的回答即使正确，读者也难以核对每句话依据了什么。带引用的回答生成要求回答中的每个陈述后附上所依据片段的编号，例如“活动于4月18日14:00在海棠楼201举行〔A〕”。最直接的实现是提示：把检索到的片段编号后放入上下文，在示例中示范“陈述后加编号”的格式。ALCE基准把这一任务系统化，并给出可自动计算的引用质量指标\scite{nlp-qa-gao2023alce}。

引用质量用自然语言推断模型判断。设回答由陈述$s_1,\ldots,s_n$组成，陈述$s_i$所附的引用集合为$C_i$；$\phi(\text{前提},\text{假设})\in\{0,1\}$是一个蕴含判别器，前提蕴含假设时输出1，ALCE用的是T5-11B规模的TRUE模型，蕴含判断的基本概念见第\ref{sec:nlp-classification-nli}小节。\term{引用召回率}（citation recall）检查每个陈述是否被其引用整体支持：
\begin{equation}
  \mathrm{rec}(s_i)=\mathbb 1\bigl[C_i\neq\varnothing\ \wedge\ \phi(\operatorname{concat}(C_i),s_i)=1\bigr],
  \qquad
  \mathrm{CitRec}=\frac1n\sum_{i=1}^{n}\mathrm{rec}(s_i).
  \label{eq:nlp-dialogue-cite-recall}
\end{equation}
\term{引用精确率}（citation precision）检查每条引用是否多余。引用$c\in C_i$称为不相关，当且仅当它单独不能支持$s_i$，且去掉它后其余引用仍能支持$s_i$，即$\phi(c,s_i)=0$且$\phi(\operatorname{concat}(C_i\setminus\{c\}),s_i)=1$。一条引用的精确度为1，当且仅当所在陈述的召回为1且它不是不相关引用；引用精确率是全部引用上的平均。召回为1被设为前提，否则条件(b)无从判断。

用教学回答演算。回答为“活动于4月18日14:00在海棠楼201举行〔A〕〔C〕。校外来宾须在4月16日18:00前报名〔A〕。”第一个陈述引用A与C，A蕴含它，召回为1；C单独不支持它，去掉C后A仍支持，所以C是不相关引用，A的精确度为1。第二个陈述只引用A，A不含报名条件，召回为0，A的精确度也为0。于是引用召回率为$1/2$，引用精确率为$1/3$。正确的写法是第二句引用B。这个例子也说明了三种情况的差别：答案正确且引用正确；答案正确但引用错误，例如第二句；答案正确但没有引用，召回同样为0。三者在答案准确率上无法区分，引用指标把它们分开。

引用指标衡量的是可追溯性，不是正确性：一个错误陈述若恰好被某个错误片段蕴含，召回仍为1，所以答案的正确性要另行评价。蕴含判别器本身也会出错，跨句推理、数字与日期比较是它的薄弱处；ALCE也指出，这套判定忽略了一条引用部分支持陈述的情形。

\subsubsection{长上下文与检索的比较}
\label{sec:nlp-dialogue-long-context}

上下文窗口扩展到数十万词元以后，一个自然的问题是：能否不检索，把全部资料直接放进上下文？这样可以避免检索遗漏，也省去切分、索引与重排。但它带来两个新问题。

第一是位置效应。第\ref{sec:nlp-inference-lost-middle}小节的控制实验表明，含答案文档处在长输入中间时最难被利用\cite{liu2024middle}；可容纳的长度不等于被有效利用的长度，资料越多，答案落在中间的机会越大。

第二是成本。预填充阶段每个输入词元约需$2N$次浮点运算（$N$为参数量，见第\ref{sec:nlp-inference-autoregression}小节），注意力部分还随长度平方增长，键值缓存随长度线性增长。以教学设定计：资料共20万词元，检索增强只送入5个各300词元的片段，即1500词元，仅线性层部分的预填充计算就相差约133倍，且这一代价在每次提问时都要重新支付，除非资料前缀可以缓存复用。

已有比较研究给出的结论与设置有关。Xu等在相同模型上比较发现，4K上下文的模型配合简单检索，可以达到经位置插值扩展到16K并微调后的模型的水平，而检索对扩展到32K的模型仍有帮助\scite{nlp-qa-xu2024retrieval-long}。Li等在资源充足时观察到长上下文模型的平均效果优于检索增强，但后者成本低得多；据此提出的Self-Route先让模型判断检索到的片段能否回答问题，能则直接作答，不能再转用全文长上下文，以较小的成本接近长上下文的效果\scite{nlp-qa-li2024selfroute}。表\ref{tab:nlp-dialogue-long-context}汇总两条路线的差别；长上下文结构本身见第\ref{sec:nlp-llm-long-context}节。

\begin{table}[htbp]
  \centering
  \small
  \begin{tabular}{@{}lll@{}}
    \toprule
    比较项 & 检索增强生成 & 长上下文直接阅读\\
    \midrule
    证据召回 & 受检索器限制，遗漏无法补回 & 资料全部可见\\
    位置敏感性 & 片段少，可经重排放到首尾 & 中间位置利用较差\\
    每次提问成本 & 只处理少量片段 & 随资料长度增长\\
    资料更新 & 重建或增量更新索引 & 直接替换输入\\
    引用粒度 & 片段编号天然可用 & 需另行定位来源\\
    \bottomrule
  \end{tabular}
  \caption{检索增强生成与长上下文直接阅读的比较。两者可以组合：先检索，判断片段不足时再改读全文。}
  \label{tab:nlp-dialogue-long-context}
\end{table}

\subsection{结构化证据问答}
\label{sec:nlp-dialogue-structured-qa}

证据若存于数据库、知识图谱或表格，答案常需要筛选、连接或聚合，“截止前提交的校外报名有几条、合计几人”这类问题无法从任何一个文本跨度中读出。一条路线把问题转成可执行查询，执行结果即答案，查询本身就是可复核的推导；它要求查询语言覆盖问题，且需要模式链接与解析。另一条路线如TAPAS直接在表格上选单元格并预测聚合算子，不必写出查询，训练可只用答案监督，但只支持少数几种聚合。两条路线都在固定模式、数据快照和只读权限下工作，并保存查询、所选单元格与计算依据。

\subsubsection{SQL与图查询的执行式问答}
\label{sec:nlp-dialogue-execution-qa}

设有额外构造的报名表\code{registration}，字段为活动编号、报名记录编号、参加者编号、身份类别、提交时刻和随行人数。四条记录为：R1，校外，4月16日17:30，2人；R2，校外，4月16日18:00，1人；R3，校内，4月17日09:00，1人；R4，校外，4月16日16:00，3人。它们都属于\code{ML-0418}，时区统一为北京时间。问题是“按时提交的校外报名有几条，合计几人”；提交记录不等于已批准的资格。

从问题到查询的解析属于文本到SQL，模式链接、关系感知编码与大语言模型方法见第\ref{sec:nlp-analysis-text-to-sql}小节。问答一侧要处理的是解析之后的步骤。模式链接把“校外”对应到身份字段，把“按时”对应到通知中“截止前”的严格条件，把“几条”和“几人”分别对应到记录数与人数之和；解析出的候选查询还要检查字段、类型、只读性和条件，再绑定参数执行：
\begin{codeblock}[language=SQL,label={code:nlp-dialogue-sql},unbreakable]{报名统计的只读查询}
SELECT COUNT(*) AS records,
       SUM(party_size) AS people
FROM registration
WHERE event_id = :event
  AND visitor_type = :external
  AND submitted_at < :deadline;
\end{codeblock}
绑定的活动为\code{ML-0418}，截止为2025年4月16日18:00。筛选留下R1、R4，结果为2条、5人。回复应为“按给定快照，截止前提交的校外报名有2条，登记人数合计5人”，并保留行编号R1、R4作为证据。若把条件误写成$\leq$，会多计R2，得到3条、6人：查询执行无错误，却没有回答原问题。

执行结果需要按答案协议解释。空结果不总是零：\code{COUNT}在无匹配行时为0，空集上的\code{SUM}在多数数据库中为\code{NULL}，答案应表达为“无符合记录”还是数值零，要事先规定。两条不等价的查询可能在当前表上偶然同值，用截止恰为18:00、同名不同编号、多表连接产生重复行等反例数据复核，比只比较一次执行结果更能暴露语义错误。

图数据可把活动、报名记录和参加者表示为节点，以“属于活动”“提交者”等边连接。问题到图路径的映射同样要绑定活动身份与时间条件，再计数通过筛选的报名节点；跨边连接时保存节点与边的编号作为路径证据。若一个人关联两条报名记录，计数报名节点与计数不同参加者的结果不同，问题要求哪一种必须在解析时确定。

\subsubsection{TAPAS单元格选择与聚合}
\label{sec:nlp-dialogue-tapas}

执行式问答训练解析器需要问题—查询对，而这类标注昂贵；许多表格问答数据只给出答案。TAPAS绕开查询：模型直接指出哪些单元格参与回答，再预测对它们施加什么聚合，训练只需答案\scite{nlp-dialogue-herzig2020}。

输入是\code{[CLS] 问题 [SEP]}后接按行展开的表格词元。每个词元的嵌入在BERT的词元、位置、段嵌入之外，再加三类与表格有关的嵌入：列编号与行编号（问题词元取0），以及数值排序编号，可解析为数值或日期的列按大小排序，最小者为1，不可比较者为0；对话场景还加一个“是否为上一问答案”的标记。排序编号让模型不必从子词中读出数值大小，就能回答“最多”“最早”这类问题。

编码之后有两个输出头。单元格选择头先对每个词元的末层表示做线性变换得logit，单元格的logit取其所含词元logit的平均，经sigmoid得选择概率$p_s(c)$。原模型还加入单列选择的归纳偏置：对每列取其单元格表示的平均，经另一个线性层得列logit，再加一个表示“不选任何列”的空列，在列之间做softmax得$p_{\mathrm{col}}(\mathit{co})$；被选列之外的单元格概率置为0。聚合头对\code{[CLS]}表示做线性变换与softmax，得$p_a(\mathit{op})$，$\mathit{op}\in\{\text{NONE},\text{COUNT},\text{SUM},\text{AVERAGE}\}$。推断时取最可能的聚合算子，选出概率超过0.5的单元格并执行。

训练按答案类型分两种情况。若答案能映射到单元格集合$C$，就选含$C$中单元格最多的列作为目标列$\mathit{col}$，损失由列选择、列内单元格选择与聚合三项组成：
\begin{equation}
\begin{aligned}
  J_{\mathrm{CS}}={}&
  \frac{1}{|\mathrm{Cols}|}\sum_{\mathit{co}}\mathrm{CE}\bigl(p_{\mathrm{col}}(\mathit{co}),\mathbb 1_{\mathit{co}=\mathit{col}}\bigr)\\
  &+\frac{1}{|\mathrm{Cells}(\mathit{col})|}\sum_{c\in\mathrm{Cells}(\mathit{col})}\mathrm{CE}\bigl(p_s(c),\mathbb 1_{c\in C}\bigr)
  -\alpha\log p_a(\text{NONE}),
\end{aligned}
  \label{eq:nlp-dialogue-tapas-cs}
\end{equation}
其中CE是二元交叉熵，$\alpha$是权重；这类样本不需要聚合，故以NONE为聚合目标。若答案是表中没有的标量$s$，哪些单元格、哪个算子都不知道，而能执行出$s$的组合可能很多，例如$s=5$时任意5个单元格的计数都对。TAPAS不去搜索这些组合，而是对每个算子构造可微的软执行：软计数为$\sum_c p_s(c)$，软求和为$\sum_c p_s(c)\,T[c]$，软平均取两者之比，$T[c]$为单元格数值。期望结果按去掉NONE后重新归一化的算子概率$\hat p_a$加权：
\begin{equation}
  s_{\mathrm{pred}}=\sum_{\mathit{op}\neq\text{NONE}}\hat p_a(\mathit{op})\operatorname{compute}(\mathit{op},p_s,T),\quad
  J_{\mathrm{scalar}}=
  \begin{cases}
    \tfrac12 a^2, & a\leq\delta,\\
    \delta a-\tfrac12\delta^2, & a>\delta,
  \end{cases}
  \label{eq:nlp-dialogue-tapas-scalar}
\end{equation}
其中$a=|s_{\mathrm{pred}}-s|$，$J_{\mathrm{scalar}}$是\term{Huber损失}（Huber loss），误差小时为平方、误差大时为线性，比平方损失对离群样本更稳定。总损失为$J_{\mathrm{SA}}=-\log\sum_{\mathit{op}\neq\text{NONE}}p_a(\mathit{op})+\beta J_{\mathrm{scalar}}$，第一项把概率质量推离NONE；$J_{\mathrm{scalar}}$超过设定的截断值时整条样本被忽略，因为这类样本多为离群点。

固定上一小节的表，问“截止前校外报名合计几人”，答案为表中没有的标量5。设人数列被选中，四个人数单元格的选择概率为$(0.95,0.10,0.05,0.90)$，聚合概率在去掉NONE后为SUM 0.9、COUNT 0.1、AVERAGE 0。软求和为$0.95\times2+0.10\times1+0.05\times1+0.90\times3=4.75$，软计数为$2.0$，期望结果$s_{\mathrm{pred}}=0.9\times4.75+0.1\times2.0=4.475$，误差$a=0.525$；取教学值$\delta=1$，$J_{\mathrm{scalar}}=0.5\times0.525^2=0.138$。推断时阈值0.5留下R1的2与R4的3，SUM执行为离散的5。训练时的软结果与推断时的离散结果不同，这是为了让答案监督可微而付出的近似。

TAPAS的优势是只用答案监督、不写查询，选出的单元格本身就是证据；代价是单列选择的偏置与三种聚合限定了能回答的问题，跨列比较、多表连接与多步计算都不在其中。弱监督还可能让一组错误的单元格偶然算出同样的总数，评价时除答案外应检查所选单元格与单位。

\subsection{混合证据问答}
\label{sec:nlp-dialogue-hybrid-qa}

财报这类材料中，正文限定对象与统计口径，数值却在表格里，问题还常要多步计算。此时仅找出数字不足以形成答案，还需要把对象、时间、单位和比较范围对齐。TAGOP在一次证据选择后套用预定义算子，适合单步计算；FinQANet生成多步程序，能表达更长的计算链，代价是需要程序监督与执行器。本节只处理已可用的文本与结构化表格，不涉及页面图像与版面恢复。

表\ref{tab:nlp-dialogue-operands}是独立的教学统计，与共同通知的报名记录无关。附带正文为：“比较甲校读书会的年度到场人数；两年口径相同，均按去重人数统计，不含旁听记录。”问题是“甲校2025年到场人数比2024年增长百分之几”。

\begin{table}[htbp]
  \centering
  \begin{tabular}{llll}
    \toprule
    单元格编号 & 对象 & 年份 & 到场人数\\
    \midrule
    T1-R1-C3 & 甲校读书会 & 2024 & 120人\\
    T1-R2-C3 & 甲校读书会 & 2025 & 150人\\
    T1-R3-C3 & 乙校读书会 & 2025 & 180人\\
    \bottomrule
  \end{tabular}
  \caption{增长率问题的教学操作数。正文限定甲校、年度和去重口径，表格提供数值及年份；乙校一行只用作对象对齐的反例。}
  \label{tab:nlp-dialogue-operands}
\end{table}

\subsubsection{TAGOP标签选择与符号计算}
\label{sec:nlp-dialogue-tagop}

TAT-QA基准收集了财报中表格与正文混合的问题，TAGOP是随之提出的模型\scite{nlp-dialogue-zhu2021}。它把混合问答拆成四个分类任务：选证据、选算子、定操作数顺序、定尺度，数值计算交给符号执行。问题、按行展开的表格与相关段落依次输入RoBERTa编码器，得各子词表示$\bm h_t$。

证据选择采用IO标注：每个子词独立预测标签，
\begin{equation}
  p^{\mathrm{tag}}_t=\operatorname{softmax}\bigl(\mathrm{FFN}(\bm h_t)\bigr)\in\mathbb R^{2},
  \label{eq:nlp-dialogue-tagop-tag}
\end{equation}
FFN是带GELU激活的两层前馈网络。一个单元格或一个词只要有子词被标为I就算正候选，段落中连续的正词合并为跨度。IO只区分“是否属于证据”，不像BIO那样标注边界起点，因为这里要的是证据集合而非实体边界。算子分类器以\code{[CLS]}表示预测十种算子之一：段落跨度、单元格、多项、求和、计数、平均、乘法、除法、差额与变化率，$p^{\mathrm{op}}=\operatorname{softmax}(\mathrm{FFN}(\bm h_{\mathrm{[CLS]}}))$。除法、差额与变化率对操作数顺序敏感，因此再加一个顺序分类器，输入为选择概率最高的两个词的表示的平均：$p^{\mathrm{order}}=\operatorname{softmax}(\mathrm{FFN}(\operatorname{avg}(\bm h_{t_1},\bm h_{t_2})))$。尺度分类器以$[\bm h_{\mathrm{[CLS]}};\bm h_{\mathrm{tab}};\bm h_{\mathrm{p}}]$为输入，在无、千、百万、十亿与百分比中选择，$\bm h_{\mathrm{tab}}$与$\bm h_{\mathrm{p}}$分别是表格与段落词元表示的平均。训练损失是四项负对数似然之和：
\begin{equation}
\begin{aligned}
  L={}&\mathrm{NLL}(p^{\mathrm{tag}},G^{\mathrm{tag}})+\mathrm{NLL}(p^{\mathrm{op}},G^{\mathrm{op}})\\
   &+\mathrm{NLL}(p^{\mathrm{scale}},G^{\mathrm{scale}})+\mathrm{NLL}(p^{\mathrm{order}},G^{\mathrm{order}}),
\end{aligned}
  \label{eq:nlp-dialogue-tagop-loss}
\end{equation}
各项标签由标注的支持证据、推导式与答案尺度恢复。

在教学问题上，模型应把T1-R2-C3的150与T1-R1-C3的120标为I，并依靠正文确认比较对象是甲校。设两者的选择置信为$0.93$与$0.89$，乙校180的置信为$0.08$；算子分类为变化率，顺序分类确定当年数在前，尺度分类为百分比。执行得
\begin{equation}
 r=\frac{150\ \text{人}-120\ \text{人}}{120\ \text{人}}
   =0.25,\qquad \text{百分比答案}=25\%.
 \label{eq:nlp-dialogue-growth}
\end{equation}
人数单位在比值中消去，尺度头给出的百分比使0.25输出为25\%，而不是0.25\%。若问题改问增加多少人，算子应为差额，结果为30人。乙校180与甲校120也能执行出一个变化率，但回答了错误的对象；对象对齐依赖证据选择是否读懂了正文的限定。

TAGOP把数值计算交给确定的符号操作，避免了让神经网络直接输出数字；但每题只套用一个算子，算子的输入又只取概率最高的两个数，“两个部门先合计、再减去旁听人数、然后求比值”这类多步计算超出了它的表达范围。

\begin{table*}[htbp]
  \centering
  \small
  \begin{tabularx}{\linewidth}{@{}>{\raggedright\arraybackslash}p{24mm}
    >{\raggedright\arraybackslash}X>{\raggedright\arraybackslash}X
    >{\raggedright\arraybackslash}X@{}}
    \toprule
    方法 & 解决的问题 & 核心机制 & 新代价\\
    \midrule
    BiDAF & 问题信息过早压缩为单向量 & 相似矩阵上的双向注意力，逐位置表示流入建模层 & 手工交互层，无预训练，只输出单跨度\\
    BERT阅读器 & 交互层需手工设计 & 问题与材料联合编码，跨度头与空答案阈值 & 长度上限需分窗，仍只能抽取\\
    DPR加阅读器 & 词汇检索漏掉改写表述 & 双编码器内积检索，阅读器兼做片段选择 & 需片段标注，编码器更新需重建索引\\
    REALM & 检索器不知阅读器需要什么 & 片段作潜变量，答案似然的梯度更新检索器 & 异步刷新索引，训练成本高\\
    Atlas & 少样本下联合训练 & 阅读器困惑度蒸馏为检索分布 & 阅读器偏差传给检索器\\
    RAG & 答案受跨度限制 & 生成器条件于片段，按检索概率边缘化 & 每片段一个分支，无来源标注\\
    FiD & 片段间信息只在概率层汇合 & 独立编码、拼接后由一个解码器交叉注意 & 无显式片段后验，片段编码时互不可见\\
    查询改写、HyDE & 问题与文档表述不一致 & 改写器以强化学习优化；假设文档向量检索 & 多一次生成调用；领域陌生时偏离\\
    交叉编码器重排 & 双编码器排序粗 & 问题与片段联合编码打分 & 每个候选一次完整编码\\
    IRCoT & 一次检索不足以多跳 & 推理句作为下一次检索的查询 & 调用次数随步数增长，错误推理带偏检索\\
    Self-RAG & 检索时机固定，结果不经评判 & 反思标记决定检索并评判相关与支持 & 依赖离线标注，每片段一个候选\\
    ALCE式引用 & 回答无法追溯 & 陈述附片段编号，蕴含模型算引用召回与精确率 & 不衡量正确性，受蕴含模型误差影响\\
    执行式问答 & 答案需筛选与聚合 & 解析为SQL或图查询并执行 & 需模式链接与解析，查询语言需覆盖问题\\
    TAPAS & 查询标注昂贵 & 单元格选择加聚合头，软执行用答案监督 & 聚合种类少，软结果与离散执行不一致\\
    TAGOP & 文本与表格混合 & IO标注选证据，分类算子、顺序与尺度 & 每题一个算子，难以多步\\
    FinQANet & 多步数值计算 & 检索事实后生成带步骤引用的程序 & 需程序标注，有步骤与终止错误\\
    \bottomrule
  \end{tabularx}
  \caption{证据问答方法对照。“解决的问题”指该方法针对前一类方法或任务的哪一项困难，“新代价”是它为此付出的计算、标注或表达能力上的代价。}
  \label{tab:nlp-dialogue-evidence-methods}
\end{table*}

\subsubsection{FinQANet程序生成与执行}
\label{sec:nlp-dialogue-finqa}

多步计算需要把中间结果写出来。FinQA数据集为财报问题标注了完整的计算程序，FinQANet随之提出，由检索器和程序生成器两部分组成\scite{nlp-dialogue-chen2021}。

报告全文常超过编码器的长度限制，所以先检索。检索器用模板把表格的每一行转成句子，例如“甲校读书会2024年的到场人数为120人”，把每个候选事实与问题拼接后用BERT类分类器判断是否相关，取前$n$个事实并按其在原报告中的顺序排列，作为生成器的输入。检索器以标注的支持事实为监督。

程序生成器逐个输出程序词元，词元有三个来源：输入中的数字与表格行名，其表示$\bm h_{e_i}$来自预训练编码器；限定语言中的函数名、括号与常数，嵌入$\bm h_{s_i}$随机初始化；表示已完成步骤结果的记忆标记\code{\#0}、\code{\#1}等，嵌入$\bm h_{m_i}$同样随机初始化。三者组成候选词元矩阵$\bm H=[\bm h_{e_i};\bm h_{s_i};\bm h_{m_i}]$。解码器是LSTM，第$T$步的输出$\bm h_T$分别注意输入与已解码历史，得注意力向量$\bm{att}_p$与$\bm{att}_h$，再组合成上下文向量；候选词元也按输入注意力调整，最终在候选上做softmax：
\begin{equation}
\begin{aligned}
  \bm c_T&=\bm W_c[\bm{att}_p;\bm{att}_h;\bm h_T],\qquad
  \bm H'_T=\bm W_h[\bm H;\bm H\odot\bm{att}'_p],\\
  \bm w_T&=\operatorname{softmax}(\bm H'_T\bm c_T).
\end{aligned}
  \label{eq:nlp-dialogue-finqa}
\end{equation}
每当生成一个操作的右括号，就用$\bm a_T=\bm W_a[\bm{att}_p;\bm{att}_h;\bm h_T]$更新当前步骤记忆标记的嵌入，后续步骤便能以\code{\#0}引用它。记忆标记的嵌入只帮助生成器选择引用，实际数值仍由执行器计算。推断时按限定语言的文法在每步屏蔽非法词元，保证程序结构合法；训练以参考程序的词元负对数似然为目标，生成器的输入用检索器在训练集上的结果，若漏掉了标注事实就补上。

固定表\ref{tab:nlp-dialogue-operands}与相同正文，教学程序为：
\begin{codeblock}[language=Python,label={code:nlp-dialogue-finqa}]{从差额到增长率的限定程序}
subtract(150, 120)   # step 0: 30 persons
divide(#0, 120)      # step 1: 0.25
\end{codeblock}
这是限定语言，不是可直接执行的Python。150绑定T1-R2-C3，两个120都绑定T1-R1-C3，\code{\#0}只能引用已完成的第0步。执行器先得30人，再以120人为分母得0.25，格式化为25\%并附上正文限定与单元格来源。若生成\code{divide(120,\#0)}，程序合法但结果为4；若把120换成碰巧同值的其他部门记录，执行答案不变，证据绑定却已错误。FinQA因此同时报告执行准确率与程序准确率，程序准确率允许数学上等价的程序。

相对于TAGOP内置的变化率算子，FinQANet把同一计算展开成两个可组合的步骤，能表达限定语言支持的任意多步程序；代价是需要程序级标注，并引入步骤选择、引用与终止的错误。文法掩码能阻止参数位输出非法函数，却不能判断某个数字是否来自正确的年份。通用的程序生成、测试与修订循环见第\ref{chap:nlp-solving-and-execution}章。

表\ref{tab:nlp-dialogue-evidence-methods}汇总本节方法。从上到下，每一行大致回应上一行留下的问题：交互从手工设计走向预训练，检索从词汇走向稠密并与阅读器联合训练，答案从跨度走向生成，再由查询、排序、检索时机与引用逐环节改进；证据形式变为表格与混合材料时，答案的形成又从抽取转为执行。

\FloatBarrier

\section{主张与回答核验}
\label{sec:nlp-dialogue-verification}

上一节的系统根据问题寻找证据并形成答案；本节反过来，把一段已经写成的回答当作待查对象，逐项判断其中的陈述有无依据。这项工作称为\term{事实核验}（fact verification）：给定待查文本、允许使用的知识源和时点，找出相关证据，判断每条陈述被支持、被反驳还是信息不足，再据此修订回答。第\ref{sec:nlp-generation-faithfulness}小节检查的是生成结果与给定输入是否一致，本节的证据则需要自行检索，而且回答中的陈述可能超出任何单份输入。

核验面临四个挑战。其一，主张粒度难定：一句话常同时包含日期、地点和人数，整句判断会让一处错误掩盖其余正确内容，切得过细又会丢失命题之间的绑定。其二，证据需要检索：资料库大，一个主张的证据可能分散在几份文档的几个句子里。其三，反驳与信息不足要区分：找不到证据不等于主张为假。其四，修订不能引入新错误：改正一处时，应尽量保留原回答中已获支持的部分。

针对这些挑战，方法沿三条线发展。FEVER数据集\cite{nlp-intro-thorne2018}确立了“文档检索—句子选择—三分类推断”的流水线，解决主张级别的证据定位与判定；FActScore和SAFE把长回答拆成原子事实，分别对照固定知识源和网络搜索逐条判定，解决长回答的粒度问题；RARR和CoVe则在核验之后修订回答，前者依靠外部检索做最小编辑，后者依靠模型自己提出并独立回答验证问题。本节沿用章首通知，并另设一份教学签到统计S：“ML-0418实际到场人数为120人，统计时点为活动结束后。”待查回答为“读书会4月18日在海棠楼201举行，现场有150人”。

\subsection{主张拆解与证据定位}
\label{sec:nlp-dialogue-claims}

待查回答中的日期和地点可由通知核对，人数却需要另一份资料。若把整句作为一个单位，就只能得出“部分正确”这样无法操作的结论。拆解的目标是得到可以分别检索、分别判断的命题，并保持对象、否定、数量词和时间条件不变；例如“至少150人”不能拆成“150人”，“不是所有校外来宾都可直接参加”也不能变成“所有校外来宾都不可参加”。句式稳定的业务文本可以复用第\ref{chap:nlp-information-extraction}章的事件抽取，先得到$(\text{ML-0418},\text{日期}=4/18,\text{地点}=\text{海棠楼201},\text{人数}=150)$这样的记录，再把每个字段还原为一个命题；开放文本则要依靠语言模型拆解。拆出命题之后，第二步是为每个命题找到证据文档并选出证据句。

\subsubsection{FActScore原子事实分解}
\label{sec:nlp-dialogue-factscore}

长回答的事实性很难用一个“对”或“错”概括：一篇人物传记可能九成内容正确，只错了出生年份。FActScore的做法是先把回答分解为\term{原子事实}（atomic fact），即每条只承载一个可独立检查的信息点的短句，再统计其中被指定知识源支持的比例\scite{min-etal-2023-factscore}。设模型$M$对输入$x$的回答为$y=M_x$，$A_y$是$y$的原子事实集合，$C$是规定的知识源，则
\begin{equation}
 f(y)=\frac{1}{|A_y|}\sum_{a\in A_y}\mathbb{I}[C\text{支持}a],\qquad
 \operatorname{FActScore}(M)=\mathbb{E}_{x}\bigl[f(M_x)\mid M_x\text{作答}\bigr].
 \label{eq:nlp-dialogue-factscore}
\end{equation}
期望只对作答的样本计算，因此原文同时报告作答比例和每个回答的平均事实数。原研究以人物传记为任务、英文维基百科为知识源；自动估计器先用语言模型按少样本提示拆出原子事实，再针对每条事实从知识源检索若干相关段落，提示语言模型判断该事实是否为真。

对本节的待查回答，先把“读书会”绑定\code{ML-0418}，补出年份2025，得到三条原子事实：活动在4月18日举行；活动地点为海棠楼201；现场人数为150人。若知识源只有通知，前两条获得支持，第三条未获支持；若知识源包括签到统计S，第三条被反驳。两种情形下$f(y)=2/3$，但原因不同：前者是缺少资料，后者是资料矛盾。FActScore只区分“支持”与“未获支持”，这一差别要留到第\ref{sec:nlp-dialogue-evidence-relation}小节的三分类中处理。

这一比例依赖拆解质量。分解器漏掉人数条目，分数虚升为$2/2$；把“林宁主持活动，他来自甲校”的“他”错绑到另一人，会对错误对象检索出支持。回答越短、越含糊，越容易得到高分，所以公开的评估工具还提供长度惩罚：事实数少于阈值$\gamma$时把$f(y)$乘以$\exp(1-\gamma/|A_y|)$。FActScore衡量的是精确性，即说出的内容有多少有据，不衡量回答是否完整。

\subsubsection{证据文档检索与句子选择}
\label{sec:nlp-dialogue-evidence-selection}

FActScore默认每条事实对应一个已知条目，可以直接到该条目下找证据。一般的核验没有这个便利：主张可能涉及任意实体，证据分散在大规模文集中。FEVER把这一问题拆成两步检索：先找可能相关的文档，再在文档中选出证据句。两步的划分是出于计算考虑：文档级检索可以用倒排索引快速完成，逐句打分代价高，只能在少量候选文档内进行。

文档检索有两种常见做法。词汇检索把主张$c$与文档$d$表示为一元与二元词的TF-IDF向量$\bm{\phi}(c),\bm{\phi}(d)$，按余弦相似度
\begin{equation}
 \operatorname{score}_{\mathrm{doc}}(c,d)=
 \frac{\bm{\phi}(c)^\top\bm{\phi}(d)}{\lVert\bm{\phi}(c)\rVert_2\lVert\bm{\phi}(d)\rVert_2}
 \label{eq:nlp-dialogue-doc-retrieval}
\end{equation}
取前$k$篇，记为$D_k$；FEVER基线即用此法，取$k=5$。另一种做法利用维基百科类资料“一个实体一篇文档”的结构：从主张中抽取名词短语或实体提及，按第\ref{sec:nlp-extraction-linking}节的实体链接把它们映射到文档标题，直接取回对应文档。实体链接对专名准确，但主张提到的实体若在资料中没有独立条目，就会漏检；两种方法常合并候选。

句子选择对$D_k$中每个句子$x$计算与主张的相关分$s_{\bm\theta}(c,x)$，取前$l$句作为证据集合：
\begin{equation}
 \hat E=\operatorname*{top\text{-}l}_{x\in\mathcal{S}(D_k)}\ s_{\bm\theta}(c,x),
 \label{eq:nlp-dialogue-sentence-selection}
\end{equation}
$\mathcal{S}(D_k)$是候选文档中的全部句子。最简单的$s$同样是TF-IDF余弦；学习式打分器把主张与句子成对输入句对模型（ESIM或BERT交叉编码器，结构见第\ref{sec:nlp-classification-nli}小节），输出一个实数分。训练标签来自标注的证据句：对同一主张，令$x^+$为证据句，$x^-$为同批候选中的非证据句，以间隔为$m$的排序损失
\begin{equation}
 L_{\mathrm{sel}}=\sum_{x^+}\sum_{x^-}
 \max\bigl(0,\ m-s_{\bm\theta}(c,x^+)+s_{\bm\theta}(c,x^-)\bigr)
 \label{eq:nlp-dialogue-selection-loss}
\end{equation}
要求证据句比非证据句至少高出$m$分\scite{nlp-dialogue-hanselowski2018ukp}。也可以把每句当作独立的二分类样本，用二元交叉熵训练；排序损失只关心相对次序，正好对应“取前$l$句”的用法。

选句的好坏要用召回评价，因为漏掉的证据后面无法补回。FEVER为每个可判定主张标注若干组证据，每组是一个足以作出判断的句子集合$G$，有些组只含一句，有些需要多句联合。设$\mathcal{G}_c$为主张$c$的证据组，证据召回率定义为
\begin{equation}
 R_{\mathrm{ev}}=\frac{1}{|\mathcal{C}|}\sum_{c\in\mathcal{C}}
 \mathbb{I}\bigl[\exists G\in\mathcal{G}_c:\ G\subseteq\hat E_c\bigr],
 \label{eq:nlp-dialogue-evidence-recall}
\end{equation}
$\mathcal{C}$为支持或反驳类主张的集合。端到端评价使用FEVER共享任务定义的FEVER分数：只有标签正确，且对支持、反驳类主张至少有一组完整证据落在预测的前五句中，该主张才计为正确\scite{nlp-dialogue-thorne2018shared}

以主张“读书会4月18日在海棠楼201举行”为例。设检索取回通知与一份检修公告，候选句为：$x_1$“原定于4月17日举行，现改至4月18日（周五）14:00”；$x_2$“地点为海棠楼201”；$x_3$“校内师生可直接参加”；$x_4$“海棠楼201正在检修”。教学打分为$s(x_1)=2.1$、$s(x_2)=1.6$、$s(x_4)=1.2$、$s(x_3)=-0.5$。日期与地点需要$x_1,x_2$两句联合支持，证据组为$G=\{x_1,x_2\}$。取$l=2$时$\hat E=\{x_1,x_2\}\supseteq G$，召回成功；取$l=1$时只剩$x_1$，召回失败，后续判定只能得到信息不足。若$m=1$，句对$(x_2,x_4)$的损失为$\max(0,1-1.6+1.2)=0.6$：$x_2$虽然排在前面，但领先不足一个间隔，仍会受到推动。

流水线的主要缺点是错误单向传递：文档检索漏掉的证据，选句无法恢复；选句漏掉的证据，判定阶段只能看到信息不足。实体链接与TF-IDF合并、适当放大$k$与$l$，是用后续阶段的计算换取召回。

\subsection{支持、反驳与信息不足}
\label{sec:nlp-dialogue-evidence-relation}

取得证据后，判定采用三类标签：“支持”表示证据能推出主张，“反驳”表示证据推出与主张不相容的内容，“信息不足”表示现有证据不足以作出前两种判断。三分类与第\ref{sec:nlp-classification-nli}小节的自然语言推断相对应，蕴含、矛盾、中立分别对应支持、反驳和信息不足；核验在此之外增加了证据是否取到、来源是否适用的问题。无关段落不是反证，同一句话被多次转载也不构成独立支持。判定方法有三条路线：学习式句对推断适合语言变化多的证据；规则与结构化比较适合日期、数值等可规范化字段；检索增强的逐事实判定把检索与判断交给语言模型，适合开放领域的长回答。

\subsubsection{可分解注意力的证据判定}
\label{sec:nlp-dialogue-decomposable}

FEVER基线把选出的证据句拼接为前提、主张作为假设，交给第\ref{sec:nlp-classification-decomposable}小节的可分解注意力模型三分类：
\begin{equation}
 p(r\mid c,\hat E)=f_{\mathrm{DA}}\bigl(x_{(1)}\oplus x_{(2)}\oplus\cdots\oplus x_{(l)},\ c\bigr),\qquad
 r\in\{\text{支持},\text{反驳},\text{信息不足}\},
 \label{eq:nlp-dialogue-decomposable}
\end{equation}
$\oplus$表示按选句顺序拼接，$f_{\mathrm{DA}}$内部依次做软对齐、比较与求和聚合。拼接让模型能够把$x_1$中的日期和$x_2$中的地点同时对齐到主张，处理多句证据组。

训练时有一个核验特有的问题：信息不足类主张没有标注证据，但模型推断时总会拿到某些句子。FEVER为这类样本构造输入证据，比较了两种取法：从检索到的最近文档中取句子，或从全部资料中随机取句子\scite{nlp-intro-thorne2018}。基线流水线最终采用前者，因为它更接近测试时“检索到了相关但不充分的句子”的情形；后者过于容易，模型可能只学会“句子与主张无关就判信息不足”，遇到检索失败时行为与训练不一致。

拼接也有代价：句子边界和来源信息被抹平，模型无法报告是哪一句支持了主张。另一种做法先对每句与主张分别做推断，再用注意力加权汇总各句结论，既保留逐句依据，也能联合多句\scite{nlp-dialogue-hanselowski2018ukp}。设模型对三类证据给出教学概率：通知中的改期句对“当前读书会在4月18日举行”为$(0.90,0.04,0.06)$；一份假设的权威更新“当前活动改为4月19日”为$(0.05,0.90,0.05)$；检修公告为$(0.15,0.10,0.75)$。按最大值恢复支持、反驳与信息不足。可分解注意力参数少、可并行，但基础版本不含词序，“甲通知乙”与“乙通知甲”难以区分；后续系统多用BERT类交叉编码器替换判定部件，输入协议不变。

\subsubsection{规则与结构化事实核验}
\label{sec:nlp-dialogue-structured-check}

日期、数值和明确字段可以先规范化再比较，不必交给学习模型。规则的输入是带对象标识、属性、时间和单位的主张记录，以及通过只读查询取得的证据记录。只有主体、属性和适用范围一致时才比较值：相同输出匹配，不相容输出冲突，缺字段或范围未对齐输出信息不足，并附上原因。

主张“报名可以在4月16日18:00提交”经规范化得到提交时刻$t$为2025-04-16 18:00，通知条件为$t<t_{\mathrm{deadline}}$，代入边界时刻为假，因此判为冲突。“到场人数为0.12千人”与S中的“120人”换算到同一单位后匹配；“120组”缺少每组人数，判为信息不足，不能只比较数字120。“原定4月17日”与“现改4月18日”分属原计划日期和当前日期两个属性，保留属性后不冲突；若抽取时只留下两个日期再任选其一，就会把时间更新误判为来源矛盾。

规则精确、可解释，但只覆盖已经编码的比较。“人数增长是因为改到了周五”增加了因果主张，日期和人数分别正确也不能推出它；规则应返回“不在覆盖范围内”，交给有相应证据的其他判定方法。

\subsubsection{检索增强的逐事实判定}
\label{sec:nlp-dialogue-safe}

FActScore对照的是固定知识源，适用于人物传记这类有对应条目的任务；开放问题的长回答涉及的事实没有现成条目，FEVER流水线又需要为每个领域训练检索与判定模型。SAFE（search-augmented factuality evaluator）把整个过程交给语言模型执行：模型拆出事实，为每条事实自己生成搜索查询，读取搜索结果后判断是否被支持\scite{nlp-dialogue-wei2024safe}。它是FActScore的扩展：判定单位仍是原子事实，证据来源从固定知识源换成多步网络搜索。

SAFE的步骤如算法\ref{alg:nlp-dialogue-safe}。拆分之后先做自足化改写，把代词和省略替换为明确对象，使每条事实脱离原回答也能被检索；再判断每条事实与原问题是否相关，无关事实不计入后续统计。检索是多步的：每一步的查询由语言模型根据事实和此前已取得的结果生成，因此第二个查询可以针对第一次结果中缺少的部分。

\begin{algorithm}[SAFE检索增强的逐事实判定]
  \label{alg:nlp-dialogue-safe}
  \begin{algorithmic}[1]
    \Require 问题$x$，长回答$y$，语言模型$\mathrm{LM}$，搜索接口$\mathrm{Search}$，每条事实的查询上限$J$
    \Ensure 支持数$S$，未获支持数$N$，无关数$I$
    \State $A\gets\mathrm{LM}_{\mathrm{split}}(y)$\Comment{拆成单个事实}
    \State $S,N,I\gets0$
    \For{$a\in A$}
      \State $a\gets\mathrm{LM}_{\mathrm{revise}}(a,y)$\Comment{替换代词与省略，使事实自足}
      \If{$\mathrm{LM}_{\mathrm{relevant}}(a,x)$为假}
        \State $I\gets I+1$；\textbf{continue}
      \EndIf
      \State $R\gets\varnothing$
      \For{$j=1,\ldots,J$}
        \State $q_j\gets\mathrm{LM}_{\mathrm{query}}(a,R)$；$R\gets R\cup\mathrm{Search}(q_j)$
      \EndFor
      \If{$\mathrm{LM}_{\mathrm{judge}}(a,R)=\text{支持}$}
        \State $S\gets S+1$
      \Else
        \State $N\gets N+1$
      \EndIf
    \EndFor
    \State \Return $S,N,I$
  \end{algorithmic}
\end{algorithm}

与SAFE同时提出的指标还考虑了回答是否说得够多。精确率$P=S/(S+N)$衡量说出的事实有多少有据；用户期望的事实数记为超参数$K$，召回率$R_K=\min(S/K,1)$衡量有据事实是否达到期望数量，二者合成
\begin{equation}
 F_1@K=
 \begin{cases}
 \dfrac{2P\,R_K}{P+R_K}, & S>0,\\[1ex]
 0, & S=0.
 \end{cases}
 \label{eq:nlp-dialogue-safe-f1}
\end{equation}
它补上了FActScore鼓励少说的缺口：只说一句正确事实的回答精确率为1，但$R_K$很低。例如某回答有$S=8$、$N=2$，则$P=0.8$；取$K=10$时$R_K=0.8$，$F_1@10=0.8$；取$K=5$时$R_K=1$，$F_1@5=2\times0.8/1.8\approx0.889$。$K$表达的是对回答长度的偏好，比较模型时必须固定。

SAFE的优点是领域开放、无需训练专门模型；原文在约一万六千条事实上与众包标注的一致率为72\%。代价同样明显：每条事实需要多次搜索与模型调用，搜索结果随时间变化，同一回答隔月重评可能得到不同结论，复现时应保存查询与结果快照。它的判定只有“支持”与“未获支持”两类，反驳与信息不足没有区分，与FActScore相同。

\subsection{回答修订与回答边界}
\label{sec:nlp-dialogue-answer-boundary}

核验得到逐条结论后，还要把它们变成可以交付的回答：保留被支持的部分，对信息不足的部分局部弃答，对被反驳的部分按来源修订，再重新核验新增的陈述。对本节的待查回答，整句弃答会丢掉已获支持的日期和地点；局部弃答可得“通知说明活动于4月18日在海棠楼201举行，但未给出到场人数”；若已取得S，则修订为“签到统计记录为120人”，并引用S而不是通知。

是否接受模型的判定，需要在独立开发样本上选阈值，并同时报告覆盖率与接受风险。设四条教学主张的支持置信为$0.95,0.85,0.75,0.40$，独立核验结果依次为对、错、对、错。阈值$0.8$接受两条，覆盖率$2/4$，接受样本的错误率$1/2$；阈值$0.9$只接受一条，覆盖率$1/4$，错误率为0。后一数字只描述这四个样本，阈值选择与最终测试须使用不同数据。

修订本身也可以自动化。RARR依靠外部检索，先查证再做最小编辑；CoVe不访问外部资料，让模型提出验证问题并独立回答，借自我一致性发现错误。两者都以“初稿已经存在”为前提，与第\ref{chap:nlp-text-generation}章先规划再生成的方法互补。

\subsubsection{检索–核验–最小修订}
\label{sec:nlp-dialogue-rarr}

语言模型的回答多数内容通常可用，只有少数片段有误。若为改正一个数字而重写整段，既浪费已有内容，也可能引入新错误。RARR（retrofit attribution using research and revision）的目标因此有两个：修订后的文本要能归因到检索到的证据，同时尽量保留原文\scite{nlp-dialogue-gao2023rarr}。它不需要训练，各模块均由少样本提示的语言模型实现，可以接在任何生成模型之后。

RARR分为研究阶段和修订阶段（算法\ref{alg:nlp-dialogue-rarr}）。研究阶段让模型针对原文$x$提出一组覆盖其各方面内容的核验问题，例如“ML-0418在哪天举行”“ML-0418现场有多少人”；对每个问题调用搜索引擎，把结果切成段落，按与问题的相关性排序，保留前几段作为证据。修订阶段逐一处理问题—证据对：一致性模型判断当前文本在该问题上是否与证据相符，不相符时编辑模型只改写相关片段，其余部分原样保留。最后从全部证据中选出至多$M$段组成\term{归因报告}（attribution report），选择准则是使每个核验问题都有尽量相关的证据覆盖。

\begin{algorithm}[RARR检索–核验–最小修订]
  \label{alg:nlp-dialogue-rarr}
  \begin{algorithmic}[1]
    \Require 原文$x$，语言模型$\mathrm{LM}$，搜索接口$\mathrm{Search}$，归因报告上限$M$
    \Ensure 修订文本$y$，归因报告$\mathcal{A}$
    \State $Q\gets\mathrm{LM}_{\mathrm{question}}(x)$\Comment{研究阶段：生成核验问题}
    \For{$q\in Q$}
      \State $E_q\gets$按与$q$的相关性排序$\mathrm{Search}(q)$中的段落，保留前若干段
    \EndFor
    \State $y\gets x$
    \For{$q\in Q$，$e\in E_q$}\Comment{修订阶段}
      \If{$\mathrm{LM}_{\mathrm{agree}}(y,q,e)$判为不一致}
        \State $y\gets\mathrm{LM}_{\mathrm{edit}}(y,q,e)$\Comment{只改与$q$相关的片段}
      \EndIf
    \EndFor
    \State $\mathcal{A}\gets$从$\bigcup_qE_q$中选至多$M$段，使各$q$的最大相关分之和最大
    \State \Return $y,\mathcal{A}$
  \end{algorithmic}
\end{algorithm}

评价同样分两项。归因度检查修订文本$y$中的陈述能否由报告$\mathcal{A}$支持，可由人工判断，也可用自然语言推断模型自动估计。\term{保持度}（preservation）检查改动是否克制：人工判断原意是否保留，记为$\operatorname{Pres}_{\mathrm{intent}}\in\{0,1\}$；字符层面用编辑距离计算
\begin{equation}
 \operatorname{Pres}_{\mathrm{Lev}}(x,y)=\max\Bigl(1-\frac{\operatorname{Lev}(x,y)}{|x|},\ 0\Bigr),\qquad
 \operatorname{Pres}=\operatorname{Pres}_{\mathrm{intent}}\cdot\operatorname{Pres}_{\mathrm{Lev}},
 \label{eq:nlp-dialogue-rarr-preservation}
\end{equation}
$\operatorname{Lev}$是字符级编辑距离，$|x|$是原文长度。两项指标此消彼长：整段重写可以让归因度很高而保持度很低，什么都不改则相反，所以原文用二者的调和平均作为综合分。

对待查回答，核验问题“现场有多少人”检索到S，一致性模型判为不一致，编辑模型把“150”改为“120”，其余不动。原文共25个字符，编辑距离为1，$\operatorname{Pres}_{\mathrm{Lev}}=1-1/25=0.96$；若改写为“签到记录显示ML-0418当日到场120人，活动在海棠楼201举行”，日期信息丢失，编辑距离也大得多。RARR的局限在于它只会编辑、不会删除：如果没有S，搜索找不到人数证据，“150人”会原样保留而缺少归因，这时需要本小节开头的局部弃答规则补充。

\subsubsection{验证问题链}
\label{sec:nlp-dialogue-cove}

外部检索并非总是可用。\term{验证问题链}（chain-of-verification, CoVe）观察到，模型在长回答中编造的细节，单独作为短问题提问时往往答得更准；利用这种不一致，不访问外部资料也能发现部分错误\scite{nlp-dialogue-dhuliawala2024cove}。其流程分四步：
\begin{equation}
 \begin{aligned}
 r_0&=\mathrm{LM}(x), &
 \{v_1,\ldots,v_m\}&=\mathrm{LM}_{\mathrm{plan}}(x,r_0),\\
 a_j&=\mathrm{LM}(v_j),\quad j=1,\ldots,m, &
 r&=\mathrm{LM}_{\mathrm{final}}\bigl(x,r_0,\{(v_j,a_j)\}_{j=1}^m\bigr).
 \end{aligned}
 \label{eq:nlp-dialogue-cove}
\end{equation}
先起草初稿$r_0$；再根据问题和初稿规划一组验证问题$v_j$，每个问题针对初稿中的一处事实；然后独立回答每个问题；最后把初稿与问答对一起交给模型，删去或改正与验证答案矛盾的内容，生成终稿$r$。

关键在第三步“独立”二字。若回答验证问题时初稿仍在上下文中，模型倾向于复述初稿，错误被原样确认。原文比较了几种执行方式：联合式在同一次生成中规划并回答问题，答案受初稿影响最大；两步式把规划与回答分开，回答时不看初稿；分解式进一步让每个问题单独调用一次模型，问题之间也互不干扰；分解加修订式再增加一步，显式比较每个答案与初稿是否一致。从联合式到分解式，答案受初稿的影响逐步减小，调用次数随之增加。

设初稿为“读书会4月18日在海棠楼201举行，主讲人为林宁”。规划出验证问题“ML-0418的主讲人是谁”，模型在不看初稿时若回答“材料未给出主讲人”，终稿删去主讲人一句。CoVe的边界也由此可见：验证答案仍出自同一模型的参数知识，模型若在短问题上同样确信错误事实，验证无法发现；它也不产生可供引用的外部证据。因此CoVe适合在无检索条件下降低编造，需要可归因结论时应与RARR类检索方法结合。

表\ref{tab:nlp-dialogue-verification-methods}汇总本节方法。从左到右的变化是：判定单位从整条主张变为原子事实，证据来源从固定文集扩展到网络搜索乃至模型自身，输出从判定标签扩展到修订后的回答。

\begin{table*}[htbp]
\centering
\small
\begin{tabular}{p{18mm}p{30mm}p{28mm}p{32mm}p{36mm}}
\toprule
方法 & 解决的问题 & 证据来源 & 核心机制 & 新增代价 \\
\midrule
FEVER流水线 & 大文集中的证据定位与三分类 & 固定文集 & 文档检索、句子打分选择、句对推断 & 错误逐级传递；需训练各阶段 \\
FActScore & 长回答的整体判断过粗 & 固定知识源 & 原子事实分解、逐条判支持 & 依赖分解质量；鼓励少说 \\
SAFE & 开放领域无现成条目 & 网络搜索 & 自足化改写、多步查询、逐事实判定、$F_1@K$ & 调用多；结果随时间变化 \\
RARR & 修订时整段重写引入新错 & 网络搜索 & 核验问题、一致性判断、最小编辑、归因报告 & 只编辑不删除；多轮调用 \\
CoVe & 无检索时发现编造 & 模型自身 & 规划验证问题、独立回答、据答案修订 & 无外部证据；共享模型盲点 \\
\bottomrule
\end{tabular}
\caption{主张与回答核验方法对照。FEVER流水线输出主张级标签，FActScore与SAFE输出逐事实统计，RARR与CoVe输出修订后的回答；三分类只在FEVER流水线中完整保留。}
\label{tab:nlp-dialogue-verification-methods}
\end{table*}

\section{任务型对话}
\label{sec:nlp-dialogue-task}

用户若只问活动何时举行，查通知即可；若要核对自己一行人的报名，答案还取决于身份、查询条件和业务记录，而且这些条件往往分几轮才说清，中途还可能更正。\term{任务型对话}（task-oriented dialogue）系统在多轮交流中帮助用户完成预订、查询一类业务，它面对四个挑战：理解每轮话语的意图与槽位；跨轮累积用户需求，并正确处理更正和撤回；根据当前需求和服务条件选择下一步动作，正确调用真实服务；把服务返回的结果如实转成自然语言，失败时给出可行的下一步。

经典做法是把这四项分给四个模块组成流水线：轮次理解、对话状态追踪、对话策略和回复生成。各模块随后分别神经化，又被合并为单个序列模型端到端生成；大语言模型出现后，状态追踪可以只靠上下文中的示例完成，整个流水线也可以由一个模型配合函数调用实现。本节按流水线的模块顺序展开，每个模块内部从规则到学习方法排列，最后一小节讨论大语言模型如何改变这一分工。

全节使用同一个只读业务。咨询时点为2025年4月17日，用户已登录，身份令牌为U17，目标是核对已有报名而非新增报名。只读接口\code{check\_registration}接受活动编号、认证身份、活动日期和同行人数。假设其快照中有一条U17的记录：活动ML-0418，日期4月18日，人数2，提交于4月16日17:30，状态有效。三轮用户输入为“帮我核对读书会的报名记录”“4月17日，两人”“说错了，是4月18日，人数不变”：第一轮缺日期与人数，第二轮条件下查无匹配，第三轮更正日期后查到记录（图\ref{fig:nlp-dialogue-state}）。

\begin{figure*}[htbp]
  \centering
  % Book-styled flowchart; solid arrows are data, dashed arrows update state.
\begin{tikzpicture}[
  every node/.style={font=\figurefont\small,align=center},
  utterance/.style={book node,text width=38mm,minimum height=16mm,
    fill=FigureInput!10,draw=FigureInput},
  state/.style={book node,text width=42mm,minimum height=16mm,
    fill=FigureModel!10,draw=FigureModel},
  result/.style={book node,text width=48mm,minimum height=16mm,
    fill=FigureOutput!10,draw=FigureOutput},
  flow/.style={book arrow,rounded corners=1.2mm},
  update/.style={flow,dashed,draw=FigureModel},
  labeltext/.style={text=BookMuted,font=\figurefont\small}
]
  \node[labeltext] at (0,14mm) {用户输入};
  \node[labeltext] at (48mm,14mm) {当前需求状态};
  \node[labeltext] at (104mm,14mm) {行为与观察};
  \node[utterance] (u1) at (0,0) {第1轮：核对报名\\日期、人数未给出};
  \node[state] (s1) at (48mm,0) {日期：缺失\\人数：缺失};
  \node[result] (a1) at (104mm,0) {缺必需字段\\询问日期与人数};
  \node[utterance] (u2) at (0,-28mm) {第2轮：新增约束\\4月17日、两人};
  \node[state] (s2) at (48mm,-28mm) {日期：4月17日\\人数：2};
  \node[result] (a2) at (104mm,-28mm) {只读查询：查无匹配\\核对是否指4月18日};
  \node[utterance] (u3) at (0,-56mm) {第3轮：更正日期\\4月18日、人数不变};
  \node[state] (s3) at (48mm,-56mm) {日期：4月18日\\人数：2（保持）};
  \node[result] (a3) at (104mm,-56mm) {只读查询：匹配1条\\核验两人已按时提交};
  \foreach \i in {1,2,3} {
    \draw[update] (u\i.east) -- (s\i.west);
    \draw[flow] (s\i.east) -- (a\i.west);
  }
  \draw[flow] (a1.south) -- (104mm,-14mm)
    -- (0,-14mm) -- (u2.north);
  \draw[flow] (a2.south) -- (104mm,-42mm)
    -- (0,-42mm) -- (u3.north);
  \node[labeltext,fill=white,inner sep=1mm] at (48mm,-14mm) {回复：补全必需条件};
  \node[labeltext,fill=white,inner sep=1mm] at (48mm,-42mm) {回问：让用户确认修改};
  \node[book node,text width=110mm,minimum height=12mm,
    fill=BookAmber!10,draw=BookAmber] (failure) at (53mm,-84mm)
    {独立故障分支：TIMEOUT不是查无记录\\说明失败，按预算重试或回问是否稍后继续};
  \draw[flow,draw=BookAmber] (a2.east) -- (138mm,-28mm)
    -- (138mm,-84mm) -- (failure.east);
  \node[labeltext] at (53mm,-98mm)
    {固定上下文：认证U17、活动ML-0418；不新增或修改报名};
\end{tikzpicture}

  \caption{只读报名核验的三轮对话。用户目标为核对U17在ML-0418的已有报名；实线传递话语、动作和接口结果，虚线表示状态写入或更正。第二轮查无匹配后回问，第三轮只覆盖日期而保持人数。接口故障另走重试或说明失败的分支，不产生报名结论。}
  \label{fig:nlp-dialogue-state}
\end{figure*}
\FloatBarrier

\subsection{轮次理解}
\label{sec:nlp-dialogue-turn}

轮次理解把单轮话语解析为结构化的语义：\term{意图}（intent）是用户这一轮想做的事，如“核对报名”；\term{槽位}（slot）是完成该意图所需的参数，如日期和人数；操作说明本轮对槽位做了什么，如新增、更正或保持。日期解析以给定的咨询时点为参照，省略和指代较多的话语可先按第\ref{sec:nlp-analysis-utterance-rewriting}小节改写成独立句。方法谱系是：意图分类与槽位标注分别训练；JointBERT用共享编码器联合训练两者；模式约束的记录生成进一步输出操作类型和来源轮次。

\subsubsection{意图分类与CRF槽位标注}
\label{sec:nlp-dialogue-intent-crf}

流水线把轮次理解拆为两个已知任务：句级意图用分类器预测，以交叉熵训练；槽位转为BIO标签序列，用线性链CRF标注，以正确路径的负对数似然训练。BIO编码、CRF的路径打分与合法转移掩码见第\ref{sec:nlp-analysis-paradigm-tagging}小节，这里只说明对话中增加的部分。

第二轮按“4月／17日／，／两／人”分词，正确标签为\code{B-DATE I-DATE O B-PARTY I-PARTY}。解码后按分词到字符的映射把相邻标签合并回原文跨度“4月17日”和“两人”，再规范化为2025-04-17和整数2。设意图分类器在“核对报名、取消报名、查询通知”上输出$(0.8,0.1,0.1)$，选择核对报名。意图只决定走哪个业务流程，不代表权限：本节服务在业务层禁用取消动作，意图误判为取消时也不会执行。

这种组合的缺陷在第三轮暴露出来。“说错了，是4月18日”中，普通DATE标签只能标出新日期的位置，无法表示它要替换哪个旧值；“人数不变”没有任何槽位跨度，标注器的输出为空。更正、撤回与保持需要另行定义标签或规则。

\subsubsection{JointBERT联合意图与槽位}
\label{sec:nlp-dialogue-jointbert}

意图和槽位彼此相关：知道意图是“核对报名”，日期更可能是活动日期；反过来，出现日期和人数也提示意图。分别训练两个模型时这种相关无法共享。JointBERT让两项任务共用一个BERT编码器\cite{nlp-dialogue-chen2019}：句首\code{[CLS]}的表示$\bm{h}_0\in\mathbb{R}^{d}$进入意图分类器，每个原词首个子词的表示进入槽位分类器，两项交叉熵相加：
\begin{equation}
\begin{aligned}
 p(c\mid u)&=\operatorname{softmax}(\bm{W}_{c}\bm{h}_0+\bm{b}_{c}),\\
 p(s_i\mid u)&=\operatorname{softmax}(\bm{W}_{s}\bm{h}_{a(i)}+\bm{b}_{s}),\\
 L&=-\log p(c^*\mid u)-\sum_i\log p(s_i^*\mid u).
\end{aligned}
\label{eq:nlp-dialogue-jointbert}
\end{equation}
其中$a(i)$是原词$i$的首子词位置，$\bm{W}_c\in\mathbb{R}^{|\mathcal{C}|\times d}$、$\bm{W}_s\in\mathbb{R}^{|\mathcal{T}|\times d}$分别映射到意图类别和BIO标签。两个头的梯度都回传到共享编码器，这是与分别训练的直接差别。

对第二轮，设正确意图概率为0.8，五个正确标签的概率均为0.9，则$L=-\log0.8-5\log0.9\approx0.750$。基础模型的槽位是逐位置独立分类，可能输出句首\code{I-DATE}这类非法序列；原文也比较了在槽位头上加CRF的版本，以引入标签转移约束。JointBERT与上一小节一样，输出空间仍是单意图加BIO标签，第三轮的“人数不变”需要新增\code{KEEP(party)}一类操作标签才能表达。

\subsubsection{模式约束的语义记录生成}
\label{sec:nlp-dialogue-semantic-record}

要表达更正、撤回和来源，最直接的办法是让模型生成结构化记录。给定当前轮、限定范围的历史和字段模式，模型输出\code{intent}、\code{operation}、\code{slot}、\code{old\_value}、\code{new\_value}和\code{source\_turn}等字段；有标注时以记录词元的负对数似然训练，提示式实现则固定示例和模型版本。解码时用第\ref{sec:nlp-inference-grammar}小节的文法约束屏蔽非法键和枚举值，生成后再用解析器检查类型。

第三轮应生成两条操作：日期\code{REPLACE(2025-04-17, 2025-04-18)}，人数\code{KEEP}；新日期的依据在第3轮，旧日期的依据在第2轮。若上一轮同时讨论了活动日期和报名截止，“改成周五”未指明对象，就应输出待澄清，而不是同时覆盖两个字段。开放值生成省去了为每个取值预设类别，但会出现格式合法、含义错误的记录：原话若是“不是4月18日”，\code{REPLACE(date, 2025-04-18)}键值都合法，含义却正好相反。因此解析后仍要核对来源跨度、否定范围和历史绑定。

\subsection{对话状态追踪}
\label{sec:nlp-dialogue-state}

轮次理解只看一轮，业务调用却需要截至当前的全部有效要求。\term{对话状态追踪}（dialogue state tracking, DST）在每轮之后维护\term{对话状态}（dialogue state）$B_t$，即当前有效的领域—槽位—值集合，而不是历史上出现过的所有词。MultiWOZ提供多领域对话及逐轮状态标注\cite{budzianowski2018multiwoz}，SGD进一步用服务与槽位的自然语言描述定义接口，并区分训练中见过和未见过的服务\scite{rastogi2020sgd}。常用指标是\term{联合目标准确率}（joint goal accuracy），即预测状态与标注状态在所有槽位上完全一致的轮次比例；它要求每轮全部槽位正确，一处错误会在后续各轮持续计错。

槽位值需要区分四种特殊情况：尚未提及为缺失，用户明确说不知道为未知，表示任意值都可接受为不限定，明确拒绝某个值为否定约束。状态追踪有三个挑战：用户更正时要覆盖旧值而非并存；槽位取值开放，无法预先列出所有日期和名称；新服务不断出现，标注数据跟不上。方法谱系为：规则增量更新，可审计但依赖理解结果；TRADE用复制机制生成开放值；SimpleTOD把状态作为整个对话序列的一部分生成；IC-DST用检索到的示例让大语言模型直接输出状态变化，几乎不需要训练数据。

\subsubsection{规则槽位写入与撤回}
\label{sec:nlp-dialogue-state-rule}

规则更新器接收旧状态$B_{t-1}$和本轮解析出的操作集合$\Delta_t$，计算$B_t=U(B_{t-1},\Delta_t)$，并为每个值记录来源轮号。规则来自字段模式与业务含义，不需要训练。表\ref{tab:nlp-dialogue-state-update}列出本节使用的转换；遇到冲突时不按“最后出现者覆盖”处理，而是标记待澄清。

\begin{table}[htbp]
\centering
\begin{tabular}{p{22mm}p{74mm}}
\toprule
操作 & 状态转换\\
\midrule
新增 & 旧槽位缺失时写入值与来源；已有不相容值则待澄清\\
更正 & 明确指向旧槽位时覆盖，旧值移入历史而非有效值\\
保持 & 不改变原值及其来源\\
撤回 & 删除有效值；若只拒绝某值，则写入否定集合\\
未知／不限定 & 分别写入不同标记，不用同一空字符串表示\\
\bottomrule
\end{tabular}
\caption{对话状态的规则更新。缺失是从未提及，撤回是有来源的交流事件，两者在当前有效值表中可能相同，在历史记录中不同。}
\label{tab:nlp-dialogue-state-update}
\end{table}

活动编号和身份由已认证的上下文提供，以下只演算用户可修改的两个槽位：$B_0=B_1=\varnothing$；$B_2=\{(\text{日期},4/17),(\text{人数},2)\}$，两值均来自第2轮；$B_3=\{(\text{日期},4/18),(\text{人数},2)\}$，日期来自第3轮，人数仍来自第2轮。第三轮的变化集只有一个覆盖操作，不能因为本轮没有提到人数就把2清空。增量规则每次改动都可审计，但前面误写的值会一直保留下去，其可靠性完全取决于输入操作是否正确。

\subsubsection{TRADE复制式状态生成}
\label{sec:nlp-dialogue-trade}

规则更新依赖轮次理解的输出，而早期神经状态追踪器又多把每个槽位当作在预定义取值表上的分类，遇到表外的日期或名称就无能为力。TRADE（transferable dialogue state generator）直接读取对话历史，为每个领域—槽位对生成值，生成时可以从历史中复制词元，因此不需要预定义取值表\scite{nlp-dialogue-wu2019}。

模型结构如下。对话历史的词元经嵌入和双向GRU编码为$\bm{H}=[\bm{h}_1,\ldots,\bm{h}_n]\in\mathbb{R}^{n\times d}$。对第$j$个领域—槽位对，以领域嵌入与槽位嵌入之和作为解码器的首个输入，所有槽位共享同一个GRU解码器，各自独立解码。第$k$步解码状态为$\bm{g}_{j,k}$，对历史的注意力为$\bm{\alpha}_{j,k}=\operatorname{softmax}(\bm{H}\bm{g}_{j,k})\in\mathbb{R}^{n}$。输出分布按第\ref{sec:nlp-generation-copy-pointer}小节的复制门混合词表分布与历史复制分布，同一词在历史中多次出现时复制概率相加：
\begin{equation}
 p(y\mid j,k)=
 \gamma_{j,k}\,p_{\mathrm{vocab}}(y\mid j,k)
 + (1-\gamma_{j,k})\sum_{i:x_i=y}\alpha_{j,k,i}.
 \label{eq:nlp-dialogue-trade}
\end{equation}
$\gamma_{j,k}\in(0,1)$由解码状态、当前输入嵌入和注意力摘要$\bm{c}_{j,k}=\bm{H}^\top\bm{\alpha}_{j,k}$经sigmoid计算。另有一个\term{槽门}（slot gate）用首步的$\bm{c}_{j,0}$做三分类：\code{none}表示该槽位无值，\code{dontcare}表示用户不限定，\code{ptr}表示采用解码出的值。训练目标为$L=\alpha L_{\mathrm{gate}}+\beta L_{\mathrm{value}}$，前者是槽门交叉熵，后者是正确值词元的负对数似然，监督来自每轮的完整状态。

第三轮解码日期时，设$\gamma=0.2$，词表对“4月18日”的概率为0.1，历史中对应位置的复制概率为0.7，则混合概率为$0.2\times0.1+0.8\times0.7=0.58$；旧日期“4月17日”若词表概率0.2、复制概率0.1，混合后为0.12，解码选择新日期。人数槽位仍可从第2轮复制“两人”，规范化后保留为2。

TRADE的优点是开放取值和跨领域共享：所有槽位共用解码器，新领域中与旧领域相似的槽位（如不同领域的“日期”）可以零样本迁移。它的弱点同样来自复制：历史中失效的旧值随时可被复制出来，生成完全合法却过期的状态；槽门三类也不足以表达未知、撤回和否定约束，“用户从未提及”与“用户明确撤销”都会映射为\code{none}。此外，每轮都从全部历史重新生成全部槽位，计算量随槽位数和历史长度增长。

\subsubsection{SimpleTOD完整状态序列生成}
\label{sec:nlp-dialogue-simpletod}

TRADE只负责状态，策略和回复仍需其他模块。SimpleTOD把整个流水线重写为一个序列：对话历史$C$、对话状态$B$、数据库结果摘要$D$、系统行为$A$和去词汇化回复$S$依次用边界标记连接，由一个预训练因果语言模型按下一词元负对数似然训练\scite{nlp-dialogue-hosseini2020}。这样做的动机是让各模块共享同一预训练模型的参数和语言知识，把模块间的接口变成序列中的位置。

推断时分两段生成。先由$C$生成$B$，在状态结束标记处停下，解析、校验后实际查询数据库，得到$D$（匹配条目数等摘要）；再以$[C;B;D]$为条件生成$A$和$S$。$D$必须来自真实查询，而不是让模型续写；具体返回行保留在程序侧，供最后回填。原文把数据库段作为可选输入，其端到端最佳设置甚至不读数据库结果；本节保留它，是为了让回复以真实返回为依据。第三轮的教学序列可读作“历史：三轮；状态：日期4/18、人数2；数据库：匹配1条；行为：告知已有有效记录；回复：查到您在[date]的[party]人报名记录”，其中“匹配1条”是调用后的观察，不是人数。

与TRADE相比，SimpleTOD共享更多参数，结构也最简单；代价是状态、动作和回复的错误相互影响，状态生成仍需训练数据，且序列随对话增长而变长。

\subsubsection{上下文学习式状态追踪}
\label{sec:nlp-dialogue-icdst}

以上方法都需要大量逐轮标注的状态，每增加一个服务就要补充标注。大语言模型可以通过第\ref{sec:nlp-inference-icl-shots}小节的上下文学习，只看少量示例完成新任务；困难在于状态是结构化对象，示例又必须与当前轮相似才有效。IC-DST（in-context learning for dialogue state tracking）针对这两点设计了输出格式和示例检索\scite{nlp-dialogue-hu2022icdst}。

第一个设计是把状态写成SQL。服务模式写成\code{CREATE TABLE}语句，每个领域一张表，槽位为列；本轮的状态变化写成一条\code{SELECT}查询，槽位约束写在\code{WHERE}子句中。代码模型在预训练中见过大量SQL，这种格式比自定义的槽值列表更容易被正确生成，也可以用SQL解析器检查合法性。第二个设计是只预测本轮变化：输入是上一轮状态$B_{t-1}$（以文本列出）和最近一对系统与用户话语$(A_t,U_t)$，输出本轮的状态变化$\Delta_t$，完整状态由
\begin{equation}
 B_t=B_{t-1}\oplus\Delta_t
 \label{eq:nlp-dialogue-icdst}
\end{equation}
得到，$\oplus$表示用$\Delta_t$中的槽值覆盖或新增$B_{t-1}$中的对应槽位。这样输入长度不随对话轮数增长，示例也只需包含一轮。代码\ref{code:nlp-dialogue-icdst}是按这一格式写的教学提示。

\begin{codeblock}[language=SQL,label={code:nlp-dialogue-icdst}]{IC-DST格式的状态追踪提示（教学示例）}
CREATE TABLE registration(
  event_id text,   -- 活动编号
  date text,       -- 活动日期
  party_size int   -- 同行人数
);
-- 检索到的示例
-- [context] registration-date: 2025-04-20, registration-party_size: 3
-- [system] 按4月20日、三人未查到记录。
-- Q: [user] 记错了，是4月21日
SELECT * FROM registration WHERE date = '2025-04-21';
-- 当前轮
-- [context] registration-date: 2025-04-17, registration-party_size: 2
-- [system] 按4月17日、两人未查到匹配记录，您要核对4月18日吗？
-- Q: [user] 说错了，是4月18日，人数不变
SELECT * FROM registration WHERE
\end{codeblock}

模型续写出\code{date = '2025-04-18';}，解析为$\Delta_3=\{(\text{日期},4/18)\}$，代入式\eqref{eq:nlp-dialogue-icdst}得到与规则更新相同的$B_3$。

示例选择决定了效果。IC-DST从有标注的样本池中按相似度检索最近的$k$个轮次作为示例，相似度的定义是关键：两个轮次的话语可能用词相近，状态变化却完全不同。因此它以两轮状态变化的重合程度作为相似度标签：分别按槽位集合与槽位—值集合计算两轮变化之间的F1，取两者平均$s=\frac12(F_{\mathrm{slot}}+F_{\mathrm{slot\text{-}value}})$；再从每个样本的近邻中取$s$最高与最低的一部分构成正、负样本对，微调一个句向量检索器，使状态变化相似的轮次在向量空间中接近。例如当前轮的变化只有“日期改为4/18”，候选甲只把日期改为4/21，候选乙同时改日期和人数。槽位F1甲为1，乙为$2\times1\times\frac12/(1+\frac12)\approx0.667$；槽位—值F1两者都为0（值不同）；于是$s$甲为0.5、乙约为0.333，检索器应把甲排在前面。

IC-DST的优点是几乎不需要训练数据（原文以Codex为语言模型，在只用MultiWOZ百分之一到百分之十的标注时超过了此前的少样本方法；换用较小的开源代码模型时则不及这些基线），新增服务只需改写表结构和补充少量示例。缺点是增量更新会累积错误：$\Delta_t$一旦漏掉更正，错误会通过$B_{t-1}$传到所有后续轮次；每轮都要调用大模型，成本和延迟高于小模型；撤回、否定约束同样需要在SQL格式中另作约定。

表\ref{tab:nlp-dialogue-dst-methods}汇总四种状态追踪方法。

\begin{table*}[htbp]
\centering
\small
\begin{tabular}{p{18mm}p{30mm}p{36mm}p{34mm}p{26mm}}
\toprule
方法 & 解决的问题 & 核心机制 & 新增代价 & 监督需求 \\
\midrule
规则更新 & 更正与保持的可审计处理 & 按操作类型增量改写状态 & 依赖轮次理解；旧错持续 & 无需训练 \\
TRADE & 预定义取值表无法覆盖开放值 & 共享解码器、复制门、三类槽门 & 可能复制过期值 & 逐轮完整状态 \\
SimpleTOD & 模块分离、参数不共享 & 状态、查询摘要、行为、回复串成一个序列 & 错误相互影响；序列变长 & 逐轮状态、行为与回复 \\
IC-DST & 新服务缺少标注 & 状态写成SQL、检索相似轮次、只预测变化 & 增量误差累积；大模型调用成本 & 少量示例 \\
\bottomrule
\end{tabular}
\caption{对话状态追踪方法对照。前三种方法的状态由训练好的专用模型或规则产生，IC-DST用冻结的大语言模型加检索示例产生。}
\label{tab:nlp-dialogue-dst-methods}
\end{table*}

\subsection{对话策略与业务调用}
\label{sec:nlp-dialogue-action}

状态完整之后，系统还要决定下一步做什么。\term{对话策略}（dialogue policy）根据当前状态和此前的服务返回，从询问、澄清、查询、告知、结束等动作中选择一个。服务模式规定了字段类型、必需条件和可用动作。规则策略可控，但难以覆盖复杂情形；监督学习从对话语料中模仿动作；强化学习以任务成功为奖励优化长程决策。三种方法共有一个前提：动作被选中不等于已经执行，业务结论只能来自真实返回。

\subsubsection{有限状态与槽位驱动策略}
\label{sec:nlp-dialogue-rule-policy}

规则策略由服务设计者写定：缺必需字段就询问；字段冲突或指代不明就澄清；参数齐备且认证有效才允许查询；取得与目标相符的结果后告知并结束。允许动作掩码把不可用动作的分数设为$-\infty$；由于已过报名截止且服务只读，本例任何轮次都不允许新增报名。

第一轮日期与人数均缺失，动作为“请问您核对哪天、几人的报名记录”。第二轮以\code{(ML-0418, U17, 2025-04-17, 2)}查询，返回\code{OK}且记录为空，策略转为条件核对：“按4月17日、两人未查到匹配记录；通知中的活动日期为4月18日，您要核对这个日期吗？”系统不应自行把日期改成通知日期再报告结果。第三轮以更正后的参数查询，返回一条有效记录，规则另验提交时点$17{:}30<18{:}00$，告知并结束。若返回\code{TIMEOUT}，则没有任何关于记录的结论，可选动作是说明失败、按预算重试或请用户稍后再试。

\subsubsection{监督动作预测与结构化行动生成}
\label{sec:nlp-dialogue-learned-policy}

规则数量增多后，可以从带动作标注的对话中学习策略。分类路线把历史摘要、当前状态和此前服务返回编码为$\bm{v}_t$，以$\operatorname{softmax}(\bm{W}\bm{v}_t)$预测离散动作，最小化交叉熵，参数再由类型化的槽位回填；生成路线以同样的输入直接生成动作名与参数组成的结构记录，以记录词元的负对数似然训练，并约束动作名、字段和枚举值。两条路线都与允许动作掩码配合：第一轮若未加掩码时$\Pr(\text{询问})=0.7$、$\Pr(\text{查询})=0.3$，缺必需字段使查询不可用，掩码后只能询问。生成路线若写出人数\code{"两"}，先由规范化器转为整数2或拒绝，不能依赖服务容忍任意格式。

监督学习只能模仿语料中的动作分布。它不知道哪个动作最终带来任务成功，语料中没出现过的情形（如本例第二轮的空结果）也学不到恰当反应；对话越长，早期动作的后果越难从逐轮标签中体现出来。

\subsubsection{强化学习对话策略}
\label{sec:nlp-dialogue-rl-policy}

对话策略的好坏最终体现在整段对话是否完成任务、用了多少轮，这正是序列决策问题。把对话建模为\term{部分可观测马尔可夫决策过程}（partially observable Markov decision process, POMDP）的思路很早就已提出\cite{nlp-dialogue-levin2000,nlp-dialogue-young2013}：用户的真实目标不可直接观测，系统只能看到对话状态和服务返回。

形式化如下。时刻$t$的状态$s_t=(B_t,D_t)$由对话状态和最近一次服务返回组成；动作$a_t$取自对话行为集合，如\code{request(date)}、\code{confirm(date)}、\code{query}、\code{inform}、\code{bye}；用户回应和服务返回决定下一状态。奖励按任务设计，本节采用教学取值：每轮$-1$，鼓励简短；成功完成目标时$+20$；失败结束时$-10$。策略$\pi_{\bm\theta}(a\mid s)$的目标是最大化期望回报
\begin{equation}
 J(\bm\theta)=\mathbb{E}_{\tau\sim\pi_{\bm\theta}}\Bigl[\sum_{t=1}^{T}\gamma^{t-1}r_t\Bigr],
 \label{eq:nlp-dialogue-rl-objective}
\end{equation}
$\tau$是一整段对话轨迹，$\gamma\in(0,1]$为折扣因子。

这里只用到策略梯度的一个结论：目标对策略参数的梯度可写为$\nabla_{\bm\theta}J=\mathbb{E}\bigl[\sum_t\nabla_{\bm\theta}\log\pi_{\bm\theta}(a_t\mid s_t)\,(G_t-\beta(s_t))\bigr]$，其中$G_t$是从$t$起的回报，基线$\beta(s_t)$不依赖$a_t$，减去它不改变梯度期望、只降低方差\scite{rl1-williams1992-reinforce}。也可以改用深度Q学习估计动作价值，早期端到端任务型对话系统即采用此法\scite{nlp-dialogue-li2017e2e}。策略梯度、价值方法及其成立条件可参考\BookChapterRef{chap:rl-learning}{卷七“未知模型的求解”章}。

用一个两动作的教学例子看梯度方向。第二轮查询为空后，策略在“澄清是否指4月18日”与“告知无记录并结束”之间选择，当前概率各为0.5，取$\gamma=1$。澄清的轨迹从第二轮起依次得到$-1$与$-1+20$，$G_2=18$；结束的轨迹得到$-1-10$，$G_2=-11$。取两者均值$3.5$为基线，优势分别为$+14.5$与$-14.5$。对“澄清”的logit求导，$\nabla\log\pi(\text{澄清})=1-0.5=0.5$，$\nabla\log\pi(\text{结束})=-0.5$，两条轨迹的贡献都是$+7.25$：成功的轨迹抬高澄清，失败的轨迹压低结束，二者方向一致。

强化学习需要大量交互，与真人交互成本太高，因此通常先在\term{用户模拟器}（user simulator）上训练。议程式模拟器为每段对话抽样一个用户目标，包括约束（日期4月18日、人数2）和要询问的信息，再维护一个待说动作的栈：系统询问日期时把“提供日期”压栈，每轮从栈顶弹出若干动作作为用户回应\scite{nlp-dialogue-schatzmann2007}。模拟器还可以加入误差模型，按一定概率替换槽值，模拟理解模块的识别错误，使策略学会确认。

这一方法的问题来自模拟器与奖励。模拟器的行为比真人规整，策略会利用它的规律，换到真人时性能下降；奖励设计不当会被钻空子，例如成功判定只看是否给出了“告知”动作而不核对内容，策略就会学到尽早告知。强化学习策略的优势是能针对任务成功优化长程决策、处理语料中少见的情形，代价是需要可靠的模拟器、可检查的成功判定和大量训练交互。

\subsection{回复生成与交流恢复}
\label{sec:nlp-dialogue-response}

最后一个模块把系统行为和真实返回转为自然语言。这里要区分两种状态：对话状态描述用户当前要什么，业务状态描述服务中实际有什么。用户更正日期只改变前者，不会改动报名记录；回复必须在两者之间建立有依据的对应。模板保证字段来自真实返回，但表达单一；SC-LSTM学习按语义条件生成，表达更自然，需要另外检查遗漏与多余。

\subsubsection{去词汇化模板与槽位回填}
\label{sec:nlp-dialogue-template}

模板把结果类型与语言表达分开。设计者为成功、空结果和故障三类返回各写模板，规定每个占位符取自接口返回还是查询条件。成功模板为“查到[date]的[party]人有效报名，提交于[submitted]”；空结果模板为“按[date]、[party]人未查到匹配记录，请核对条件”；故障模板不包含关于记录是否存在的任何断言。去词汇化指训练与生成时用占位符代替具体槽值，最后再回填。

第二轮的\code{OK, records=[]}选空结果模板，日期和人数取自本次请求；第三轮取得匹配行后，成功模板填入“4月18日”“2”“4月16日17:30”。回填前检查返回行的身份与活动是否与请求一致、所有占位符是否消除；若用户已改变目标而旧的异步结果才返回，应按请求版本丢弃。模板没有训练成本、不会编造字段，但表达固定，覆盖新情形要人工增写，也无法纠正上游的错误状态。

\subsubsection{SC-LSTM语义条件回复}
\label{sec:nlp-dialogue-sclstm}

模板数量随行为与槽位组合增长，表达也显得生硬。语义控制LSTM（semantically controlled LSTM, SC-LSTM）从带对话行为标注的语料中学习“哪种语义用哪种说法”，不需要逐词对齐\scite{nlp-dialogue-wen2015}。输入是一个控制向量$\bm{d}_0$，由行为类型和槽位的独热特征组成，表示待表达的全部语义。它在普通LSTM单元之外增加读取门$\bm{r}_t$，每生成一个词就从$\bm{d}$中“划掉”已表达的部分：
\begin{equation}
\begin{aligned}
 \bm{r}_t&=\sigma(\bm{W}_{wr}\bm{w}_t+
                   \alpha\bm{W}_{hr}\bm{h}_{t-1}),\\
 \bm{d}_t&=\bm{r}_t\odot\bm{d}_{t-1},\\
 \bm{c}_t&=\bm{f}_t\odot\bm{c}_{t-1}
       +\bm{\imath}_t\odot\widetilde{\bm{c}}_t
       +\tanh(\bm{W}_{dc}\bm{d}_t).
\end{aligned}
\label{eq:nlp-dialogue-sclstm}
\end{equation}
$\bm{w}_t$为当前输入词向量，$\bm{\imath}_t$、$\bm{f}_t$为输入门与遗忘门，其余与普通LSTM相同。剩余语义$\bm{d}_t$经$\tanh(\bm{W}_{dc}\bm{d}_t)$进入细胞状态，持续影响后续生成。训练目标为
\begin{equation}
 L=-\sum_{t}\log p(w_t^*\mid w_{<t}^*,\bm{d}_0)+\lVert\bm{d}_T\rVert+\sum_{t=0}^{T-1}\eta\,\xi^{\lVert\bm{d}_{t+1}-\bm{d}_t\rVert},
 \label{eq:nlp-dialogue-sclstm-loss}
\end{equation}
第一项是词元负对数似然；第二项惩罚结束时仍未表达的语义，防止漏说；第三项惩罚相邻步之间$\bm{d}$的剧烈变化，$\xi$取较大的底数，使一次划掉多个语义特征的代价急剧增大。生成时先为每个对话行为过生成20个候选，再按$R=-(F_{\mathrm f}+F_{\mathrm b}+\lambda\,\mathrm{ERR})$重排并保留前5个，其中$F_{\mathrm f},F_{\mathrm b}$是前向与反向生成器给出的代价，$\mathrm{ERR}=(p+q)/N$是槽位错误率，$p,q$为遗漏与多余的槽位数，$N$为应表达的槽位总数。

对第三轮结果，设待表达语义为日期、人数、提交时点，$\bm{d}_0=(1,1,1)$。生成日期附近时读取门为$(0.1,1,1)$，$\bm{d}$变为$(0.1,1,1)$；生成人数时门为$(1,0.1,1)$，$\bm{d}$变为$(0.1,0.1,1)$。若句子到此结束，$\lVert\bm{d}_T\rVert_2\approx1.01$，主要来自未表达的提交时点，训练会推动模型继续生成。读取门是软约束，不保证每个槽位恰好出现一次，因此重排时仍要检查遗漏与重复，并用真实记录回填。

\subsection{基于大语言模型与函数调用的任务型对话}
\label{sec:nlp-dialogue-llm-function}

前面四个模块各需一套标注和模型，新增一个服务就要更新理解标签、状态模式、动作集合和回复模板。经过指令微调的大语言模型已经能够读懂业务说明、追踪多轮内容并写出自然的回复，于是可以把整条流水线交给一个模型：系统提示给出业务规则和可调用函数的说明，模型在对话中决定何时调用哪个函数、填什么参数，执行结果回填到上下文后再生成回复。这里的\term{函数调用}（function calling）指模型输出一个结构化的调用请求，由外部执行器运行并返回结果，模型本身不直接访问服务。

第一步是把槽位和接口写成\term{函数模式}（function schema）。函数模式用JSON Schema描述函数名、用途和参数的类型与约束，必需参数对应原先的必需槽位，枚举与格式约束对应槽位取值范围（代码\ref{code:nlp-dialogue-function-schema}）。用户身份不出现在参数中，由执行器从已认证会话注入，避免模型填写或篡改他人身份。

\begin{codeblock}[label={code:nlp-dialogue-function-schema}]{只读报名核验的函数模式（教学示例）}
{
  "name": "check_registration",
  "description": "只读查询当前登录用户在指定活动的已有报名记录；不能新增或修改报名。",
  "parameters": {
    "type": "object",
    "properties": {
      "event_id":   {"type": "string", "enum": ["ML-0418"]},
      "date":       {"type": "string", "format": "date",
                     "description": "用户所说的活动日期，YYYY-MM-DD"},
      "party_size": {"type": "integer", "minimum": 1,
                     "description": "同行人数"}
    },
    "required": ["event_id", "date", "party_size"]
  }
}
\end{codeblock}

第二步是调用—结果—回复循环（算法\ref{alg:nlp-dialogue-function-loop}，图\ref{fig:nlp-dialogue-function-loop}）。每轮用户话语到来后，模型以系统提示、函数模式和完整对话历史为条件生成：若输出普通文本，就作为回复交给用户；若输出函数调用，执行器先按模式校验参数、检查权限，再调用服务，把返回结果作为一条工具消息追加到历史，模型继续生成，直到给出回复或达到调用次数上限。

\begin{algorithm}[函数调用式任务型对话的单轮处理]
  \label{alg:nlp-dialogue-function-loop}
  \begin{algorithmic}[1]
    \Require 系统提示$P$（业务规则），函数模式集合$\mathcal{F}$，对话历史$H$，用户话语$u$，已认证身份$z$，调用上限$K$
    \Ensure 回复$o$，更新后的历史$H$
    \State $H\gets H\oplus u$
    \For{$k=1,\ldots,K$}
      \State $m\gets\mathrm{LM}(P,\mathcal{F},H)$；$H\gets H\oplus m$
      \If{$m$为普通文本}
        \State \Return $o=m$, $H$
      \EndIf
      \State 解析$m$得函数名$f$与参数$\bm{a}$
      \If{$f\notin\mathcal{F}$或$\bm{a}$不符合$f$的模式或$f$需要未授予的权限}
        \State $\rho\gets$错误说明\Comment{不执行，交给模型改正或询问}
      \Else
        \State $\rho\gets\mathrm{Execute}(f,\bm{a},z)$\Comment{身份由执行器注入}
      \EndIf
      \State $H\gets H\oplus\rho$\Comment{工具消息，含成功、空结果、超时等状态}
    \EndFor
    \State \Return $o=$“暂时无法完成查询”的说明，$H$
  \end{algorithmic}
\end{algorithm}

\begin{figure*}[htbp]
  \centering
  % 第九卷 13.3.5：基于大语言模型与函数调用的任务型对话（跨栏，约160 mm宽）。
% FigureInput：用户；FigureModel：大语言模型；FigureOutput：执行器与服务；虚线框：可选的显式状态记录。
\begin{tikzpicture}[
  x=1mm,y=1mm,
  figure picture,
  user/.style={figure node,draw=FigureInput,fill=FigureInputFill,minimum height=14mm,
    minimum width=16mm,inner xsep=2pt},
  model/.style={figure node,draw=FigureModel,fill=FigureModelFill,minimum height=14mm,
    minimum width=40mm,inner xsep=2pt,line width=.9pt},
  exec/.style={figure node,draw=FigureOutput,fill=FigureOutputFill,minimum height=14mm,
    minimum width=30mm,inner xsep=2pt},
  record/.style={figure node,draw=FigureMuted,fill=white,minimum height=8mm,
    minimum width=56mm,inner xsep=2pt,dash pattern=on 2pt off 1.5pt}
]
  \node[user] (user) at (8,30) {用户};
  \node[model] (llm) at (56,30) {大语言模型\\系统提示＋函数模式＋历史};
  \node[exec] (exe) at (110,30) {执行器\\校验参数与权限\\注入认证身份};
  \node[exec] (svc) at (148,30) {服务\\\texttt{check\_registration}};
  \node[record] (state) at (56,6) {状态记录（可选）：日期、人数及来源轮次};

  \draw[figure flow] ([yshift=3mm]user.east) -- node[figure edge label,above=1pt] {话语} ([yshift=3mm]llm.west);
  \draw[figure flow] ([yshift=-3mm]llm.west) -- node[figure edge label,below=1pt] {回复} ([yshift=-3mm]user.east);
  \draw[figure flow] ([yshift=3mm]llm.east) -- node[figure edge label,above=1pt] {函数调用} ([yshift=3mm]exe.west);
  \draw[figure flow] ([yshift=-3mm]exe.west) -- node[figure edge label,below=1pt] {工具消息} ([yshift=-3mm]llm.east);
  \draw[figure flow] ([yshift=3mm]exe.east) -- node[figure edge label,above=1pt] {调用} ([yshift=3mm]svc.west);
  \draw[figure flow] ([yshift=-3mm]svc.west) -- node[figure edge label,below=1pt] {返回} ([yshift=-3mm]exe.east);
  \draw[figure flow,dash pattern=on 2pt off 1.5pt] (llm.south) -- node[figure note,right=1pt] {写出与核对} (state.north);
\end{tikzpicture}

  \caption{基于大语言模型与函数调用的任务型对话。模型以系统提示、函数模式和对话历史为条件，或直接回复用户，或输出函数调用；执行器校验参数与权限、注入认证身份后调用服务，把返回结果作为工具消息追加到历史，模型再据此回复。图下方的状态记录是可选的显式状态，用于审计与核对。}
  \label{fig:nlp-dialogue-function-loop}
\end{figure*}

用三轮案例走一遍。第一轮，模型发现三个必需参数缺两个，不调用函数，直接询问日期和人数，相当于规则策略的询问动作。第二轮，模型输出调用\code{check\_registration(event\_id="ML-0418", date="2025-04-17", party\_size=2)}，执行器校验通过，返回空记录；模型据此回复“未查到匹配记录，您要核对4月18日吗”。第三轮，模型从历史中读出旧参数，只替换日期后再次调用，返回一条有效记录，模型生成告知。若第二次调用超时，工具消息写明\code{TIMEOUT}，系统提示中的规则要求模型说明失败而不推断记录是否存在。这一过程中，轮次理解、状态追踪和策略都隐含在一次生成里：模型从历史中读取槽值，自己判断参数是否齐备，再决定调用还是询问。

表\ref{tab:nlp-dialogue-pipeline-vs-llm}按模块对照两种实现。

\begin{table*}[htbp]
\centering
\small
\begin{tabular}{p{20mm}p{58mm}p{70mm}}
\toprule
环节 & 模块化流水线 & 大语言模型加函数调用 \\
\midrule
轮次理解 & 意图分类与槽位标注，需逐轮标注 & 隐含在生成中，由模型从话语和历史中读取 \\
状态追踪 & 显式状态$B_t$，规则或专用模型更新 & 状态隐含在上下文中；可另要求模型输出状态记录 \\
对话策略 & 规则、监督或强化学习，在固定动作集上选择 & 由系统提示中的规则和模型判断决定询问、调用或回复 \\
业务调用 & 槽位回填到固定接口 & 模型按函数模式生成参数，执行器校验与执行 \\
回复生成 & 模板或条件生成模型 & 模型依据工具消息直接生成 \\
新服务适配 & 补标注、扩动作集、增写模板，常需重训 & 改写函数模式与系统提示，少量示例 \\
可审计性 & 各模块输出可单独检查 & 只有调用记录天然可查，中间判断需另行记录 \\
典型错误 & 模块误差逐级传递；覆盖不到语料外情形 & 参数幻觉、误读旧值、未调用就声称结果、成本与延迟较高 \\
\bottomrule
\end{tabular}
\caption{模块化流水线与大语言模型加函数调用的对照。后者把四个模块合并到一次生成中，以函数模式取代固定的槽位与动作集合；业务正确性仍由执行器校验和真实返回保证。}
\label{tab:nlp-dialogue-pipeline-vs-llm}
\end{table*}

这种实现减少了模块标注和误差传递，新增服务只需描述新函数，也能应对流水线语料中没有的表达。它的问题集中在状态与执行两处。状态隐含在上下文中，模型可能从历史中读出已被更正的旧日期，或把用户说的人数与通知中的数字混淆；对专门评测的研究发现，大语言模型在显式状态追踪上仍落后于专用模型，但给定正确状态时能较好地把对话引向完成\scite{nlp-dialogue-hudecek2023}。因此实践中常让模型在每轮额外输出一份状态记录（相当于把SimpleTOD的状态段保留下来），或调用专门的状态更新函数，使状态可以检查和回放。执行一侧，模型可能编造参数、未调用就声称查到结果，或在用户与环境都能改变共享状态时失去一致；评测因此需要检查最终的环境状态，而不只是回复文本\scite{nlp-intro-barres2025}。写操作还应要求用户显式确认，并限制模型可调用的函数集合。

模型怎样获得按模式生成调用、读取结果并继续推理的能力，属于工具使用的训练问题，第\ref{chap:nlp-solving-and-execution}章讨论工具调用数据的构造、工具检索与参数生成，以及多步执行中的失败恢复。本节的结论是：函数调用把任务型对话的四个模块压缩为“模型决策、执行器校验、真实返回作依据”的循环，流水线中显式维护的状态和权限检查在新框架中并没有消失，而是转移到了执行器和状态记录中。

\section{开放对话}
\label{sec:nlp-dialogue-open}

任务型对话有一张待填写的业务表，成败由状态和服务结果判断。\term{开放对话}（open-domain dialogue）没有这样的目标：用户可能延续一个知识话题、聊自己的兴趣，或说出一次不顺利的经历，系统要做的是给出与上文相关、信息充分而又得体的回复。这一任务的挑战可以归为五项：回复要贴合上文并带有具体内容，避免“好的”“我不知道”一类放在哪里都成立的话；用到知识时不编造；保持指定人物设定和对真实用户记忆的一致；识别用户的情绪与求助目的并给出合适的支持；评价时把事实、连贯性和用途分开，而不是压成一个分数。

本节先在第\ref{sec:nlp-dialogue-open-lineage}小节沿方法演进处理第一项挑战，并说明大规模模型和对齐训练怎样改变开放对话；其后三小节分别以知识、人物与情绪为条件，处理其余挑战。多方会话还需保留说话者与回应对象：“我第一次参加”属于当前发言人，不一定属于群聊中上一位同名成员。下面各方法都假定历史范围和身份绑定已经给定。

\subsection{开放对话的挑战与模型谱系}
\label{sec:nlp-dialogue-open-lineage}

开放对话的模型演进可以按“前一步留下的主要缺陷”串起来。检索式回复从候选库中挑选人写的句子，流畅但受限于库内已有的说法；生成式回复逐词产生新句子，却偏好高频的通用回复；互信息重排把“回复是否特定于上文”加入打分，抑制通用回复；大规模预训练对话模型用更多数据和参数提高相关性与知识量，并开始把知识、人物和共情技能混合训练；指令微调与偏好对齐再让模型遵循用户要求、减少有害回复，形成今天的通用聊天模型。表\ref{tab:nlp-dialogue-open-lineage}按同一口径列出这些方法，以及后三小节的条件化方法。

\begin{table*}[htbp]
  \centering
  \small
  \setlength{\tabcolsep}{3pt}
  \renewcommand{\tabularxcolumn}[1]{>{\raggedright\arraybackslash}p{#1}}
  \begin{tabularx}{\linewidth}{@{}>{\raggedright\arraybackslash}p{30mm}XXX@{}}
    \toprule
    方法 & 解决的问题 & 核心机制 & 新代价或遗留问题\\
    \midrule
    检索式回复 & 生成句子不通顺 & 编码上文与候选回复，按匹配分选人写回复 & 无法说出库外内容；候选库需维护\\
    \addlinespace
    序列到序列生成 & 候选库覆盖不足 & 编码器—解码器以最大似然拟合上文到回复 & 偏好通用回复，前后不一致\\
    \addlinespace
    互信息目标与重排 & 通用回复 & 惩罚无条件高概率回复，或用反向模型$p(S\mid T)$重排 & 需要反向模型与候选列表；反语言模型解码易不合语法\\
    \addlinespace
    DialoGPT、Meena & 小语料下相关性与知识不足 & 在大规模社交对话上训练大模型，Meena以SSA做人工评价 & 事实错误、重复与人物矛盾仍在\\
    \addlinespace
    BlenderBot & 单一语料只练一种技能 & 预训练后在知识、人物、共情混合数据上微调，可接检索 & 解码需长度与重复控制；知识仍可能编造\\
    \addlinespace
    指令微调与对齐的聊天模型 & 不遵循指令、内容不安全 & 多轮示范的有监督微调，再用偏好与安全反馈对齐 & 标注与评判成本；过度拒答与奖励偏离\\
    \addlinespace
    Transformer Memory Network & 回复缺少知识依据 & 检索知识句并注意，排序或生成回复 & 只用一句知识；选择错误向后传播\\
    \addlinespace
    Profile Memory Network与用户记忆 & 人物与偏好前后不一致 & 人物条目作为独立记忆注意；用户记忆按身份与有效期筛选 & 资料来源、许可与过期需另行规定\\
    \addlinespace
    MoEL、ESConv & 不辨情绪与求助目的 & 情绪专家混合，或先预测支持策略再生成 & 情绪与策略标注成本；不等于专业帮助\\
    \bottomrule
  \end{tabularx}
  \caption{开放对话的方法谱系。前六行沿“检索—生成—互信息—大规模预训练—对齐”的演进，每行解决上一行的主要缺陷；后三行是第\ref{sec:nlp-dialogue-knowledge}至\ref{sec:nlp-dialogue-empathy}小节以知识、人物和情绪为条件的方法。}
  \label{tab:nlp-dialogue-open-lineage}
\end{table*}

\subsubsection{检索式与生成式回复}
\label{sec:nlp-dialogue-retrieval-generation}

最直接的做法是不写新句子，而从人写的回复中挑一句。\term{检索式回复}（retrieval-based response selection）预先收集候选回复库$\mathcal{R}$，对上文$c$和每个候选$r\in\mathcal{R}$打匹配分，输出分数最高者。Lowe等人在Ubuntu技术支持聊天记录上给出的双编码器基线是这一路线的标准形式\scite{nlp-dialogue-lowe2015ubuntu}：两个LSTM分别读入上文与候选，取末步隐状态$\bm{c},\bm{r}\in\mathbb{R}^{d}$，候选是正确回复的概率为
\begin{equation}
 p(\text{正确}\mid c,r)=\sigma\bigl(\bm{c}^{\top}\bm{M}\bm{r}+b\bigr),
 \label{eq:nlp-dialogue-dual-encoder}
\end{equation}
其中$\bm{M}\in\mathbb{R}^{d\times d}$和$b$是可学习参数。$\bm{M}\bm{r}$可以理解为由候选“预测”出的上文表示，再与真实上文做内积。训练时把真实下一轮作为正例，从语料中随机抽取的其他回复作为负例，用二元交叉熵训练；评价常用$R_{10}@1$，即在1个正确回复与9个随机回复中把正确者排第一的比例。

式\eqref{eq:nlp-dialogue-dual-encoder}属于第\ref{sec:nlp-classification-matching-architectures}小节的表示型匹配：候选向量与上文无关，可以离线算好并建索引，所以库很大时仍然快。把上文与候选拼接后联合编码的交叉编码器更准确，但每个候选都要重算；Poly-encoder用若干个上文“代码”向量与候选向量做一次注意，在两者之间折中\scite{nlp-matching-humeau2020polyencoder}。无论用哪种打分器，检索式回复的句子都来自人写语料，因而流畅；它的根本限制是只能复述库中已有的说法，用户提到库外的新组合时只能找一句“差不多”的回复。

\term{生成式回复}（generative response）去掉候选库，直接把上文当作源序列、回复当作目标序列，用第\ref{sec:nlp-generation-encoder-decoder}小节的编码器—解码器建模
\begin{equation}
 p(T\mid S)=\prod_{i=1}^{|T|}p(t_i\mid S,t_{<i}),
 \label{eq:nlp-dialogue-seq2seq}
\end{equation}
其中$S$是上文词元序列，$T=(t_1,\ldots,t_{|T|})$是回复。训练对每个（上文，回复）对最小化$-\log p(T\mid S)$，推断时用束搜索取高概率回复。Shang等人在微博短文本对话上、Vinyals与Le在技术支持记录和电影字幕上最早展示了这条路线：模型能够生成语料中没有原样出现过的回复\cite{nlp-dialogue-shang2015,nlp-dialogue-vinyals2015}。这条路线随后暴露出两个缺陷：回复常常是与上文关系不大的通用句；Vinyals与Le也指出，对“你多大了”“你是哪里人”之类的问题，换一种问法就给出互相矛盾的回答，因为模型只拟合每一对上文与回复，没有一致的“说话者”。第一个缺陷由下一小节处理，第二个缺陷在第\ref{sec:nlp-dialogue-personalization}小节处理。

\subsubsection{通用回复问题与互信息重排}
\label{sec:nlp-dialogue-mmi}

为什么最大似然训练会偏好通用回复？语料中“哈哈，是的”“我不知道”这类句子能接在许多不同的上文之后，因而在几乎任何$S$下都有不低的$p(T\mid S)$；与上文紧密相关的具体回复只适合少数上文，每一种具体说法分到的概率都较小。束搜索取最高概率的回复，就系统地落到通用句上。问题不在于模型没学会，而在于目标$\log p(T\mid S)$本身不区分“在这个上文下合适”与“在所有上文下都合适”。

Li等人把选择准则改为上文与回复之间的点互信息\scite{nlp-dialogue-li2016mmi}。\term{最大互信息目标}（maximum mutual information objective, MMI）要求回复在给定上文时的概率明显高于它不看上文时的概率：$\log\frac{p(S,T)}{p(S)p(T)}=\log p(T\mid S)-\log p(T)$。引入权重$\lambda$控制对无条件概率的惩罚强度，得到
\begin{equation}
 \hat T=\argmax_{T}\bigl\{\log p(T\mid S)-\lambda\log p(T)\bigr\}.
 \label{eq:nlp-dialogue-mmi-antilm}
\end{equation}
$p(T)$由一个只看回复的语言模型估计，所以这一形式称为MMI-antiLM：通用回复的$p(T)$大，被扣分多。用贝叶斯公式$\log p(T\mid S)=\log p(S\mid T)+\log p(T)-\log p(S)$代入，$\log p(S)$与$T$无关，可以去掉：
\begin{equation}
 \hat T=\argmax_{T}\bigl\{(1-\lambda)\log p(T\mid S)+\lambda\log p(S\mid T)\bigr\}.
 \label{eq:nlp-dialogue-mmi-bidi}
\end{equation}
这一形式称为MMI-bidi。$p(S\mid T)$由一个反向的序列到序列模型估计，它问的是“看到这句回复，能否猜回原来的上文”；“哈哈，是的”几乎猜不出任何上文，具体回复则能。两式在$\lambda$相同时选出同一个回复，差别在于计算上怎样近似。

用一个教学数例看重排效果。上文$S$为“我这周第一次去读书会，讨论的是注意力机制”，候选$T_1$为“哈哈，是的”，$T_2$为“注意力那一章的推导你跟上了吗”。假设正向模型给出$p(T_1\mid S)=0.05$、$p(T_2\mid S)=0.02$，回复语言模型给出$p(T_1)=0.04$、$p(T_2)=0.0005$。按最大似然，$T_1$胜出。取$\lambda=0.5$代入式\eqref{eq:nlp-dialogue-mmi-antilm}：$T_1$的分数为$\ln0.05-0.5\ln0.04=-2.996+1.609=-1.386$，$T_2$为$\ln0.02-0.5\ln0.0005=-3.912+3.800=-0.112$，$T_2$胜出。$T_1$之所以落败，是因为它在给定上文时的概率只比无条件概率高25\%，而$T_2$高40倍。

两种形式各有实现困难。MMI-antiLM可以在束搜索中逐词计算，但惩罚整个$\log p(T)$会偏向语言模型认为不像话的序列，也就是不合语法的回复。原文观察到上文对解码的影响主要体现在开头几个词元，之后解码器更多依赖自己已生成的前缀，于是只惩罚前$\gamma$个位置：把$p(T)$换成$U(T)=\prod_{k=1}^{|T|}p(t_k\mid t_{<k})^{g(k)}$，其中$g(k)=1$当$k\le\gamma$，否则为0。MMI-bidi则无法直接解码，因为$p(S\mid T)$要在整句$T$生成完后才能计算。实际做法是先用正向模型生成候选列表，再重排，如算法\ref{alg:nlp-dialogue-mmi-rerank}。重排分数中额外加入长度项$\eta|T|$：对数概率逐词累加，短句天然占优，长度项给每个词元一份固定奖励来抵消这一偏向。$\lambda$与$\eta$在开发集上调节。

\begin{algorithm}[MMI-bidi候选重排]
  \label{alg:nlp-dialogue-mmi-rerank}
  \begin{algorithmic}[1]
    \Require 上文$S$；正向模型$p(T\mid S)$；反向模型$p(S\mid T)$；候选数$N$；权重$\lambda,\eta$
    \Ensure 回复$\hat T$
    \State 用正向模型做束宽为$N$的束搜索，得到候选$T^{(1)},\ldots,T^{(N)}$及$\log p(T^{(n)}\mid S)$
    \For{$n=1,\ldots,N$}
      \State 用反向模型对固定的$S$做一次教师强制前向，得到$\log p(S\mid T^{(n)})$
      \State $\mathrm{score}_n\gets\log p(T^{(n)}\mid S)+\lambda\log p(S\mid T^{(n)})+\eta\,|T^{(n)}|$
    \EndFor
    \State \Return $\hat T=T^{(n^*)}$，$n^*=\argmax_n\mathrm{score}_n$
  \end{algorithmic}
\end{algorithm}

算法中的重排分数把式\eqref{eq:nlp-dialogue-mmi-bidi}的$(1-\lambda)$系数吸收进了$\lambda$的取值，排序不受影响。它是第\ref{sec:nlp-inference-reranking}节候选重排的一个特例：打分器不是独立训练的质量模型，而是同一语料上反向训练的生成模型，与统计机器翻译中的噪声信道模型同构。评价多样性时常用distinct-$n$，即生成结果中不同的$n$元组数目除以生成词元总数；原文报告MMI解码在BLEU与distinct-1、distinct-2上都优于最大似然基线。

互信息重排的代价有三项：要另训一个与正向模型同规模的反向模型；重排要对$N$个候选各做一次反向前向；束搜索的候选彼此相近，若列表里根本没有具体回复，重排也选不出来，所以后来的工作常改用采样产生更分散的候选。它也只处理“与上文的特定性”，不检查回复中的事实是否正确，后者需要第\ref{sec:nlp-dialogue-knowledge}小节的知识条件。

\subsubsection{大规模预训练对话模型}
\label{sec:nlp-dialogue-pretrained}

互信息重排修正的是选择准则，模型本身仍在几百万对话对上训练，知识量与表达方式都有限。第\ref{chap:nlp-large-language-models}章的预训练经验提示了另一条路：在规模大得多的对话语料上训练大模型，让相关、具体的回复本身获得更高概率。

DialoGPT沿用GPT-2的因果Transformer解码器，在从Reddit讨论串中抽取的约1.47亿段对话交流上训练\scite{nlp-dialogue-zhang2020dialogpt}。它不区分编码器和解码器：一段多轮会话$x_1,\ldots,x_N$的各轮依次拼接，每轮末尾加结束标记，整体当作一条长文本。若前$m$个词元构成上文$S$、其余构成回复$T$，目标为
\begin{equation}
 p(T\mid S)=\prod_{n=m+1}^{N}p(x_n\mid x_1,\ldots,x_{n-1}).
 \label{eq:nlp-dialogue-dialogpt}
\end{equation}
对多轮会话，第2轮到第$K$轮回复的联合概率分解为各轮在全部前文条件下的概率之积，一条会话因而为每一轮回复都提供监督。与式\eqref{eq:nlp-dialogue-seq2seq}相比，上文与回复共享同一组参数和同一个注意力栈，上文中的每个位置也能通过因果掩码被后续所有位置读取。DialoGPT发布了1.17亿、3.45亿和7.62亿参数三种规模，解码时先用top-$k$采样生成一组候选，再用反向模型$p(S\mid T)$按算法\ref{alg:nlp-dialogue-mmi-rerank}的思路重排，可见互信息重排在大模型上仍被沿用。

Meena进一步把规模推到26亿参数，使用1个编码器块和13个解码器块的Evolved Transformer序列到序列结构，上文最多取前7轮，训练语料是约341GB的公开社交媒体对话，目标仍是最小化回复的困惑度\scite{nlp-dialogue-adiwardana2020meena}。它的另一项贡献是人工评价协议。\term{敏感性与具体性平均}（sensibleness and specificity average, SSA）让评审对每条回复回答两个是非问题：在上文中是否说得通（敏感性），是否特定于这个上文（具体性）；不敏感的回复直接记为不具体，SSA是两项比例的平均。“我喜欢网球”后接“那挺好”说得通但不具体，接“我也是，费德勒的比赛我每场都看”两项都满足。具体性一项正是针对通用回复设计的。Meena报告困惑度与SSA高度相关，完整版本的SSA为79\%，人类对话为86\%。解码采用“采样再排序”：基础版本以温度0.88独立采样20条回复，取对数概率最高的一条，SSA为72\%；79\%的完整版本改用$k=40$的top-$k$采样，并加入检测跨轮重复的规则。

BlenderBot指出规模之外还要有技能\scite{nlp-dialogue-roller2021blender}。它先在约15亿条Reddit评论上预训练Transformer序列到序列模型（9000万、27亿与94亿参数），再在混合技能数据上微调：Persona-Chat练人物一致性，Wizard of Wikipedia练使用知识，EmpatheticDialogues练共情\cite{nlp-intro-zhang2018persona,nlp-intro-dinan2019,nlp-dialogue-rashkin2019empathetic}，Blended Skill Talk则在同一段对话中要求三种技能交替出现\scite{nlp-dialogue-smith2020bst}。BlenderBot比较了三类结构：Poly-encoder检索式模型、生成式模型，以及\term{检索后精炼}（retrieve and refine）模型——先检索一条回复或一句知识，把它拼接到上文后面再交给生成器。原文的人工比较发现解码设置影响很大：默认束搜索倾向过短的通用回复，强制最短20个词元后回复明显更受评审偏好；屏蔽重复三元组也有改善，但差异不显著。最终设置为束宽10、最短20词元加三元组屏蔽。后续工作研究了跨会话的长期记忆与互联网检索，分别对应第\ref{sec:nlp-dialogue-user-memory}与第\ref{sec:nlp-dialogue-knowledge}小节的问题\cite{nlp-dialogue-xu2022goldfish,nlp-dialogue-komeili2022internet}。

三者的关系可以这样概括：DialoGPT和Meena说明规模能显著提高相关性与具体性，BlenderBot说明混合技能数据和解码约束同样重要。它们共同的缺陷也在原文中被明确列出：模型仍会编造事实、自相矛盾、在长对话中重复，并模仿论坛语料中的冒犯性内容。前两项需要外部知识与记忆，后一项需要下一小节的对齐训练。

\subsubsection{聊天模型的指令微调与对齐}
\label{sec:nlp-dialogue-chat-alignment}

预训练对话模型学到的是“论坛里的人接下来会说什么”，这与“助手应当怎样回应用户”不同：前者可以跑题、挖苦或给出未经核实的说法，也不会按用户指令改变篇幅或格式。今天的通用聊天模型由预训练语言模型出发，先在人写的指令示范上做有监督微调，再用成对偏好训练奖励模型并做强化学习，或用直接偏好方法对齐\scite{ouyang2022instructgpt}。示范数据、损失构造与对齐算法分别在第\ref{chap:nlp-supervised-finetuning}章与第\ref{chap:nlp-preference-alignment}章展开，此处只写开放对话特有的三处变化。

第一处是多轮数据的监督位置。一条训练样本是按第\ref{sec:nlp-sft-context}小节对话模板拼接的完整会话，设第$k$轮助手回复的词元位置集合为$A_k$，损失只覆盖助手位置：
\begin{equation}
 L=-\sum_{k=1}^{K}\sum_{t\in A_k}\log p_{\theta}(x_t\mid x_{<t}),
 \label{eq:nlp-dialogue-multiturn-sft}
\end{equation}
这正是第\ref{sec:nlp-sft-masked-loss}小节的掩码损失在多轮上的写法。一条$K$轮会话因而同时监督$K$次回复，每次都以全部前文为条件，模型能学到延续前面约定的行为。它带来的困难是长会话中早期约束容易被遗忘。Llama~2-Chat为此提出Ghost Attention：构造训练数据时把系统指令（如“始终用一句话回答”）附加到每一轮用户消息后生成回复，训练时只在第一轮保留该指令，并把前面各轮的损失置零，使模型学会在后续轮次中仍遵守只出现过一次的指令\scite{nlp-align-touvron2023llama2}。

第二处是偏好反馈的单位。对话中的比较通常以“同一段历史下的两个候选回复”为单位收集，奖励模型以完整历史为条件给回复打分，比较关系按第\ref{sec:nlp-align-bradley-terry}小节的Bradley--Terry模型转为训练目标。Llama~2-Chat分别训练帮助性与安全性两个奖励模型：提示被标为安全相关，或安全奖励低于0.15时，取安全奖励；否则取帮助奖励\scite{nlp-align-touvron2023llama2}这样做的原因是两项要求常常冲突，一个分数难以同时刻画。

第三处是安全约束。开放对话没有业务接口来限定输出范围，用户也可能在多轮中逐步诱导模型。LaMDA的做法有代表性：在人工标注的安全与质量数据上微调同一个模型，使其能给候选回复输出安全分$s_{\mathrm{safe}}(y)$和质量分$q(y)$（原文取$3P(\text{合理})+P(\text{具体})+P(\text{有趣})$），推断时先采样若干候选，再按
\begin{equation}
 y^{*}=\argmax_{y\in\mathcal{Y},\ s_{\mathrm{safe}}(y)\ge\tau}q(y)
 \label{eq:nlp-dialogue-safety-filter}
\end{equation}
选择回复，其中$\mathcal{Y}$为候选集合、$\tau$为安全阈值；涉及事实时调用检索工具，把查到的内容放入上下文后再生成\scite{nlp-dialogue-thoppilan2022lamda}。更普遍的做法是把有害请求和多轮诱导样本纳入示范与偏好数据，或按第\ref{sec:nlp-align-cai}小节的原则批评生成AI反馈，再用对抗性的多轮测试检查遗漏。

与前面几代模型相比，对齐后的聊天模型能按指令改变回复方式，有害输出也更少；代价是示范、偏好与安全标注的成本，以及对齐目标本身带来的问题：评审偏爱长而自信的回复会让奖励被利用，安全约束过强会让模型拒答无害请求。事实依据、人物一致性与情绪支持也不会因对齐自动解决，仍需下面三小节的条件与检查。

\subsection{知识支持的交流}
\label{sec:nlp-dialogue-knowledge}

用户追问时，回复应当用上正确的知识。这里有两个子问题：从知识库中选出与当前话题相关的知识句；在回复中准确使用它，不添加没有来源的内容。Wizard of Wikipedia让一名“向导”在交流中检索维基百科并标出所用的知识句，为两个子问题分别提供了监督\scite{nlp-intro-dinan2019}。与知识配套提出的两种Transformer Memory Network分别对应第\ref{sec:nlp-dialogue-retrieval-generation}小节的检索式与生成式两条路线。以下用共同通知构造教学例子：用户先问活动时间，下一轮问“那地点呢”，第二轮应沿用同一活动而切换待答属性，不重复时间。

\subsubsection{检索式Transformer Memory Network}
\label{sec:nlp-dialogue-retrieval-memory}

检索式版本先用信息检索取得知识句，再据知识和历史给候选回复排序。Wizard of Wikipedia的各模型共用同一检索步骤：以话题及最近两轮为查询检索百科条目的首段，把它们拆成附带标题的知识句；召回器也可换成第\ref{sec:nlp-dialogue-dpr}小节的稠密检索。共享Transformer把历史编码为$\bm{q}$，把第$j$句知识编码为$\bm{m}_j$，对知识做注意并与历史相加：
\begin{equation}
 \alpha_j=\frac{\exp(\bm{q}^{\top}\bm{m}_j)}
                  {\sum_k\exp(\bm{q}^{\top}\bm{m}_k)},
 \qquad
 \bm{v}=\bm{q}+\sum_j\alpha_j\bm{m}_j.
 \label{eq:nlp-dialogue-retrieval-memory}
\end{equation}
另一Transformer把候选回复编码为$\bm{r}_l$，按$\cos(\bm{v},\bm{r}_l)$取最高者。训练以正确回复为正例、批内其他回复为负例做交叉熵；知识注意权重在回复目标下间接学习。

取知识$K_1$为“活动4月18日14:00举行”，$K_2$为“地点海棠楼201”，教学向量$\bm{m}_1=(1,0)$、$\bm{m}_2=(0,1)$。第一轮问时间，$\bm{q}=(1,0)$，知识权重为$(0.731,0.269)$，$\bm{v}=(1.731,0.269)$；时间回复与地点回复的向量取$(1,0)$与$(0,1)$，余弦分数约为0.988与0.154，选择时间回复。下一轮若历史被正确编码为$(0,1)$，权重与得分对调，选出地点回复。真实模型的错误可能出在召回错活动、选错属性或候选回复空泛，三者应分别计算知识召回、知识选择和回复选择；最高注意权重便于检查模型偏向，但不是逐句引用。

\subsubsection{生成式Transformer Memory Network}
\label{sec:nlp-dialogue-generative-memory}

检索式版本只能从候选回复中选，不能把知识组织成新的措辞。生成式版本逐词生成回复，有两种变体。两阶段版本单独训练知识选择器，取最高分知识句，与历史拼接后交给另一个编码器，再训练解码器预测参考回复；两个模型可分别改进，但第一阶段选错会向后传播。端到端版本共享历史与知识的编码器，计算知识注意分布后仍取最高分的一句，把它的词元表示与历史表示拼接供解码器交叉注意，并加一项知识选择的辅助损失：
\begin{equation}
 L=(1-\lambda)L_{\mathrm{response}}+\lambda L_{\mathrm{knowledge}}.
 \label{eq:nlp-dialogue-generative-memory}
\end{equation}
$L_{\mathrm{response}}$是参考回复的词元负对数似然，$L_{\mathrm{knowledge}}$是人类所选知识句的交叉熵。由于最终只取一句，这一版本不是对全部知识句的软加权生成；辅助损失直接监督知识注意，解码损失训练实际采用的表示。

上例第二轮中，设选择器给地点句的概率为0.8，人类标签也是地点句，则知识损失为$-\log0.8=0.223$。生成器据该句输出“地点是海棠楼201”。若输出为“海棠楼201，现场免费提供资料”，前半句有依据，后半句没有，知识选择正确并不使整句通过。两种变体都只用一句知识，需要组合多个片段时可改用第\ref{sec:nlp-dialogue-evidence}节的RAG或FiD接口。

\subsection{社交与个性化交流}
\label{sec:nlp-dialogue-personalization}

第\ref{sec:nlp-dialogue-retrieval-generation}小节提到，生成模型对同一问题的不同问法会给出互相矛盾的个人信息。人物条件为模型指定一个主体，使同一问题在不同主体下有不同而各自一致的回答。Persona-Chat给交流双方各指定几条人物资料，比较带资料与不带资料的回复\scite{nlp-intro-zhang2018persona}；它评价的是与指定人物的一致性，不涉及资料在现实中是否为真。用于真实用户时，问题变为长期记忆：哪些资料允许使用、从何而来、何时失效。

\subsubsection{Profile Memory Network人物条件回复}
\label{sec:nlp-dialogue-profile-memory}

Profile Memory Network把人物条目作为独立记忆，而不是与历史无区别地拼成一句长文本。排序式版本把历史表示$\bm{q}$与每条人物资料$\bm{p}_j$做相似度softmax得到权重，形成$\bm{q}^{+}=\bm{q}+\sum_j\alpha_j\bm{p}_j$，再按与候选回复的相似度排序，以正确回复和负回复的排序损失训练。生成式版本用LSTM编码历史并初始化解码器，人物条目的词向量按逆词频权重$1/(1+\log(1+\mathrm{tf}))$加权求和为向量后组成矩阵$\bm{P}$；解码第$i$步按$\operatorname{softmax}(\bm{P}\bm{W}_a\bm{h}_i)$对人物条目注意，所得摘要与词元输入结合预测下一词。与把全部资料直接前置的序列到序列基线相比，区别在于模型可以按当前话题选择资料条目。

教学角色甲的资料为“喜欢机器学习”“去年做过读书会志愿者”，角色乙为“喜欢机器学习”“这次第一次参加读书会”。对“以前来过吗”，兴趣条目权重应低，经历条目决定回答：甲两条权重若为$(0.2,0.8)$，回复“我去年做过志愿者”；乙则回复“这是我第一次参加”。注意必须先绑定当前扮演的角色，只在该角色的资料上计算，不能把甲乙资料合在一起。资料只说做过志愿者，就不支持“我去年主持了活动”；人物一致性、话题相关性和现实真实性应分开评分。

\subsubsection{用户记忆检索与规则更新}
\label{sec:nlp-dialogue-user-memory}

对真实用户，人物资料变成跨会话积累的记忆。Xu等构造了多会话对话数据，把此前各次会话摘要成记忆条目，下次对话时检索相关条目作为条件\scite{nlp-dialogue-xu2022goldfish}；这解决了“记得住”，但没有规定记忆的身份、来源和有效期。这里采用本书的教学协议：每条记忆保存主体标识、属性、值、来源轮次、确认时点、有效范围和使用许可；用户明确陈述的内容才进入候选写入，推测保留为未确认。语义编码器提供相关性分数，其参数可来自第\ref{chap:nlp-classification-matching-organization}章的匹配训练；身份、有效期和更新优先级由规则决定。

假设两位同名“林宁”的标识为U17、U28。记录M1为U17在4月10日确认“默认先给简短回复”，M2为U28偏好“完整推导”，M3为U17在4月17日说“这次请详细解释”。当前请求来自U17。即使M2的语义相似度0.95高于M1的0.80，也因主体不同被过滤；M3只对本次任务有效，优先于长期默认M1但不删除它，下一次请求仍用M1。若U17说“以后都改成详细解释”，确认后把M1标为失效并写入新的长期记录；只说“这个回答太短”不足以推出所有未来任务都要详细。同一时点相互矛盾的明确指令保留冲突并询问，不以向量相似度裁决。评价可在同名、过期、更正和临时覆盖样本上分别检查命中、误用和错误写入。

\subsection{共情与支持性交流}
\label{sec:nlp-dialogue-empathy}

“我没赶上报名”可能表达失落，也可能只是请求核对截止时间。系统需要同时辨认情绪线索与交流目的，而不是看见负面词就统一安慰。合适的回应承认已表达的处境，不添造经历或诊断；用户要信息时给有依据的信息，要倾听时不急着替对方作决定。方法路线是先识别情绪或支持策略，再以其为条件生成。这里研究的是普通支持性交谈，评审的主观感受不能当作专业疗效的证据。

\subsubsection{情绪条件生成与MoEL}
\label{sec:nlp-dialogue-moel}

基线先从历史预测情绪标签$e$，再以$(C,e)$生成回复，分别用分类交叉熵和词元负对数似然训练。硬选一个标签会丢失混合感受；测试时若直接提供真实标签，又会高估实际的情绪识别能力。

MoEL把硬选择换成软组合\scite{nlp-dialogue-lin2019}。它由情绪跟踪器、多个情绪回应专家和一个汇总解码器组成。Transformer编码器用词、位置与说话者角色嵌入读取历史，查询表示$\bm{q}$对各情绪键$\bm{k}_e$做softmax得到分布$p_e$。每个专家是一个独立的Transformer解码层，在同一历史与回复前缀上输出$\bm{V}_e$；另有一个共享专家输出$\bm{V}_0$。组合与训练目标为
\begin{equation}
 \bm{V}_{\mathrm{mix}}=\bm{V}_0+\sum_e p_e\bm{V}_e,
 \qquad
 L=\alpha[-\log p_{e^*}]+\beta[-\log p(y^*\mid C)],
 \label{eq:nlp-dialogue-moel}
\end{equation}
其中$e^*$是标注情绪，$y^*$是参考回复，汇总解码器再变换$\bm{V}_{\mathrm{mix}}$并输出词元分布。混合发生在隐表示上，不是对多条完整回复或词元概率做加权平均。

“我第一次参加读书会，很期待”与“我第一次参加读书会，怕听不懂”事件相同而情绪不同。压缩成期待、焦虑两个教学专家，两句的情绪权重分别取$(0.8,0.2)$与$(0.2,0.8)$。若某步$\bm{V}_0=(0.5,0.5)$、$\bm{V}_{\text{期待}}=(1,0)$、$\bm{V}_{\text{焦虑}}=(0,1)$，两次混合分别为$(1.3,0.7)$与$(0.7,1.3)$，汇总层据此生成不同回应：前者可以问感兴趣的话题，后者可以承认担心并询问最担心哪类内容。若用户接着说“先别安慰，我只想知道地点”，情绪分布可能仍偏焦虑，但当前目的要求回答地点；情绪识别准确率、回复相关性与人工感受到的理解程度应分别比较。

\subsubsection{ESConv支持策略条件生成}
\label{sec:nlp-dialogue-esconv}

情绪标签说明用户怎样感受，却没有说明系统接下来该做什么。ESConv把支持行为标成更接近回复职责的策略，如提问、复述、反映感受、肯定与安抚、提供建议或信息，标注者依据指南选择策略并组织回应\scite{nlp-dialogue-liu2021}。原文以预训练对话模型为骨干，把策略特殊词元放在参考回复之前，以展平的历史为条件训练
\begin{equation}
 p(s,y\mid C)=p(s\mid C)\prod_i p(y_i\mid C,s,y_{<i}),
 \label{eq:nlp-dialogue-esconv}
\end{equation}
其中$s$为策略、$y$为回复。联合版本先预测策略再生成回复；给定正确策略的版本衡量表达能力，随机策略版本检查策略选择的贡献。三者输入信息不同，报告时应分开。

同一事件可对应三种目的。“错过报名让我很失落，只想说说”宜先反映感受；“报名什么时候截止”宜提供通知中的精确信息；“能帮我想想怎样避免下次再错过吗”才适合给建议。假设后一句的策略分布在反映感受、信息、建议上为$(0.2,0.1,0.7)$，选建议后生成“可以在确认下一次截止时间后设一个提前提醒”。检查时既看策略标签，也看文字是否真的执行了该策略。策略中的“自我披露”不授权普通助手虚构亲身经历；原研究的参与者评价也不能推广为临床或专业建议的有效性。

\section{协作式对话}
\label{sec:nlp-dialogue-collaboration}

用户说“把通知写得简单些”，并没有完整规定交付物。“简单”可能指篇幅、阅读难度或语气；直接改短有时合适，有时删掉了用户最关心的条件。协作式对话与用户共同确定并完成一件产物，挑战有三项：要求含混时，决定是否询问、问什么；多个可行方案难以抽象描述时，让用户在看得见的候选中选择，并把选择正确地记入要求；反馈只说“不满意”而不说怎么改，修改一处又可能牵动别处。对应的方法谱系是：目标澄清（规则询问、信息增益、学习的澄清问题排序）→候选讨论（差异候选构造、成对偏好排序）→互动修订（用户反馈条件编辑、模型自反馈修订）。生成与编辑的算法在第\ref{chap:nlp-text-generation}章，本节研究交流怎样决定要生成什么，以及产物是否落实了当前要求。

\subsection{目标澄清}
\label{sec:nlp-dialogue-clarification}

澄清的输入是含混请求、已有材料和已确认条件，输出是值得优先询问的问题。对“简单些”，可以直接问希望多短，可以给两个版本让用户比较，也可以先按默认理解交稿；选哪一种取决于误解的代价、问题是否好回答和用户愿意投入多少。已确定的活动事实不再询问，尚未确定的目标也不能由默认候选冒充用户确认。下面先给可解释的规则与信息增益基线，再介绍从数据中学习问题价值的方法。

\subsubsection{神经EVPI澄清问题排序}
\label{sec:nlp-dialogue-evpi}

规则和信息增益都需要人写出可能的目标与回答，换一个领域就要重写。Rao与Daumé III改为从论坛数据中学习：一个澄清问题有价值，是因为它的回答补上了原帖缺的信息\scite{nlp-dialogue-rao2018}。他们借用决策论中的\term{完全信息价值}（expected value of perfect information, EVPI）：如果知道问题$q$的回答$a$，原帖会变得多有用；对可能的回答取期望，就得到问$q$的价值。这一方法对候选问题排序，不生成新问题。

设当前帖子为$p$，候选问题为$q_i$。从相似帖子中检索出的问答对$(q_j,a_j)$提供可能的回答。LSTM编码帖子和问题，前馈网络$F_{\mathrm{ans}}$预测一个回答表示；$\widehat{\bm{a}}_j$与$\widehat{\bm{q}}_j$分别是候选回答与问题的平均词向量。记距离$d_{i,j}=1-\cos(F_{\mathrm{ans}}(p,q_i),\widehat{\bm{a}}_j)$，排序分数为
\begin{equation}
\begin{aligned}
 w_{i,j}&=\exp(-d_{i,j})\,
            \cos(\widehat{\bm{q}}_i,\widehat{\bm{q}}_j),\\
 S(q_i;p)&=\sum_j w_{i,j}U(p,q_j,a_j).
\end{aligned}
\label{eq:nlp-dialogue-evpi}
\end{equation}
$w_{i,j}$衡量“$a_j$像是$q_i$的回答”的程度，扮演回答概率的角色；$U$是效用网络，把帖子、问题与回答的LSTM表示送入前馈网络后经sigmoid输出，表示补上该回答后帖子的有用程度。原式未把$w_{i,j}$归一化为概率，这里保持原形式。训练有两个目标：回答表示损失让$F_{\mathrm{ans}}$的输出接近原配对回答；效用模型以原始的帖子—问题—回答三元组为正例、检索到的其他组合为负例做二分类。负例带噪声，因为别的问题也可能有用；监督来自论坛的交互记录，而非逐问题测得的真实效用。

对“把通知写得简单些”，取$q_1$为“需要限制篇幅吗”，$q_2$为“是否面向首次参加者”，候选回答$a_1$为“限制在80字左右”，$a_2$为“面向首次参加者”。设两个问题的交叉余弦为0.5、自身为1，距离矩阵为$\left[\begin{smallmatrix}0.2231&0.9163\\0.5108&0.5108\end{smallmatrix}\right]$，权重近似为$(0.8,0.2)$和$(0.3,0.6)$。给定效用$U_1=0.7$、$U_2=0.9$，得$S(q_1)=0.74$、$S(q_2)=0.75$，先问受众。展示问题后要等待用户的真实回答，不能把候选回答$a_2$写成已确认的事实。这一方法的限制也来自它的形式：候选里没有关于语气的问题就选不出它，效用分数也没有扣除用户回答的负担。

\subsubsection{规则询问与信息增益基线}
\label{sec:nlp-dialogue-information-gain}

规则基线先查必需条件：用途未定则问用途，两个硬要求冲突则先澄清冲突，其余采用声明的默认值并交付初稿。另一个可手算的基线把未知目标表示为有限个假设$G$，计算问题$q$的回答$A_q$能减少多少不确定性，再扣去询问成本$c(q)$：
\begin{equation}
 J(q)=H(G)-\sum_a p(a\mid q)H(G\mid A_q=a)-c(q),
 \label{eq:nlp-dialogue-information-gain}
\end{equation}
熵以2为底，前两项之差是回答带来的\term{信息增益}（information gain），即知道回答后目标熵的期望下降量。回答概率由假设先验与回答模型求得；若回答可能误解或拒答，需在$p(a\mid q,G)$中计入。把成本换算成与比特可比的单位是教学约定，不是普适标准。与上一小节相比，这个基线要求显式写出先验、回答和成本，换来的是每一步都可以复核。

规定四个互斥的教学目标：只缩短篇幅$g_1$，只降低难度$g_2$，只放松语气$g_3$，同时缩短并降低难度$g_4$，先验为$(0.4,0.3,0.1,0.2)$。问题$q_L$问“是否主要需要更短”，把$g_1,g_4$归入“是”，回答概率为$(0.6,0.4)$；$q_D$问“是否需要降低阅读难度”，把$g_2,g_4$归入“是”，概率为$(0.5,0.5)$。假设回答确定且诚实。先验熵为1.846比特。$q_L$的“是”后验为$(2/3,1/3)$，“否”后验为$(3/4,1/4)$，预期剩余熵$0.6\times0.918+0.4\times0.811=0.875$，信息增益0.971；$q_D$两组后验为$(0.6,0.4)$与$(0.8,0.2)$，预期剩余熵0.846，信息增益1。若$c(q_L)=0.10$、$c(q_D)=0.25$，净分数为0.871与0.750，优先问篇幅。

直接按最高先验目标$g_1$交稿，猜对概率为0.4；问$q_L$后各分支取最大后验目标，命中概率升到$0.4+0.3=0.7$。多问一轮的收益由此可见，但具体字数和难度等级仍未确定。熵减只衡量当前假设空间内的不确定性，可能与产物是否可用不一致；直接交付、给选项与先询问三种做法最终仍要按交付质量和用户投入比较。

\subsection{候选讨论}
\label{sec:nlp-dialogue-candidates}

用户常常说不清想要什么语气，却能看出哪个版本合适。候选讨论要让差异可辨，又不能拿一份正确、一份漏了事实的版本制造伪选择。实现上分两步：构造覆盖硬条件、彼此有实质差异的候选；把用户的选择转成对候选的偏好分数。用户点选某个版本或说“挺好”，只按当时问题的范围解释，不自动视为所有细节都已验收。

\subsubsection{约束覆盖与差异候选选择}
\label{sec:nlp-dialogue-diverse-candidates}

固定任务为发布当前活动简讯，已确认事实包括活动名、当前时间地点、校内参加条件和校外严格截止条件。生成器先以这些条件产生候选，规则检查字段覆盖；只有通过硬条件的候选才进入篇幅、语气与组织方式的比较，避免用多样性弥补事实缺失。例如版本甲为“机器学习读书会将于2025年4月18日（周五）14:00在海棠楼201举行。校内师生可直接参加，校外来宾须在4月16日18:00前提交报名信息。”版本乙为“机器学习读书会：2025年4月18日14:00，海棠楼201。校内师生可直接参加；校外来宾须在4月16日18:00前提交报名信息。”两者事实相同，差异可以说明为完整叙述与紧凑简讯。

候选较多时，用第\ref{chap:nlp-text-generation}章的最大边际相关性（MMR）逐个选入：设相关性为$R(y)$、已选集合为$S$，
\begin{equation}
 s(y)=\lambda R(y)-(1-\lambda)\max_{z\in S}\operatorname{sim}(y,z).
 \label{eq:nlp-dialogue-mmr}
\end{equation}
已选甲、$\lambda=0.7$时，乙的相关性0.9、与甲相似度0.6，得分0.45；几乎同文的丙为0.95与0.95，得分0.38，于是选乙。展示数量要计入用户的阅读负担，比较不同展示方案时应固定生成预算与展示条数。

\subsubsection{成对偏好下的候选排序}
\label{sec:nlp-dialogue-bradley-terry}

用户在多轮中反复二选一，积累的是成对比较。把它们变成候选排序，用的是第\ref{sec:nlp-align-bradley-terry}小节的Bradley--Terry模型：给每个候选一个分数$s_i$，选$i$胜过$j$的概率为$\sigma(s_i-s_j)$，用比较次数的负对数似然估计分数；模型的推导、零点不可识别与全胜发散在那里已经讲过。用在协作对话中，要处理的是三件事。

第一是读出分差。同一用户在相同目标下比较甲乙四次，乙胜3次、甲胜1次，固定$s_{\text{甲}}=0$，最大似然估计为$s_{\text{乙}}=\log3\approx1.10$，即乙的胜负比是3比1。这一估计把四次选择当作独立同分布，而顺序、记忆和疲劳都会影响重复询问，所以独立性是假设。第二是平局与带修改的反馈。“两个都行”“选乙但把日期恢复”不是一次单纯的乙胜：平局按事先约定剔除或用支持平局的扩展模型，修改意见作为独立反馈保留，交给下一小节的编辑。第三是作用范围。分数描述的是某个问题、某个用户条件下的相对选择，不证明乙的每条事实都对；聚合其他用户偏好得到的分数，也不能覆盖当前用户已经明确的要求。

用户说“第二个更短，但保留地点”，共同状态就更新为采用乙的紧凑形式，并把“地点必须保留”记为已确认条件。乙本来含地点，无需为显示“已修改”再改一次；是否进一步缩短、是否省略年份仍未确认。

\subsection{互动修订}
\label{sec:nlp-dialogue-revision}

选定候选后，交流的结果要落到产物上。系统保存稿件版本号、已确认条件和未决问题，把每条反馈绑定到它所针对的版本，避免用户评论旧稿时误改新稿。挑战在于反馈常只指出不满意之处，且修改一处可能要同步另一处。两条路线的区别在反馈从哪里来：用户反馈条件编辑依赖用户给出的意见，Self-Refine让模型自己生成反馈。两者都要先确定作用范围再编辑，以实际差异和验收判断是否完成，而不是以回复了一句“已修改”为准。

\subsubsection{用户反馈条件编辑}
\label{sec:nlp-dialogue-feedback-edit}

用户反馈条件编辑把（原稿$y$，反馈$f$）作为条件，生成修订稿$y'$。设$x$为事实材料与已确认条件，训练数据为原稿—反馈—修订稿三元组，目标是修订稿的词元负对数似然：
\begin{equation}
 L_{\mathrm{edit}}=-\log p_{\theta}(y'\mid x,y,f)
  =-\sum_{i=1}^{|y'|}\log p_{\theta}(y'_i\mid x,y,f,y'_{<i}).
 \label{eq:nlp-dialogue-feedback-edit}
\end{equation}
模型结构可以是把$x$、$y$、$f$拼接为源序列的编码器—解码器，也可以是第\ref{sec:nlp-generation-editing}节那样预测保留、删除、替换操作的局部编辑器；提示式方法则在固定接口下由大模型直接改写。与第\ref{chap:nlp-text-generation}章的一般编辑相比，新增的是反馈$f$：它只说明要改什么，没有给出改后全文。

训练目标只要求输出接近参考修订稿，不检查改动是否落在反馈要求的范围内，因此需要单独的编辑检查。用Levenshtein对齐求出原稿到修订稿的改动位置集合$E(y,y')$，再由反馈与当前状态确定必须改动的位置集合$R(f)$，定义
\begin{equation}
 \mathrm{NMR}=\frac{\lvert E(y,y')\cap R(f)\rvert}{\lvert R(f)\rvert},
 \qquad
 \mathrm{EXR}=\frac{\lvert E(y,y')\setminus R(f)\rvert}{\lvert E(y,y')\rvert}.
 \label{eq:nlp-dialogue-edit-rates}
\end{equation}
NMR是必要修改完成率，EXR是改动中未经授权的比例；两者都是本书为检查编辑范围给出的教学指标。$R(f)$必须由反馈与已确认事实确定，例如“改日期”确定的是所有表示当前活动日期的位置，而不是所有日期字样。

设教学旧稿有两处“4月17日”，用户说“改成4月18日，地点不变”，依据是共同通知中已确认的改期。用压缩词元序列表示：旧稿为“读书会／4月17日／海棠楼201／请于／4月17日／到场”，目标为“读书会／4月18日／海棠楼201／请于／4月18日／到场”，$R(f)$为第2、5个位置。若修订稿只替换了第2个位置，NMR为$1/2$；若两处都替换但把地点改成另一栋楼，NMR为1而EXR为$1/3$。若原稿另有一句“原定4月17日”，它表示旧计划，不属于$R(f)$。有歧义的“改到周五”需先绑定日期属性和年份，不能改报名截止；反馈与权威事实冲突时，要先澄清是在编辑假设通知还是修订真实材料。

\subsubsection{Self-Refine辅助自反馈修订}
\label{sec:nlp-dialogue-self-refine}

用户反馈条件编辑需要用户逐条指出问题。Self-Refine让同一个语言模型依次承担初稿生成、反馈和修订，不训练专门的反馈器，也不更新参数；三种角色分别由三类提示中的说明与少量示例指定\scite{nlp-dialogue-madaan2023}。设请求与材料为$x$，初稿为$y_0$，第$t$轮反馈为$f_t$：
\begin{equation}
\begin{aligned}
 y_0&=M(p_{\mathrm{gen}},x),\\
 f_t&=M(p_{\mathrm{feedback}},x,y_t),\\
 y_{t+1}&=M(p_{\mathrm{refine}},x,y_0,f_0,\ldots,y_t,f_t).
\end{aligned}
\label{eq:nlp-dialogue-self-refine}
\end{equation}
修订提示保留此前全部输出与反馈，循环到固定轮数或反馈中出现停止指示为止。停止指示本身也是模型输出。

假设初稿把要求写成“校外来宾在4月16日前报名”。反馈提示要求“对照给定材料，定位被删去的时间与限定词，只指出可执行的修改”，得到“补回18:00及严格的前，其他字段保持不变”；修订后为“校外来宾须在4月16日18:00前提交报名信息”。再用规则检查截止时刻、严格不等号、主体与动作，才确认条件未变。这一轮共有生成、反馈、修订三次调用。

与用户反馈条件编辑相比，Self-Refine减少了对用户逐条指出问题的依赖，代价是多次模型调用，而且模型可能察觉不到自己的错误，或在反馈中重复同一错误。自反馈的优先级低于用户已确认的要求：若它建议“删除地点以便更短”，而地点已被确认必须保留，编辑器应拒绝这条建议。评价时分别记录用户轮数与自循环调用数，收益由最终产物的独立检查衡量。

表\ref{tab:nlp-dialogue-collaboration-methods}汇总本节方法。

\begin{table*}[htbp]
  \centering
  \small
  \setlength{\tabcolsep}{3pt}
  \renewcommand{\tabularxcolumn}[1]{>{\raggedright\arraybackslash}p{#1}}
  \begin{tabularx}{\linewidth}{@{}>{\raggedright\arraybackslash}p{28mm}XXX@{}}
    \toprule
    方法 & 解决的问题 & 核心机制 & 新代价\\
    \midrule
    规则询问 & 缺少必需条件或硬要求冲突 & 按条件清单询问，其余用声明的默认值 & 规则覆盖有限\\
    \addlinespace
    信息增益 & 不知先问哪个问题 & 对有限目标假设计算熵减并扣除成本 & 需写出先验、回答模型与成本单位\\
    \addlinespace
    神经EVPI排序 & 人写假设难以迁移 & 从论坛问答学习回答表示与效用，按期望效用排序 & 只能排序已有候选；不计用户负担\\
    \addlinespace
    差异候选选择 & 用户难以抽象描述偏好 & 硬条件过滤后用MMR选有差异的候选 & 用户阅读负担；展示数量需控制\\
    \addlinespace
    成对偏好排序 & 多次选择需汇总 & Bradley--Terry把成对选择转为分数 & 独立性假设；平局与修改需另记\\
    \addlinespace
    用户反馈条件编辑 & 意见需落到具体位置 & 以原稿与反馈为条件生成修订稿，再检查改动范围 & 需要三元组数据或固定接口；依赖用户逐条反馈\\
    \addlinespace
    Self-Refine & 用户逐条反馈负担大 & 同一模型循环生成、反馈、修订 & 调用次数增加；自评不能替代核验\\
    \bottomrule
  \end{tabularx}
  \caption{协作式对话的方法对照。前三行处理目标澄清，中间两行处理候选讨论，后两行处理互动修订；后一组方法都以前一组确定的要求为条件。}
  \label{tab:nlp-dialogue-collaboration-methods}
\end{table*}

\section{问答与对话的综合比较}
\label{sec:nlp-dialogue-comparison}

前五节分别处理了五类研究内容，它们可以共享编码、检索和生成部件，但各自的挑战与完成条件不同。表\ref{tab:nlp-dialogue-lineage}把全章的方法谱系放在同一张表中：每类研究内容按“挑战—方法—解决前者什么问题—代价”展开，同一类内的方法自上而下大致对应演进顺序。读表时可以沿两个方向看：横向看某个方法为解决什么而提出、付出了什么；纵向看同一挑战在问答侧和对话侧的不同处理，例如“证据不足时弃答”与“目标未定时澄清”都是在信息不够时不硬给答案，“引用生成”与“知识支持的回复”都要求输出可追溯到来源。

\begin{BookLongTable}{@{}>{\RaggedRight\arraybackslash}p{17mm}>{\RaggedRight\arraybackslash}p{27mm}>{\RaggedRight\arraybackslash}p{34mm}>{\RaggedRight\arraybackslash}p{\dimexpr.5\linewidth-39mm-8\tabcolsep\relax}>{\RaggedRight\arraybackslash}p{\dimexpr.5\linewidth-39mm\relax}@{}}
\toprule
研究内容 & 挑战 & 方法 & 解决前者的什么问题 & 代价\\
\midrule
\endfirsthead
\multicolumn{5}{@{}l}{\small 表~\thetable（续）}\\
\toprule
研究内容 & 挑战 & 方法 & 解决前者的什么问题 & 代价\\
\midrule
\endhead
\midrule
\multicolumn{5}{r@{}}{\small 续下页}\\
\endfoot
\bottomrule
\caption{问答与对话的方法谱系。每类研究内容内，箭头表示后一方法针对前一方法的缺陷提出；“代价”一列列出新方法引入的主要成本或遗留问题。各方法的细节见第\ref{sec:nlp-dialogue-evidence}至第\ref{sec:nlp-dialogue-collaboration}节。}
\label{tab:nlp-dialogue-lineage}\\
\endlastfoot
    证据问答 & 在给定段落中定位答案并弃答 & BiDAF$\to$\allowbreak BERT跨度阅读 & 预训练联合编码取代手工交互层 & 只在已给窗口内定位\\
     & 大规模文集中找证据 & BM25$\to$\allowbreak DPR$\to$\allowbreak REALM、Atlas & 稠密检索补词项不匹配；联合训练让检索器服务阅读器 & 索引构建与刷新；假负例\\
     & 综合多个片段生成答案 & RAG、FiD & 答案不限于跨度，可融合多片段 & 生成可能无依据\\
     & 查询不准、初检不精、一次检索不够、检索时机固定、答案难追溯 & 查询改写与HyDE、monoBERT重排、IRCoT、Self-RAG、ALCE式带引用生成 & 分别改进检索查询、排序、多跳检索、检索时机与归因 & 调用次数与延迟增加\\
     & 资料能否整体放入上下文 & 长上下文直接阅读、Self-Route & 免去检索遗漏；按需在两条路线间切换 & 位置效应；计算随长度增长\\
     & 证据在表格与数据库中 & SQL与图查询执行、TAPAS、TAGOP、FinQANet & 支持筛选、聚合与多步计算 & 需要模式、算子或程序监督\\
\midrule
    核验 & 主张粒度与证据定位 & FEVER流水线$\to$\allowbreak FActScore分解 & 逐事实判定，避免一处错误拖累整句 & 分解质量决定上限\\
     & 支持、反驳与信息不足 & 可分解注意力、规则与结构化比较、SAFE & 从固定证据扩展到检索证据 & 搜索成本与可复现性\\
     & 据核验结果修订 & RARR、CoVe & 最小编辑保留原回答；自问自答查自相矛盾 & CoVe不引入外部证据\\
\midrule
    任务型对话 & 理解每轮意图与槽位 & CRF槽位标注$\to$\allowbreak JointBERT$\to$\allowbreak 模式约束记录生成 & 联合建模意图与槽位；生成可表达更正 & 生成结果需类型检查\\
     & 跨轮累积与更正需求 & 规则写入$\to$\allowbreak TRADE$\to$\allowbreak SimpleTOD$\to$\allowbreak IC-DST & 复制处理未见值；序列化统一接口；少样本迁移 & 误复制旧值；示例检索敏感\\
     & 选择动作并调用服务 & 有限状态策略$\to$\allowbreak 监督动作预测$\to$\allowbreak 强化学习策略$\to$\allowbreak 大语言模型函数调用 & 从手写规则到从数据和交互中学策略 & 用户模拟器偏差；参数与权限需校验\\
     & 如实告知结果 & 模板回填$\to$\allowbreak SC-LSTM & 语义条件生成减少模板维护 & 软门不保证不遗漏\\
\midrule
    开放对话 & 通用乏味回复 & 检索式$\to$\allowbreak 生成式$\to$\allowbreak MMI重排 & 生成新句；惩罚无条件高概率回复 & 反向模型与候选列表\\
     & 相关性、知识量与技能 & DialoGPT、Meena、BlenderBot & 规模与混合技能数据提高相关性与具体性 & 编造事实、重复与矛盾仍在\\
     & 遵循指令与安全 & 指令微调与偏好对齐的聊天模型 & 按要求回应，减少有害输出 & 标注成本；过度拒答\\
     & 知识、人物与情绪一致 & Transformer Memory Network、Profile Memory Network、用户记忆、MoEL、ESConv & 以知识、人物、情绪或策略为条件 & 条件标注；条件选错向后传播\\
\midrule
    协作对话 & 要求含混 & 规则询问、信息增益、神经EVPI & 从手写假设到学习问题价值 & 只能排序已有候选；用户负担\\
     & 方案难以抽象描述 & 差异候选、成对偏好排序 & 让用户在可见候选中选择并汇总选择 & 阅读负担；独立性假设\\
     & 反馈落到具体位置 & 反馈条件编辑、Self-Refine & 把意见转成局部修改；减少用户逐条反馈 & 改动范围需检查；自评不替代核验\\
\end{BookLongTable}

表中的方法不能用一个统一的“回答质量”比较，因为它们的完成条件不同。比较前要固定问答的资料版本、可见片段与检索和生成预算，固定对话的服务快照、接口权限、用户或模拟器协议、最大轮数与中断条件；允许多种正确表达、等价证据组合或合法交流路径时，在评分前写出等价规则。下面用两个例子说明不同指标的分母怎样决定了结论。

问答侧的分数各自针对不同对象。规范化准确匹配检查答案是否与某个允许参考一致；词元F1按预测与参考的词元多重集合交集计算精确率与召回率。参考为“4月18日／14:00”而预测只有“4月18日”时，精确率1、召回率$1/2$、F1为$2/3$，准确匹配为0。若时间与报名条件都需回答而召回只取得片段A，按两项所需证据计的证据召回率为$1/2$；若生成了两个正确陈述却都引用A，引用支持率仍为$1/2$。答案、证据覆盖和引用三个分数不能互相替代。

对话侧把每轮有效的槽位—值对作为集合，报告槽位F1和联合目标准确率。表\ref{tab:nlp-dialogue-metrics}沿用第\ref{sec:nlp-dialogue-task}节的三轮主例，假设某系统前两轮正确、第三轮错误保留4/17。两边都为空的轮次单独记为正确，非空槽位采用跨轮汇总的微平均。

\begin{table}[htbp]
\centering
\begin{tabular}{cllc}
\toprule
轮次 & 正确状态 & 假设预测状态 & 全状态正确\\
\midrule
1 & 日期、人数均缺失 & 日期、人数均缺失 & 是\\
2 & 4/17，2人 & 4/17，2人 & 是\\
3 & 4/18，2人 & 4/17，2人 & 否\\
\bottomrule
\end{tabular}
\caption{局部状态分数与完整任务的差别。活动编号和身份为固定上下文，不计入可修改槽位；第三轮旧日期使查询继续为空。}
\label{tab:nlp-dialogue-metrics}
\end{table}

三轮共4个正确正值、4个预测正值，匹配3个，微平均精确率、召回率和F1均为$3/4$，联合目标准确率为$2/3$。但第三轮唯一需要变化的日期没有改对，变化项召回率为0；系统可以每轮都给出礼貌回复，却没有查到用户要核对的报名，完整任务成功为0。三个数字都计算无误，回答的却是不同问题：槽位F1问多少槽位值对，联合准确率问多少轮整体对，任务成功问用户的目标是否达成。完整交互还应报告用户补充或更正次数、总轮数、检索与模型调用成本，自反馈调用不记为用户轮次。

比较两个系统时，以相同问题或同一初始对话任务配对。假设四个教学任务上A的最终成功为$(1,1,0,0)$、B为$(1,0,1,0)$，平均成功率都是$1/2$，逐任务差为$(0,1,-1,0)$：两者修复的是不同任务，平均值相同不说明失败机制相同。样本足够时按完整任务配对重采样计算差值区间；多轮对话不能把各轮当成独立样本来虚增样本量。人工选择、Bradley--Terry聚合与第\ref{sec:nlp-eval-judge-protocol}小节的模型评判分别记录，并以独立人工或可执行任务依据校验。诊断时依次检查资料是否取得、本轮是否理解、状态是否更新、动作是否合规、回复是否有据、目标是否达成；给定正确证据或正确状态的诊断实验有助于隔离环节，但不替代端到端比较。


\section*{文献说明}
\label{sec:nlp-dialogue-literature}

本章的原始资料按支持的内容分为方法与数据两类。各节的数值例子、通知资料、只读报名接口、用户记忆协议、通知差异候选和信息增益基线都是本书的教学实现，不归为下列论文的原算法或实验结果。

阅读理解与检索方面，BiDAF给出字符与词嵌入、双向注意力流与建模层的原始结构\cite{nlp-dialogue-seo2017}，BERT给出联合编码的跨度阅读与空答案阈值\cite{devlin2019}；检索加阅读的流水线始于DrQA\cite{nlp-qa-chen2017drqa}，DPR给出双编码器检索与带片段选择的阅读器\cite{nlp-karpukhin2020dpr}，REALM与Atlas分别给出以答案似然和阅读器蒸馏训练检索器的方法\cite{nlp-qa-guu2020realm,nlp-qa-izacard2023atlas,nlp-qa-izacard2021distilling}。检索增强生成方面，RAG与FiD给出片段边缘化与解码器融合两种接入方式\cite{lewis2020rag,nlp-dialogue-izacard2021}；查询改写、HyDE、monoBERT重排、IRCoT与Self-RAG分别针对查询、排序、多跳与检索时机\cite{nlp-qa-ma2023rewrite,rs-s30,rs-s20,rs-s40,nlp-qa-asai2024selfrag}；ALCE给出带引用生成的引用召回与精确率\cite{nlp-qa-gao2023alce}；长上下文与检索的比较参考位置效应研究与两项对照实验\cite{liu2024middle,nlp-qa-xu2024retrieval-long,nlp-qa-li2024selfroute}。结构化与混合证据方面，TAPAS给出单元格选择、聚合头与软答案损失\cite{nlp-dialogue-herzig2020}，TAGOP与FinQANet分别给出标签加算子与程序生成的做法\cite{nlp-dialogue-zhu2021,nlp-dialogue-chen2021}。核验与对话方面，FActScore给出原子事实精确率\cite{min-etal-2023-factscore}，JointBERT、TRADE与SimpleTOD依次对应轮次理解、复制式状态追踪与序列化任务对话\cite{nlp-dialogue-chen2019,nlp-dialogue-wu2019,nlp-dialogue-hosseini2020}，SC-LSTM给出语义条件的回复生成\cite{nlp-dialogue-wen2015}，Self-Refine给出生成、反馈、修订的提示接口\cite{nlp-dialogue-madaan2023}。

数据与评价入口方面，SQuAD 2.0把不可回答性纳入阅读任务\cite{rajpurkar2018squad2}，Natural Questions以真实搜索查询构成开放域问答\cite{nlp-qa-kwiatkowski2019nq}，HotpotQA为多跳问题标注支持事实\cite{yang2018hotpotqa}，TAT-QA与FinQA分别收集表格—文本混合问题与带程序标注的财报问题\cite{nlp-dialogue-zhu2021,nlp-dialogue-chen2021}，FEVER给出主张核验的三类标签与证据句\cite{nlp-intro-thorne2018}，其共享任务报告定义了FEVER分数\cite{nlp-dialogue-thorne2018shared}。对话方面，MultiWOZ与SGD提供多领域任务型对话的状态标注\cite{budzianowski2018multiwoz,rastogi2020sgd}，$\tau^2$-Bench让用户与执行器共同改变环境\cite{nlp-intro-barres2025}，Wizard of Wikipedia、Persona-Chat与ESConv分别对应知识、人物与情绪支持条件下的交流\cite{nlp-intro-dinan2019,nlp-intro-zhang2018persona,nlp-dialogue-liu2021}。引用原文结果时，各节只保留说明机制所需的个别数字，不能据此对方法作统一排名。

\section{本章小结}
\label{sec:nlp-dialogue-summary}

回到章首的挑战清单。证据问答把“找到证据”和“形成答案”分开处理：检索器与重排器负责找到证据，稠密检索弥补词汇不匹配，联合训练让检索器学到阅读器真正需要的片段；阅读器、生成器与执行器负责形成答案，抽取式阅读定位跨度，RAG与FiD综合多个片段，SQL、TAPAS与数值程序处理筛选、聚合与计算。弃答阈值与带引用生成划定回答的边界：证据不足时不答，答了就要能逐句追溯到来源。查询改写、迭代检索与Self-RAG分别回答了“用什么查”“查几次”“何时查”，长上下文则提供了另一条以计算换召回的路线。核验从另一端处理同一问题，把已有回答拆成主张，区分支持、反驳与信息不足，只对需要的部分作最小修订。

对话侧的核心是维护随轮次变化的状态。状态追踪记录当前有效的需求，区分缺失、未知、不限定与否定，并在用户更正时撤掉旧值；策略与函数调用把状态连接到真实服务，回复只陈述服务实际返回的结果。开放交流没有任务终态，互信息重排抑制通用回复，知识、人物与情绪条件分别让回复有依据、有一致性、有合适的支持方式；协作交流在目标含混时以澄清、候选讨论与版本化修订逐步确定要求。

这些方法也留下了同一类边界：合法的状态可能已经过期，流畅的回复可能缺少依据，格式正确的函数调用可能尚未执行或执行失败。它们都不等于任务完成。下一章把已明确的语言要求转为数学解答、程序与环境中的行动，以执行结果和外部验证来判断任务是否完成。
